% !TEX TS-program = xelatex
% US English Translation of A. V. Knyazev (1986)
% Build: xelatex k_translation_en && bibtex k_translation_en && xelatex k_translation_en && xelatex k_translation_en
\documentclass[11pt,oneside]{book}
% Common preamble for Knyazev (1986) US English monograph translation
\usepackage[a4paper, left=25mm, right=25mm, top=27mm, bottom=25mm, headheight=26pt]{geometry}
\usepackage{amsmath, amssymb, mathtools}
\usepackage{fontspec}
\defaultfontfeatures{Ligatures=TeX}
% \setmainfont{CMU Serif}
% \setsansfont{CMU Sans Serif}
% \setmonofont{CMU Typewriter Text}
\setmainfont{cmunrm.otf}[
  BoldFont        = cmunbx.otf,
  ItalicFont      = cmunti.otf,
  BoldItalicFont  = cmunbi.otf,
  SlantedFont     = cmunsl.otf,
  BoldSlantedFont = cmunbl.otf
]
\setsansfont{cmunss.otf}[
  BoldFont       = cmunsx.otf,
  ItalicFont     = cmunsi.otf,
  BoldItalicFont = cmunso.otf
]
\setmonofont{cmuntt.otf}[
  BoldFont       = cmuntb.otf,
  ItalicFont     = cmunit.otf,
  BoldItalicFont = cmuntx.otf
]
\usepackage[american]{babel}
\usepackage{tikz}
\usetikzlibrary{arrows.meta, positioning, calc}
\usepackage{fancyhdr}
\usepackage{needspace}
\usepackage[square,sort&compress,colon]{natbib}
\setcitestyle{aysep={,},yysep={,},citesep={;}}
\usepackage{xurl}
\usepackage[unicode, colorlinks=true, linkcolor=black, urlcolor=blue, citecolor=black, bookmarksnumbered=true, pdfstartview=FitH]{hyperref}
\usepackage{bookmark}
\urlstyle{rm}

\makeatletter
% Catcode-safe and BibTeX-line-wrap-safe clickable DOI and URL macros
\newcommand*{\bibdoi}{%
  \begingroup
  \lccode`\~=`\#\relax \lowercase{\def~{\#}}%
  \lccode`\~=`\_\relax \lowercase{\def~{\_}}%
  \lccode`\~=`\&\relax \lowercase{\def~{\&}}%
  \lccode`\~=`\<\relax \lowercase{\def~{\bib@lt}}%
  \lccode`\~=`\>\relax \lowercase{\def~{\bib@gt}}%
  \lccode`\~=0\relax
  \catcode`\~=12\relax
  \catcode`\#=\active \catcode`\_=\active \catcode`\&=\active
  \catcode`\<=\active \catcode`\>=\active
  \@bibdoi
}
\def\@bibdoi#1{%
  \edef\#{\string##}\edef\_{\string_}\edef\&{\string&}\edef\%{\@percentchar}%
  \edef\bib@lt{\string<}\edef\bib@gt{\string>}%
  \edef\x{\toks0={{#1}}}\x
  \edef\#{\@percentchar23}\edef\_{_}\edef\&{&}\edef\%{\@percentchar}%
  \edef\bib@lt{\@percentchar3C}\edef\bib@gt{\@percentchar3E}%
  \edef\x{\toks2={\noexpand\href{https://doi.org/#1}}}\x
  \edef\x{\endgroup DOI:~\the\toks2{\noexpand\nolinkurl{\the\toks0}}}\x
}
\newcommand*{\biburl}{%
  \begingroup
  \lccode`\~=`\#\relax \lowercase{\def~{\#}}%
  \lccode`\~=`\_\relax \lowercase{\def~{\_}}%
  \lccode`\~=`\&\relax \lowercase{\def~{\&}}%
  \lccode`\~=`\<\relax \lowercase{\def~{\bib@lt}}%
  \lccode`\~=`\>\relax \lowercase{\def~{\bib@gt}}%
  \lccode`\~=0\relax
  \catcode`\~=12\relax
  \catcode`\#=\active \catcode`\_=\active \catcode`\&=\active
  \catcode`\<=\active \catcode`\>=\active
  \@biburl
}
\def\@biburl#1{%
  \edef\#{\string##}\edef\_{\string_}\edef\&{\string&}\edef\%{\@percentchar}%
  \edef\bib@lt{\string<}\edef\bib@gt{\string>}%
  \edef\x{\toks0={{#1}}}\x
  \edef\#{\string##}\edef\_{_}\edef\&{&}\edef\%{\@percentchar}%
  \edef\bib@lt{\@percentchar3C}\edef\bib@gt{\@percentchar3E}%
  \edef\x{\toks2={\noexpand\href{#1}}}\x
  \edef\x{\endgroup URL:~\the\toks2{\noexpand\nolinkurl{\the\toks0}}}\x
}
\makeatother

% Typography settings
\linespread{1.04}
\setlength{\parindent}{1.5em}
\setlength{\parskip}{2pt plus 1pt}
\tolerance=1200
\emergencystretch=2.5em
\allowdisplaybreaks[1]

% Page headers and footers
\pagestyle{fancy}
\fancyhf{}
\fancyhead[C]{\parbox[b]{\textwidth}{\centering\footnotesize\itshape \nouppercase{\rightmark}}}
\fancyfoot[C]{\small\thepage}
\renewcommand{\headrulewidth}{0.4pt}
\fancypagestyle{plain}{%
  \fancyhf{}%
  \fancyfoot[C]{\small\thepage}%
  \renewcommand{\headrulewidth}{0pt}%
}

% Custom structural macros matching US English monograph & TOC conventions
\newcommand{\monounnumchapter}[1]{%
  \clearpage
  \thispagestyle{plain}%
  \phantomsection
  \addcontentsline{toc}{chapter}{#1}%
  \markboth{#1}{#1}%
  \vspace*{6mm}%
  {\centering\LARGE\bfseries #1\par}%
  \vspace{6mm}%
}

\newcommand{\monochapter}[2]{%
  \clearpage
  \thispagestyle{plain}%
  \phantomsection
  \addcontentsline{toc}{chapter}{Chapter #1. #2}%
  \markboth{Chapter #1. #2}{Chapter #1. #2}%
  \vspace*{6mm}%
  {\centering\Large\bfseries Chapter #1.\par\smallskip\LARGE\bfseries #2\par}%
  \vspace{6mm}%
}

\newcommand{\monosuppchapter}[3]{%
  \clearpage
  \thispagestyle{plain}%
  \phantomsection
  \addcontentsline{toc}{chapter}{#1. #2. #3}%
  \markboth{#1}{#1}%
  \vspace*{6mm}%
  {\centering\Large\bfseries #1\par\medskip
   \large\mdseries #2\par\medskip
   \LARGE\bfseries #3\par}%
  \vspace{6mm}%
}

\newcommand{\monosection}[2]{%
  \par\vspace{5mm}%
  \needspace{5\baselineskip}%
  \phantomsection
  \addcontentsline{toc}{section}{\texorpdfstring{\S\,#1. #2}{§ #1. #2}}%
  \markright{\S\,#1. #2}%
  {\centering\large\bfseries \S\,#1. #2\par}%
  \vspace{3mm}%
}

% natbib bibliography layout matching monograph style
\renewcommand{\bibsection}{\monounnumchapter{References}}
\setlength{\bibhang}{1.8em}
\setlength{\bibsep}{4pt plus 1pt}

% Mathematical operators & symbols (US English AMS conventions)
\newcommand{\myspan}{\operatorname{span}}
\newcommand{\diag}{\operatorname{diag}}
\newcommand{\ang}{\measuredangle}
\hypersetup{
  pdftitle={A. V. Knyazev. Computation of Eigenvalues and Eigenvectors in Mesh Problems: Algorithms and Error Estimates (English Translation of the 1986 Russian Monograph)},
  pdfauthor={A. V. Knyazev}
}
\begin{document}

% =============================================================================
% SOURCE MODULE: part_00_front_intro.tex
% =============================================================================
% === Orig. cover and title page [s01L, s01R, s02L, s02R] ===
\begin{titlepage}
\thispagestyle{empty}
\centering
\vspace*{10mm}
% {\large Academy of Sciences of the USSR\par\smallskip
% Department of Numerical Mathematics\par}

\vspace{45mm}
{\Large A.\,V.~Knyazev\par}

\vspace{12mm}
{\LARGE\bfseries COMPUTATION OF EIGENVALUES\par\smallskip
AND EIGENVECTORS IN MESH PROBLEMS:\par\smallskip
ALGORITHMS AND ERROR ESTIMATES\par}

\vspace{6mm}
{\large (English Translation of the 1986 Russian Monograph)\par}

\vfill
% {\large Moscow\quad 1986\par}
{\large 2026\par}
\vspace{10mm}
\end{titlepage}

% === Orig. p. 2 [s03L] ===
\thispagestyle{empty}
\vspace*{8mm}
{\centering\bfseries A.\,V.~Knyazev. Computation of Eigenvalues and Eigenvectors\\
in Mesh Problems: Algorithms and Error Estimates\\
{\mdseries (English Translation of the 1986 Russian Monograph)}\par}
\vspace{8mm}

This monograph is an authorized English translation of the author's 1986 Russian book published by the Department of Numerical Mathematics of the USSR Academy of Sciences (Moscow, 1986, LCCN~87157255). Methods for solving symmetric algebraic eigenvalue problems arising from the mesh discretization of continuous problems are studied. For various vector iteration methods, subspace iteration methods, and their counterparts with auxiliary operators, convergence rate estimates are proved that help clarify the efficiency of these methods. Asymptotic properties are investigated and error estimates are proved for the Temple--Lehmann method for computing two-sided eigenvalue bounds and for the Rayleigh--Ritz method. An extensive bibliography of works devoted to various aspects of solving symmetric eigenvalue problems is provided.

The monograph is intended for specialists in numerical mathematics.

\vspace{12mm}
\noindent Reviewer of the Russian original in 1986: E.\,G.~D'yakonov

\vfill
% \noindent{\small Knyazev, A.\,V. Computation of Eigenvalues and Eigenvectors in Mesh Problems: Algorithms and Error Estimates, 1986, pp.\,1--188. Price 90 kop.\par\medskip}
% \noindent\copyright\ Department of Numerical Mathematics, Academy of Sciences of the USSR, 1986.
\noindent{\small English translation of the original Russian monograph: A.\,V.~Knyazev, ``\textit{Vychislenie sobstvennykh znachenii i vektorov v setochnykh zadachakh: algoritmy i otsenki pogreshnosti}'' (``Вычисление собственных значений и векторов в сеточных задачах: алгоритмы и оценки погрешности''), Dept.\ of Numerical Mathematics, USSR Academy of Sciences (OVM AN SSSR), Moscow, 1986, pp.\,1--188, LCCN~87157255. Digitized and translated with AI assistance and verified by the author. Posted with permission of the Marchuk Institute of Numerical Mathematics RAS (INM RAS).\par\smallskip
\noindent
\clearpage

% === Orig. pp. 3--4 [s03R, s04L] ===
{
\renewcommand{\contentsname}{Contents}
\tableofcontents
}

\monounnumchapter{Preface to the English Translation}

Compared with the original typewritten and hand-inscribed print, all hand-inscribed formulas and diagrams have been typeset in \LaTeX\ using US English mathematical conventions (for example, $\tan$ in place of $\operatorname{tg}$). The alphabetical Bib\TeX\ list with English-translated Russian titles marked ``[in Russian]'' uses modernized \texttt{[Author, Year]} citations with verified entries with clickable DOIs or URLs where available. Actually cited in text are 73 references, with additionally included uncited 600 relevant items of general bibliography, matching the 1986 monograph, cleaned to the best of our ability.

To preserve historical fidelity with the 1986 Russian edition, the mathematical content, notation, equation numbering, and exposition---including original mathematical misprints and typos in the 1986 print---have been retained verbatim without substantive revisions (known misprints are documented as comments in the \LaTeX\ source; only obvious scanning and punctuation glitches were regularized). Future releases are anticipated to fix typos and errors; see all \LaTeX\ sources at:
\url{https://github.com/lobpcg/math/tree/main/2026_Computation_of_Eigenvalues_and_Eigenvectors_in_Mesh_Problems_1986}.\par}

\noindent\textit{A.\,V.~Knyazev, 2026}

% === Orig. p. 5 [s04R] ===
% \monounnumchapter{Preface}
\vspace{8mm}
\phantomsection
\addcontentsline{toc}{chapter}{Preface to the 1986 Russian Edition}
\markboth{Preface to the 1986 Russian Edition}{Preface to the 1986 Russian Edition}
{\centering\LARGE\bfseries Preface to the 1986 Russian Edition\par}
\vspace{5mm}

The partial eigenvalue problem (computing the first few eigenvalues of a matrix or matrix pencil) is one of the traditional problems of numerical linear algebra. With the development of mesh (finite-difference, projection-difference, etc.) methods for discretizing spectral problems with differential operators, the need arises to solve such problems for large sparse matrices. This gives rise to the fundamental theoretical question of investigating the efficiency of various numerical methods and identifying among them those that are optimal for solving specific classes of such problems. On the whole, this question is quite difficult.

The present work addresses only one aspect of this question, namely, estimating the accuracy of various methods. In particular, convergence rate estimates for many iterative methods are proved in unified terms. Special attention is paid to deriving estimates that are sharp (unimprovable) in a certain sense.

In their subject matter, the first three chapters overlap with the book \citep{Parlett1983}. The concluding fourth chapter is devoted to efficient iterative methods with an auxiliary operator.

The exposition of the main results of this work is intended for a reader familiar with the theory of methods for solving symmetric spectral problems (e.g., from the books \citep{FaddeevFaddeeva1963,Parlett1983}) and the fundamentals of functional analysis. The proofs of the results are accessible mainly to specialists.

The author is grateful to E.\,G.~D'yakonov and V.\,I.~Lebedev for posing a number of problems and discussing the results, to A.\,L.~Skorokhodov for assistance in compiling the bibliography and preparing the text of this work, and to G.\,A.~Trunina and N.\,S.~Stepanenko for the prompt and high-quality formatting of the manuscript. % [orig. typo in Russian print: E.G. L'yakonovu]

% === Orig. p. 6 [s05L] ===
\monounnumchapter{Introduction}

The main body of the monograph consists of four chapters. The structure of logical dependencies among these parts is shown in Fig.\,1. The exposition is organized so that each chapter can be read independently of the others.

\begin{center}
\medskip
\begin{tikzpicture}[
    box/.style={draw, thick, align=left, inner sep=5pt, font=\small},
    arr/.style={-{Stealth[length=2.2mm, width=1.6mm]}, thick}
]
  \node[box, text width=52mm, anchor=north west] (ch1) at (0, 0)
    {Ch.\,1. Vector iterations};

  \node[box, text width=114mm, anchor=north west] (ch2) at (8mm, -13mm)
    {Ch.\,2. The Rayleigh--Ritz method in subspace iterations and\\ block algorithms};

  \node[box, text width=104mm, anchor=north west] (ch3) at (18mm, -30mm)
    {Ch.\,3. The Temple--Lehmann method for two-sided eigenvalue\\ bounds};

  \node[box, text width=122mm, anchor=north west] (ch4) at (0, -47mm)
    {Ch.\,4. Iterative methods with an auxiliary operator};

  % Arrows matching original diagram (s05L)
  % Ch1 -> Ch2
  \draw[arr] ($(ch1.south west)+(18mm,0)$) -- ($(ch2.north west)+(10mm,0)$);
  % Ch1 -> Ch4 (far left vertical line)
  \draw[arr] ($(ch1.south west)+(3mm,0)$) -- ($(ch4.north west)+(3mm,0)$);
  % Ch2 -> Ch3
  \draw[arr] ($(ch2.south west)+(18mm,0)$) -- ($(ch3.north west)+(8mm,0)$);
  % Ch2 -> Ch4 (second vertical line from left)
  \draw[arr] ($(ch2.south west)+(3mm,0)$) -- ($(ch4.north west)+(11mm,0)$);
\end{tikzpicture}
\par\smallskip
{\small Fig.\,1}
\medskip
\end{center}

In the present work we consider the following formulation of the partial symmetric generalized eigenvalue problem. In a Euclidean space (i.e., a complete finite-dimensional real linear space with an inner product) $H$, it is required to compute (approximately) the $p$ largest eigenvalues
\[
\lambda_1 > \lambda_2 > \ldots > \lambda_p
\]
and the corresponding eigenvectors
\[
u_1, u_2, \ldots, u_p
\]
of the problem
\[
Mu = \lambda Lu, \qquad M = M^*, \qquad L = L^* > 0 \tag{0.1}
\]
with spectrum
\[
\lambda_1 > \ldots > \lambda_p > \lambda_{p+1} \geqslant \ldots \geqslant \lambda_{\min}.
\]
The $p$ largest eigenvalues are assumed to be simple mainly in order to simplify the statements of the results. If it is required to determine the smallest eigenvalues in (0.1), then replacing $M$ by $-M$
% === Orig. p. 7 [s05R] ===
reduces this problem to the preceding one.

As the Euclidean space $H$, one can choose the space $\mathbb{R}^N$ of column vectors with the natural inner product, in which case the operators $M$ and $L$ are symmetric matrices. This choice is in fact implicitly assumed when describing the algorithms. In theoretical investigations, however, a basis of coordinate vectors is used very rarely, and therefore it is advisable to consider problem (0.1) in an arbitrary Euclidean space without fixing a basis in it.

Many methods for solving the generalized spectral problem (0.1) are based on its implicit reduction to a standard symmetric spectral problem
\[
Au = \lambda u, \qquad A = A^*, \tag{0.2}
\]
for example, with the operator $A = L^{-1}M$ in the space $H_L$ with inner product $(\cdot, \ast)_L = (L\cdot, \ast)$. This choice of the operator $A$ is treated in the present work as the primary one. Another way of reducing the generalized eigenvalue problem to a standard one (see the Contents) is associated with choosing $A = -M^{-1}L$ in the space $H_M$ when $M > 0$.

Chapters 1--3 are devoted to solving problem (0.2) using these two ways of reduction.

Chapter 1 studies vector iterations for computing an eigenvector of problem (0.2) corresponding to the largest eigenvalue. By ``vector iterations'' we mean iterative methods in which the current ($n$th) approximation is a vector ($u^n$). This term is introduced to distinguish such methods from subspace iterations (see below).

Chapter 2 investigates the capabilities of the Rayleigh--Ritz method both as a standalone method and as applied to subspace iterations and block iterative methods for computing several eigenvectors of problem (0.2) corresponding to the largest eigenvalues. The terms ``subspace iterations'' and ``block iterative methods'' (or ``simultaneous iterations'') are often used synonymously to denote iterative methods in which the current ($n$th) approximation is a subspace ($U^n$, $\dim U^n = p$) and the transition operator from $U^n$ to $U^{n+1}$ is independent of the choice of basis in $U^n$. In the present work, the term ``subspace iterations'' is used only when the transition operator from $U^n$ to $U^{n+1}$ is also independent of $U^n$.

% === Orig. p. 8 [s06L] ===
Chapter 3 is devoted to the Temple--Lehmann method for computing two-sided bounds for the largest eigenvalues of problem (0.2).

The concluding Chapter 4 studies iterative methods for solving problem (0.1) that do not rely on a reduction to (0.2). Their distinguishing feature is that during the iterative process one has to solve systems of equations of the form $Bu = f$ with some auxiliary operator $B = B^* > 0$, which in general coincides neither with $L$ nor with $M$. If the operator $B$ ``effectively models'' the operator $L$, then these methods converge as rapidly as the methods with the operator $A = L^{-1}M$ from Chapters 1--2.

Throughout the analysis, roundoff errors are assumed to be absent; therefore, normalization of the approximations is not indicated in the schemes of the iterative methods: while necessary in practice, it is immaterial for the theory. If in two vector iteration methods, given identical initial approximations ($u^0 = v^0$), the one-dimensional subspaces $\myspan\{u^n\}$ and $\myspan\{v^n\}$ coincide for all subsequent approximations $u^n$ and $v^n$, then these methods are considered equivalent.

Based on the results of this work, one can investigate the efficiency of various methods for solving mesh problems of the form (0.1) arising from the discretization of differential spectral problems, provided that the computational cost of performing the basic operations of the algorithm (for example, solving equations of the form $Lu = f$ in methods with the operator $A = L^{-1}M$) is known.

The choice of terminology (norms, constants) is geared toward the case where the mesh operator $L^{-1}M$ corresponds to a compact operator in the continuous problem. For example, $M$ is a mesh analogue of the operator of multiplication by a function, and $L = -\Delta_h$, where $\Delta_h$ is the mesh Laplace operator. Thus, it is implicitly assumed that for a small mesh size $h$ the spectrum of problem (0.1) exhibits ``clustering'' only at zero and that $\lambda_{p+1} > 0$. It follows, in particular, that the quantity $(\lambda_p - \lambda_{\min}) / (\lambda_p - \lambda_{p+1})$ may be regarded as bounded above as $h \to 0$, whereas $\lambda_p / \lambda_{\min} \to \infty$ as $h \to 0$ if $M > 0$.

In almost all results of this work, one can trace the nature of the dependence on $h$. The only exceptions are asymptotic statements (valid for a ``good'' approximation), such as Theorem~2.2.4, which the author nevertheless deemed appropriate to include in the text.

The reference list includes works that may be of interest to specialists engaged in the theory and/or practice of computing eigenvalues of ``large'' problems. In compiling
% === Orig. p. 9 [s06R] ===
the list, use was made of the bibliographies from the books \citep{FaddeevFaddeeva1963,Collatz1968ru,Parlett1983,Weinberger1974,Fichera1978,Chatelin1983,GolubVanLoan1983,CullumWilloughby1985}, the surveys \citep{Ikramov1982,Stewart1976b,KuttlerSigillito1984}, and the bibliographic indexes:

1. Faddeev, D.\,K. and Faddeeva, V.\,N., Computational methods of linear algebra, in: \textit{Computational Methods of Linear Algebra. Parallel Computations}, Zap.\ Nauchn.\ Sem.\ LOMI, Leningrad: Nauka, 1975, pp.\,3--228 [in Russian].

2. \textit{Computational Methods of Linear Algebra. Bibliographic Index 1975--1980} (ed.\ V.\,V.~Voevodin), Leningrad: LOMI USSR Acad.\ Sci., 1982, 334\,pp.\ [in Russian].

3. \textit{Computational Methods of Linear Algebra. Bibliographic Index 1981--1984} (ed.\ V.\,V.~Voevodin), Leningrad: Nauka, 1986, 480\,pp.\ [in Russian].

The bibliography offered here makes no claim to completeness; its large size is due rather to the diversity of topics covered by the included works.

% =============================================================================
% SOURCE MODULE: part_01_ch1.tex
% =============================================================================
% === Orig. p. 10 [s07L] ===
\monochapter{1}{Vector Iterations}

In this chapter, for problem (0.1) with $p = 1$, error measures are introduced in \S\,1.1 and convergence rate estimates are stated for the abstract vector iteration method. Specific algorithms are then examined: simple iteration and iteration with Chebyshev parameters (\S\,1.2), steepest ascent (\S\,1.3), the Lanczos method (\S\,1.4), and the conjugate gradient method (\S\,1.5). A more detailed description of these well-known algorithms can be found, for example, in \citep{FaddeevFaddeeva1963,Wilkinson1970ru,Parlett1983,LebedevBakhvalov1984,GolubVanLoan1983,Chatelin1983,CullumWilloughby1985}.

The presentation in \S\S\,1.1--1.5 deals with problem (0.2) with $A = L^{-1}M$ in the space $H_L$. In \S\,1.6 we discuss the effect on the convergence rate of the algorithms of another way of reducing problem (0.1) to (0.2) for $M > 0$, namely with $A = -M^{-1}L$ in the space $H_M$. In the concluding \S\,1.7, the method of inverse iteration with shifts for problem (0.1) is considered separately. The Supplement to Chapter~1, written jointly with A.\,L.~Skorokhodov, contains proofs of several results for the method of steepest ascent from \S\,1.3.

Many results of this chapter are fairly elementary and are given without proof.

On the other hand, certain theoretical aspects of the application of the classical algorithms listed above have not been sufficiently studied. This applies in particular to the Lanczos method and the conjugate gradient method. Therefore, alongside rigorously justified facts, we also state conjectures and hypotheses, mostly based on analogies with the theory of iterative methods for solving systems of linear algebraic equations or minimizing quadratic functionals.

% === Orig. p. 11 [s07R] ===
\monosection{1.1}{Abstract vector iteration method}

Let $u^0$ be a nonzero initial approximation to the vector $u_1$. An operator $F_k$, which is a function $F_k(\cdot)$ of the operator $A = L^{-1}M$, will be regarded as the transition operator taking $u^0$ to the $k$th approximation in abstract vector iterations:
\[
u^k = F_k u^0. \tag{1.1.1}
\]
Define the quantity
\[
\sigma^k \equiv \max_{i > 1} |F_k(\lambda_i)|^2 / |F_k(\lambda_1)|^2.
\]

Assume that $\sigma^k < 1$ and that we know (as a rule, theoretically) a number
\[
\bar{\sigma}^k \in [\sigma^k, 1). \tag{1.1.2}
\]
As measures of proximity of a vector $u \neq 0$ to $u_1$, we shall use the functionals
\[
\begin{aligned}
\tan^2\varphi(u) &\equiv \tan^2\ang_L(u; u_1), \\
r^2(u) &\equiv \{\lambda_1 - \lambda(u)\} / \cos^2\ang_L(u; u_1), \\
\Delta\lambda(u) &\equiv \{\lambda_1 - \lambda(u)\} / \{\lambda(u) - \lambda_2\},
\end{aligned}
\]
where $\lambda(u) \equiv (Mu, u)/(Lu, u) = (Au, u)_L / (u, u)_L$ is the Rayleigh quotient and the angle is measured in $H_L$, that is,
\[
\cos^2\ang_L(u; u_1) \equiv (u; u_1)_L^2 \|u\|_L^{-2} \|u_1\|_L^{-2}.
\]
Note that one of these functionals is the square of the ``energy'' seminorm of the error:
\[
r^2(u) = \|u - u_1\|_{\lambda_1 L - M}^2 = \|u\|_{\lambda_1 L - M}^2,
\]
provided $u$ and $u_1$ are normalized so that $\|u_1\|_L = (u; u_1)_L = 1$.

The well-known relation
\[
\lambda_1 - \lambda_2 \leqslant \{\lambda_1 - \lambda(u)\} / \sin\ang_L(u; u_1) \leqslant \lambda_1 - \lambda_{\min} \tag{1.1.3} % [math note: the original omits the exponent 2 on sine: \sin\ang_L(u; u_1) instead of \sin^2\ang_L(u; u_1)]
\]
% === Orig. p. 12 [s08L] ===
makes it possible to establish a connection among the functionals and, in particular, to conclude that as $u \to u_1$ all three functionals have the same order of smallness:
\[
\tan^2\varphi(u) \asymp \Delta\lambda(u) \asymp r^2(u).
\]

It is also easy to see that
\[
r^2(u) - \{\lambda_1 - \lambda_2\}\Delta\lambda(u) \asymp \tan^4\varphi(u). \tag{1.1.4}
\]
We set $\tan^2\varphi_n \equiv \tan^2\varphi(u^n)$, $r_n^2 \equiv r^2(u^n)$, $\Delta\lambda_n \equiv \Delta\lambda(u^n)$, $n = 0, k$.

An important role in establishing convergence rate estimates for vector iterations will be played by the following theorem.

\textbf{Theorem 1.1.1}~\citep{LebedevBakhvalov1984}\textbf{.} If at least one of the two equivalent conditions
\[
\tan^2\varphi_0 < \infty, \qquad r_0^2 < \infty, \tag{1.1.5}
\]
holds, then $u^k \neq 0$ and
\begin{gather*}
\tan^2\varphi_k \leqslant \bar{\sigma}^k \tan^2\varphi_0, \tag{1.1.6} \\
r_k^2 \leqslant \bar{\sigma}^k r_0^2. \tag{1.1.7}
\end{gather*}

Under the stronger condition
\[
\Delta\lambda_0 \in [0, \infty), \tag{1.1.8}
\]
the estimate
\[
\Delta\lambda_k \leqslant \bar{\sigma}^k \Delta\lambda_0 \tag{1.1.9}
\]
also holds.

From (1.1.3) and (1.1.6) follows the estimate
\[
\{\lambda_1 - \lambda(u^k)\} / \{\lambda(u^k) - \lambda_{\min}\} \leqslant \bar{\sigma}^k \tan^2\varphi_0. \tag{1.1.10}
\]

\textbf{Remark on the sharpness of the estimates.} Each of the estimates (1.1.6), (1.1.7), (1.1.9), and (1.1.10) for $\bar{\sigma}^k = \sigma^k$ is attainable in the sense that for
% === Orig. p. 13 [s08R] ===
any prescribed values of the parameters on the right-hand side of the estimate, there exist $u^0$ and $F_k$ for which the nonstrict inequality becomes an equality. The form of $u^0$ and $F_k$ can be given explicitly: depending on the estimate, we choose either $u^0 = u_1 + \alpha u_2$ or $u^0 = u_1 + \alpha u_{\min}$, while for the function $F_k$ it suffices that either
\[
\bar{\sigma}^k = |F_k(\lambda_2)|^2 / |F_k(\lambda_1)|^2 \tag{1.1.11}
\]
or
\[
\bar{\sigma}^k = |F_k(\lambda_{\min})|^2 / |F_k(\lambda_1)|^2 \tag{1.1.12}
\]
hold, likewise depending on the estimate.

\textit{Proof of Theorem 1.1.1.} Verifying inequalities (1.1.6) and (1.1.7) presents no difficulty and is carried out directly. As for formulas (1.1.9) and (1.1.10), their generalizations to subspace iterations will be proved directly in Chapter~2.

Note that studying the abstract method is very useful, since it actually yields convergence rate estimates for entire classes of algorithms. The advantages of such an approach have long been recognized in the theory of iterative methods for solving systems of linear algebraic equations; see, for example, the textbook \citep{SamarskiiNikolaev1978}. In the theory of iterative methods for solving eigenvalue problems, this approach is not widely used.

It should be pointed out that an estimate analogous to (1.1.10) is proved in the book \citep[p.~260]{Parlett1983}, where it is also applied to analyze the convergence rate of the Lanczos and simple iteration methods.

We draw the reader's attention to inequality (1.1.7) involving the functional $r^2(u)$. This essentially quadratic functional is investigated in detail and widely used in the present monograph, apparently for the first time. Its importance stems from property (1.1.4), which suggests a deep analogy between the asymptotic properties of methods for solving linear systems and those for eigenvalue problems.

% === Orig. p. 14 [s09L] ===
\monosection{1.2}{Simple iteration with a constant shift and iterations with Chebyshev parameters}

Consider an iterative scheme of the form
\[
u^{n+1} = Au^n - \gamma^n u^n, \quad n = 0, 1, \ldots, k - 1 \tag{1.2.1}
\]
(the normalization of the vectors $u^{n+1}$ required in practice is immaterial for theoretical analysis, and is therefore omitted here and below in algorithm schemes). In the case of a constant shift
\[
\gamma^n = \gamma < (\lambda_1 + \lambda_{\min}) / 2, \tag{1.2.2}
\]
Theorem~1.1.1 obviously holds for the method $\{(1.2.1),\,(1.2.2)\}$ with
\[
\bar{\sigma}^k = [\max\{|\lambda_{\min} - \gamma|; |\lambda_2 - \gamma|\} / |\lambda_1 - \gamma|]^{2k}. \tag{1.2.3}
\]
In particular, one can choose $\gamma = \lambda_{\min}$, for which
\[
\bar{\sigma}^k = [1 - \xi]^{2k}; \tag{1.2.4}
\]
here and throughout Chapter~1 (exceptions will be stated explicitly),
\[
\xi \equiv (\lambda_1 - \lambda_2) / (\lambda_1 - \lambda_{\min}). \tag{1.2.5}
\]
The optimal choice (in the sense of minimizing $\bar{\sigma}^k$ in (1.2.3)), $\gamma = (\lambda_2 + \lambda_{\min}) / 2$, yields
\[
\bar{\sigma}_k = [(1 - \xi) / (1 + \xi)]^{2k}. \tag{1.2.6} % [math note: the original has a subscript \bar{\sigma}_k instead of superscript \bar{\sigma}^k]
\]

In the case of variable shifts $\gamma^n$ in (1.2.1), the optimal choice is given by the Chebyshev parameters
\[
\gamma_n = \frac{1}{2}\{\lambda_2 + \lambda_{\min} - \tau^n (\lambda_2 - \lambda_{\min})\}, \quad n = 0, \ldots, k - 1, \tag{1.2.7} % [math note: the original has a subscript \gamma_n instead of superscript \gamma^n]
\]
where $\tau^n \in [-1, 1]$ are all $k$ roots of the Chebyshev polynomial of the first kind $T_k(\tau)$ of degree $k$. Then in the estimates of Theorem~1.1.1 and in corollary (1.1.10),
% === Orig. p. 15 [s09R] ===
\[
\bar{\sigma}^k = |T_k\{(1 + \xi) / (1 - \xi)\}|^{-2}. \tag{1.2.8}
\]

It is also easy to write down an expression for $\bar{\sigma}^k$ when, instead of $\lambda_2$ and $\lambda_{\min}$, one knows approximations (possibly crude) $\lambda_2^* \in (\lambda_{\min}, \lambda_1)$ and $\lambda_{\min}^* \leqslant \lambda_{\min}$ \citep{LebedevBakhvalov1984}. Various aspects of the practical implementation of the method $\{(1.2.1),\,(1.2.7)\}$ are examined in \citep{Lebedev1977,MarchukLebedev1981}: the possibility of refining estimates for $\lambda_2$ during the iterations, the dependence of the numerical stability of the algorithm on the ordering of the parameters $\tau^n$, optimal continuation of the iterations for $n > k$ (infinitely continuable processes), and others.

A three-term variant of the method $\{(1.2.1),\,(1.2.7)\}$, optimal for every $n = 0, 1, \ldots$, is well known:
\[
u^{n+1} = u^n - \alpha^n Au^n - \beta^n (u^n - u^{n-1}), \quad n = 0, 1, \ldots
\]
with parameters $\alpha^0 = 2 / (\lambda_2 + \lambda_{\min})$, $\beta^0 = 0$, $\alpha^n = 4\delta^n / (\lambda_2 - \lambda_{\min})$, $\beta^n = -\delta^n \delta^{n-1}$, $n \geqslant 1$, where $\delta^n = (2 / \delta^0 - \delta^{n-1})^{-1}$, $\delta^0 = (\lambda_2 - \lambda_{\min}) / (\lambda_2 + \lambda_{\min})$. Theorem~1.1.1 and corollary (1.1.10) with $\bar{\sigma}^k$ from (1.2.8) obviously hold for this scheme as well.

For all the methods mentioned above, except $\{(1.2.1),\,(1.2.2)\}$ when $\gamma > (\lambda_2 + \lambda_{\min}) / 2$, all the stated estimates are sharp in the sense that for any $c \in (0, \infty)$ we have, for example, for (1.1.6):
\[
\max_{u^0 : \tan^2\varphi_0 = c} \tan^2\varphi_k / \tan^2\varphi_0 = \bar{\sigma}^k, \tag{1.2.9}
\]
similarly for (1.1.7) and (1.1.9) (in terms of $r^2$ and $\Delta\lambda$), and for (1.1.10):
\[
\max_{u_0 : \tan^2\varphi_0 = c} \{\lambda_1 - \lambda(u^k)\} / \{\lambda(u^k) - \lambda_{\min}\} \cdot \tan^{-2}\varphi_0 = \bar{\sigma}^k. % [math note: the original has a subscript u_0 instead of superscript u^0]
\]

This is due to the fact that these methods generate only functions $F_k$ for which the required condition from (1.1.11), (1.1.12) is satisfied.

% === Orig. p. 16 [s10L] ===
\monosection{1.3}{Method of steepest ascent}

In this classical algorithm, the parameter $\gamma^n$ in (1.2.1) is chosen to maximize $\lambda(u^{n+1})$, which leads to the scheme
\[
\begin{aligned}
& u^{n+1} = Au^n - \alpha^n u^n, \quad n = 0, 1, \ldots, k - 1, \\
& \alpha^n \equiv \frac{1}{2}\left\{\lambda^n + \lambda(w^n) - \sqrt{[\lambda^n - \lambda(w^n)]^2 + 4\|w^n\|_L^2 \|u^n\|_L^{-2}}\right\}, \\
& \lambda^n \equiv \lambda(u^n), \quad w^n \equiv Au^n - \lambda^n u^n.
\end{aligned} \tag{1.3.1}
\]
Better known is another, equivalent form of this algorithm \citep{KantorovichAkilov1984}:
\[
\begin{aligned}
& u^{n+1} = u^n + \gamma^n w^n, \quad w^n \equiv Au^n - \lambda^n u^n, \quad n = 0, 1, \ldots, k - 1, \\
& \gamma^n = 2\left\{\lambda^n - \lambda(w^n) + \sqrt{[\lambda^n - \lambda(w^n)]^2 + 4\|w^n\|_L^2 \|u^n\|_L^{-2}}\right\}^{-1}, \\
& \lambda^n = \lambda(u^n).
\end{aligned} \tag{1.3.2}
\]

The complicated dependence of $\gamma^n$ on $u^0$ prevents a direct application of Theorem~1.1.1, since it is difficult to determine the quantity $\bar{\sigma}^k$. However, it is obvious that under condition (1.1.8), inequality (1.1.9) holds for steepest ascent (1.3.1) and (1.3.2) with $\bar{\sigma}^k$ at least no larger than that for simple iteration with the optimal shift, that is, no larger than (1.2.6). For convenience, we write down this estimate once again,
\[
\Delta\lambda_k \leqslant \bar{\sigma}^k \Delta\lambda_0, \quad \bar{\sigma}^k = [(1 - \xi) / (1 + \xi)]^{2k}, \tag{1.3.3}
\]
and note that it improves the result of \citep{Prikazchikov1975}. Do inequalities of the type (1.1.6), (1.1.7) hold for algorithm (1.3.1) or (1.3.2)? Answering this question required a special investigation, carried out jointly by the author and A.\,L.~Skorokhodov. The results are stated below, while the proofs are deferred to the Supplement to Chapter~1.

\textbf{Theorem 1.3.1.} Under conditions (1.1.5) and (3.1.1) or (3.1.2), % [math note: typo in formula numbers in the original: (3.1.1) or (3.1.2) instead of (1.3.1) or (1.3.2)]
the following inequalities hold:
\[
\tan^2\varphi_k \leqslant \bar{\sigma}^k \tan^2\varphi_0, \quad r_k^2 \leqslant \bar{\sigma}^k r_0^2, \quad \bar{\sigma}^k = (1 - \xi)^{2k}. \tag{1.3.4}
\]

% === Orig. p. 17 [s10R] ===
The estimates (1.3.4) are sharp in the sense that
\[
\sup_{u^k} \tan^2\varphi_{k+1} / \tan^2\varphi_k = \sup_{u^k} r_{k+1}^2 / r_k^2 = (1 - \xi)^2. \tag{1.3.5}
\]

Asymptotically (as $u^0 \to u_1$), the estimate $r_k^2 \leqslant \bar{\sigma}^k r_0^2$ from (1.3.4) can be improved by taking
\[
\bar{\sigma}^k = [(1 - \xi) / (1 + \xi)]^{2k} + o(1). \tag{1.3.6}
\]

Compared with (1.3.4), estimate (1.3.3) is sharp in a stronger sense:
\[
\sup_{u^0 : u^0 \approx u_1} \Delta\lambda_k / \Delta\lambda_0 = \bar{\sigma}^k, \tag{1.3.7}
\]
where $u^0 \approx u_1$ means that the vector $u^0$ lies in an arbitrarily small neighborhood of the vector $u_1$.

For estimates (1.3.3) and (1.3.4), equalities of the form (1.2.9) do not hold in general; however, attempts to refine the estimates accordingly lead to cumbersome expressions and are hardly justified.

There is a well-known approach to studying the method of steepest ascent in eigenvalue problems, based on an analogy with the theory of the method of steepest descent for minimizing arbitrary strongly convex functionals \citep{KantorovichAkilov1984}. However, unless $u^0 \approx u_1$, this analogy is far from straightforward, as the following example shows.

Method (1.3.1) or (1.3.2) can be interpreted as steepest descent for the (nonconvex) functional $\lambda_1 - \lambda(u)$. By analogy with the theory for strongly convex functionals, one might write down the estimates
\[
\lambda_1 - \lambda(u^k) \leqslant \rho_k (\lambda_1 - \lambda(u^0)) \tag{1.3.8}
\]
and
\[
\|u^k - u_1\|_L^2 \leqslant \rho_k \|u^0 - u_1\|_L^2, \tag{1.3.9}
\]
where $\rho_k \xrightarrow[k \to \infty]{} 0$ is independent of $u^0$, and in (1.3.9) the normalization
\[
\|u^k\|_L = \|u^0\|_L = \|u_1\|_L = 1
\]
is assumed.
% === Orig. p. 18 [s11L] ===
An inequality of the form (1.3.9) is stated without proof in \citep{CullumWilloughby1985}. Obviously, however, estimates (1.3.8) and (1.3.9) do not hold for all initial approximations; for example, when $u^0 = u_2$ we actually have $u^k = u_2$, which contradicts both (1.3.8) and (1.3.9). Moreover, in general they fail to hold simultaneously for all initial approximations $u^0$ with $\tan\ang_L(u^0; u_1) < \infty$, or even with $\lambda(u^0) > \lambda_2$. As a counterexample to (1.3.8), (1.3.9) when $2\lambda_2 > \lambda_1 + \lambda_{\min}$, a sequence $\{u_i^0\}$ is constructed in the Supplement to Chapter~1 such that $\lambda(u_i^0) > \lambda_2$, $u_i^0 \to u_2$, and for every fixed $k$ one also has $u_i^k \to u_2$ as $i \to \infty$. Oddly enough, for operators $A$ satisfying $2\lambda_2 \leqslant \lambda_1 + \lambda_{\min}$, we did not succeed in constructing such a sequence $\{u_i^0\}$. The difficulty stems from the fact that the nonlinear transition operator $u^0 \to u^k$ is discontinuous at the point $u^0 = u_2$.

We emphasize once again that the form of the domain of $u^0$ plays an essential role here: under the stronger condition $u^0 \approx u_1$, estimates (1.3.8) and (1.3.9) do hold, as follows from (1.3.3) and (1.3.4). Asymptotically (for $u^0 \approx u_1$), the analogy mentioned above is meaningful; one can even state a stronger conjecture on the asymptotic equivalence of method (1.3.1) or (1.3.2) and the method of steepest descent for minimizing the functional $r^2(u)$, which is quadratic on vectors of the form $\{u : (u, u_1)_L = 1\}$. This conjecture is based on the asymptotic relation
\[
\lambda_1 - \lambda(u) - r^2(u) \asymp \tan^4\varphi(u), \quad u \to u_1, \tag{1.3.10}
\]
(cf.\ (1.1.4)) and is corroborated, for example, by the identical asymptotic properties of the method of steepest ascent (or descent) in eigenvalue problems \citep{Zhuk1981} and steepest descent for solving linear systems, e.g., \citep{SamarskiiNikolaev1978}. A similar asymptotic equivalence is also discussed in \S\,1.5 for the conjugate gradient method.

% === Orig. p. 19 [s11R] ===
\monosection{1.4}{Lanczos method}

As is well known, e.g., \citep{Parlett1983,Ikramov1982}, the $k$th approximation $u^k$ to the vector $u_1$ in the Lanczos method is defined by the condition
\[
\begin{aligned}
\lambda(u^k) &= \max_{u \in \mathcal{K}^k(u^0)} \lambda(u), \\
\mathcal{K}^k(u^0) &\equiv \myspan\{u^0, Au^0, \ldots, A^k u^0\},
\end{aligned} \tag{1.4.1}
\]
where $\mathcal{K}^k(u^0)$ is the Krylov subspace of dimension $k + 1$. An iterative variant is also known, e.g., \citep{CullumWilloughby1985}, in which the method is terminated after a certain number of steps, an approximation to $u_1$ is computed, and this vector is then taken as a new starting vector for the Lanczos method. For instance, the method of steepest ascent is a single-step iterative Lanczos-type method.

As in the preceding section, Theorem~1.1.1 is not directly applicable. Observe, however, that the value of the Rayleigh quotient $\lambda^k = \lambda(u^k)$ in (1.4.1) is no smaller than in any variant of method (1.2.1), in particular the Chebyshev variant $\{(1.2.1),\,(1.2.7)\}$. Therefore, estimate (1.1.10) with $\bar{\sigma}^k$ from (1.2.8) holds for (1.4.1) (slightly improving the classical result of Kaniel \citep{Kaniel1966,Parlett1983}), and under condition (1.1.8), estimate (1.1.9) also holds with $\bar{\sigma}^k$ from (1.2.8) (somewhat improving the estimate of \citep{GodunovProkopov1970}). % [orig: in the citation to Kaniel and Parlett, a comma before 1983 and the closing parenthesis at the end were omitted]

It is essential that estimate (1.1.9) can be applied recursively---under condition (1.1.8), this makes it possible to prove convergence and estimate its rate for the iterative Lanczos method under any restart-frequency strategy. To replace (1.1.8) by the weaker condition (1.1.5), one needs estimates of the type (1.1.6) or (1.1.7). By analogy with the situation for the method of steepest ascent, one may conjecture that the assertions of Theorem~1.3.1 remain valid for the Lanczos method (1.4.1).

Questions concerning the sharpness of convergence rate estimates for the Lanczos method have apparently not been studied. One may expect that relation (1.3.7) remains valid if a supremum is also taken over all matrices $A$ having fixed eigenvalues $\lambda_1$, $\lambda_2$, and $\lambda_{\min}$.

% === Orig. p. 20 [s12L] ===
This would imply that for a number of problems (0.2) with $p = 1$, the Lanczos method gains only slightly in convergence rate with respect to the functional $\Delta\lambda(\cdot)$ and may even lose with respect to the functionals $\tan\varphi(\cdot)$ and $r(\cdot)$ as compared with the Chebyshev iteration method (1.2.1), (1.2.7). Nor should one forget that algorithm (1.2.1), (1.2.7) is logically simple, stable, more economical, and requires less storage even when compared with the basic Lanczos scheme without reorthogonalization. Thus, a preference for the Lanczos method (often based on comparing its convergence rate with that of simple iteration, i.e., comparing $\bar{\sigma}^k$ from (1.2.8) with (1.2.4)) does not always appear justified.

It is advisable to use combined methods. For instance, one can first apply the Lanczos method up to the onset of loss of stability, and then switch to algorithm $\{(1.2.1),\,(1.2.7)\}$ using the bounds $\lambda_2^*$, $\lambda_{\min}^*$ previously computed by the Lanczos method. The reverse order is also possible, where a few steps of the Lanczos method at the end of the iterations (1.2.7) serve to accelerate convergence and refine the bounds $\lambda_2^*$, $\lambda_{\min}^*$.

In this monograph we deliberately omit questions of practical implementation of the Lanczos method. Studying the behavior of the Lanczos method in actual finite-precision computations with rounding errors is a difficult independent problem. The original version of the method proposed by Lanczos is in a certain sense unstable (see, for example, \citep{Parlett1983,Ikramov1982}). In particular, the convergence rate estimates stated above may cease to hold. Consequently, various techniques have been developed to combat such instability, described, for instance, in \citep{CullumWilloughby1985}.

In the author's view, a variant of the Lanczos method from the book \citep{FaddeevFaddeeva1963} is of interest: one carries out a simple iteration of the form
\[
u^{n+1} = (A - \lambda_{\min} E) u^n, \quad n = 0, \ldots, k - 1,
\]
and by computing the quantities
\[
(u^n, u^n)_L, \quad (u^{n+1}, u^n)_L,
\]
accumulates the moments $((A - \lambda_{\min})^i u^0, u^0)_L$, $i = 0, \ldots, 2k$. From these moments one forms the generalized eigenvalue problem corresponding to the Rayleigh--Ritz method with coordinate vectors $u^0, \ldots, u^{k-1}$. Both matrices of this problem are Hankel and ill-conditioned. Nevertheless, it appears that $\lambda_1$ will be computed accurately and rapidly enough (for small $k$).

% === Orig. p. 21 [s12R] ===
\monosection{1.5}{Conjugate gradient method}

The functional $r^2(u)$ is quadratic and positive under the condition $\|u_1\|_L = (u, u_1)_L = 1$, and hence on the set of vectors of the form $F(A)u^0$, where $\|u_1\|_L = (u^0, u_1)_L = 1$ and $F(\lambda_1) = 1$. Consequently, the conjugate gradient method can be applied to minimize it:
\[
\begin{aligned}
& u^{n+1} = u^n - \alpha^n (\lambda_1 u^n - Au^n) - \beta^n (u^n - u^{n-1}), \\
& n = 0, 1, \ldots, k - 1,
\end{aligned} \tag{1.5.1}
\]
with parameters $\alpha^n$ and $\beta^n$ chosen to minimize $r^2(u^{n+1})$ ($\beta^0 = 0$).

The standard convergence rate estimate for the conjugate gradient method on quadratic functionals takes the following form in this case:
\[
r_k^2 \leqslant \sigma^k r_0^2, \quad \sigma^k \equiv \min \max_{i > 1} |F_k(\lambda_i)|^2 / |F_k(\lambda_1)|^2, \tag{1.5.2}
\]
where the minimum is taken over all polynomials $F_k(\lambda)$ of degree $k$. Usually, $\sigma^k$ in (1.5.2) is bounded by
\[
\sigma^k \leqslant \bar{\sigma}^k \equiv |T_k\{(1 + \xi) / (1 - \xi)\}|^{-2} \tag{1.5.3}
\]
(recall that $\xi = (\lambda_1 - \lambda_2) / (\lambda_1 - \lambda_{\min})$).

Another very simple approach is proposed in \citep{Luenberger1984} and is also used in \citep{CullumWilloughby1985}: introduce the integer $N$ such that $\lambda_N = \lambda_{\min}$, choose $F_k(\cdot)$ with roots at the points $\lambda_{N-k+1}, \ldots, \lambda_N$, and define $P_{k-1}(\lambda) \equiv F_k(\lambda) / (\lambda - \lambda_{N-k+1})$. Obviously,
\[
\max_{i > 1} |P_{k-1}(\lambda_i)| / |P_{k-1}(\lambda_1)| < 1,
\]
and therefore
\[
\sigma^k \leqslant \bar{\sigma}^k \equiv (\lambda_1 - \lambda_2)^2 / (\lambda_1 - \lambda_{N-k+1})^2. \tag{1.5.4}
\]
Estimate (1.5.4) is better than (1.5.3) if the gap $\lambda_{N-k+1} - \lambda_N$ is large. In the problems of interest to us this happens extremely rarely---for $M > 0$, clustering of the spectrum of problem (0.1) at the left end is typical.

% === Orig. p. 22 [s13L] ===
Returning to algorithm (1.5.1), note that its implementation requires knowing $\lambda_1$. This drawback is shared by all other conjugate gradient schemes (equivalent to (1.5.1)) for minimizing $r^2(\cdot)$. The traditional remedy consists in applying the conjugate gradient method to maximize the Rayleigh quotient. In particular, the counterpart of (1.5.1) is the locally optimal scheme
\[
u^{n+1} = u^n - \alpha^n (\lambda(u^n)u^n - Au^n) - \beta^n (u^n - u^{n-1}), \tag{1.5.5}
\]
which can also be written in the equivalent form (with different $\alpha^n$, $\beta^n$)
\[
u^{n+1} = Au^n - \alpha^n u^n - \beta^n u^{n-1}, \quad n = 0, 1, \ldots, k - 1, \tag{1.5.6}
\]
where in (1.5.5) and (1.5.6) the coefficients $\alpha^n$ and $\beta^n$ are chosen to maximize $\lambda(u^{n+1})$, except that $\beta^0 = 0$. Counterparts of other conjugate gradient algorithms for quadratic functionals have also been proposed, for example of the form
\[
p^n = g^n + \beta^n p^{n-1}, \quad u^{n+1} = u^n + \alpha^n p^n, \quad n = 0, 1, \ldots, k - 1, \tag{1.5.7}
\]
where $g^n = 2(Au^n - \lambda(u^n)u^n) / (u^n, u^n)_L$ is the gradient of the Rayleigh quotient at the point $u^n$. The parameters $\alpha^n$ are chosen to maximize $\lambda(u^{n+1})$. The values of $\beta^n$ (except $\beta^0 = 0$) may be selected from various considerations, for example so that $([A - \lambda(u^n)E]p^{n-1}; p^n)_L = 0$. Compared with the case of a quadratic functional, many properties of the conjugate gradient method are lost. From the theoretical standpoint, what matters most is that different variants need no longer be mutually equivalent and are apparently not equivalent to the Lanczos method. In particular, a scheme of the form (1.5.7) from \citep{CullumWilloughby1985} is not equivalent to the Lanczos method, even though $\myspan\{u^0, \ldots, u^k\}$ coincides with the Krylov subspace $\mathcal{K}^k$ from (1.4.1). One can verify this in the case of a three-dimensional space $H$ by checking that two iterations of the conjugate gradient variant in question for a suitable initial approximation $u^0 : \tan\varphi_0 < \infty$ fail to yield $u^2 = u_1$ (scheme (1.5.6) requires verification in a four-dimensional space). In all likelihood, no variant equivalent to the Lanczos method exists at all.

The lack of global optimality in the conjugate gradient method makes determining its convergence rate considerably more difficult. Nevertheless, for scheme (1.5.6) one can assert that $\lambda(u^n)$ increases no
% === Orig. p. 23 [s13R] ===
more slowly than in the method of steepest ascent (1.3.2). Consequently, under condition (1.1.8), estimate (1.3.3) holds for (1.5.6).

As $u \to u_1$, we have the asymptotic relation (1.3.10), which we reproduce here for convenience:
\[
\lambda_1 - \lambda(u) - r^2(u) \asymp \tan^4\varphi(u).
\]
It provides grounds for conjecturing the asymptotic (as $u^0 \to u_1$) equivalence of the conjugate gradient method for minimizing $r^2(\cdot)$ and any ``proper'' variant of the conjugate gradient method for maximizing $\lambda(\cdot)$, as well as, probably, the Lanczos method. One may then expect asymptotic convergence rate estimates analogous to (1.5.2), with $\sigma^k$ from (1.5.2), to hold for methods of the type (1.5.6), (1.5.7):
\[
r_k^2 / r_0^2 \leqslant \sigma^k + o(1), \qquad \Delta\lambda_k / \Delta\lambda_0 \leqslant \sigma^k + o(1). \tag{1.5.8}
\]

Note that the gradient $g^n$ of the Rayleigh quotient in method (1.5.7) is computed in the space $H_L$. Computing this gradient in the space $H$ instead yields a new formula:
\[
g^n = 2(Mu^n - \lambda(u^n)Lu^n) / (u^n, u^n)_L \tag{1.5.9}
\]
and, accordingly, a new variant of the conjugate gradient method (1.5.7).

This variant is frequently used in practice for solving generalized eigenvalue problems (see, for example, \citep{PerdonGambolati1986}), since it does not require solving linear systems of the form $Lu = f$, unlike method (1.5.7). When $g^n$ is defined by (1.5.9), the theoretical properties of method (1.5.7) (including its asymptotic properties) change substantially. The convergence rate becomes dependent on the condition number of the operator $L$: the worse the conditioning, the slower the convergence. Accordingly, to improve conditioning one can introduce an auxiliary operator $B = B^* > 0$ and compute the gradient of the Rayleigh quotient in the space $H_B$, which gives
\[
g^n = 2B^{-1}(Mu^n - \lambda(u^n)Lu^n) / (u^n, u^n)_L.
\]

Conjugate gradient methods in the space $H_B$ (and, in particular, in $H$) have been little studied and are not considered in this monograph. Some asymptotic results can be found in \citep{Savinov1981}. Other iterative methods with an auxiliary operator are studied below in Chapter~4.

% === Orig. p. 24 [s14L] ===
\monosection{1.6}{Another way of reducing the generalized eigenvalue problem to a standard one}

As shown above, when reducing (0.1) to (0.2) with
\[
A = L^{-1}M, \quad A = A^* \qquad \text{in} \qquad H_L, \tag{1.6.1}
\]
the convergence rate for most algorithms is characterized by the quantity $\bar{\sigma}^k$ with $\xi$ from (1.2.5); typical formulas are
\begin{gather*}
\bar{\sigma}^k = (1 - \xi)^{2k}, \quad \bar{\sigma}^k = [(1 - \xi) / (1 + \xi)]^{2k}, \tag{1.6.2} \\
\bar{\sigma}^k = |T_k\{(1 + \xi) / (1 - \xi)\}|^{-2}.
\end{gather*}

In computational practice, when $M > 0$, another way of reduction is also frequently used:
\[
A = -M^{-1}L, \quad A = A^* \qquad \text{in} \qquad H_M. \tag{1.6.3}
\]

Passing from (1.6.1) to (1.6.3) in the algorithm schemes of \S\S\,1.2--1.5 consists not merely in using $A = -M^{-1}L$ instead of $A = L^{-1}M$, but also requires replacing all $H_L$ inner products by $H_M$ inner products.

Two approaches to analyzing such algorithms are possible: either reintroduce the functionals in accordance with (1.6.3),
\[
\begin{aligned}
\tan^2\varphi(u) &= \tan^2\ang_M(u; u_1), \\
r^2(u) &\equiv (\lambda^{-1}(u) - \lambda_1^{-1}) / \cos^2\ang_M(u; u_1), \\
\Delta\lambda(u) &\equiv \{\lambda^{-1}(u) - \lambda_1^{-1}\} / \{\lambda_2^{-1} - \lambda^{-1}(u)\},
\end{aligned}
\]
and derive convergence rate estimates for iterations with the operator $A$ from (1.6.3) in terms of these functionals by analogy with the preceding exposition; or, without changing the terminology, take functions of the form $F_k(\lambda) = P_k(\lambda^{-1})$ as $F_k(\cdot)$ in Theorem~1.1.1, where $P_k(x)$ is a polynomial in $x$ of degree $k$.

% === Orig. p. 25 [s14R] ===
The identities
\begin{gather*}
\frac{\lambda_1 - \lambda(u)}{\cos^2\ang_L(u; u_1)} \cdot \frac{\cos^2\ang_L(v; u_1)}{\lambda_1 - \lambda(v)} = \frac{\lambda^{-1}(u) - \lambda_1^{-1}}{\cos^2\ang_M(u; u_1)} \cdot \frac{\cos^2\ang_M(v; u_1)}{\lambda^{-1}(v) - \lambda_1^{-1}}, \tag{1.6.4} \\[4pt]
\frac{\lambda_1 - \lambda(u)}{\lambda(u) - \lambda_2} \cdot \frac{\lambda(v) - \lambda_2}{\lambda_1 - \lambda(v)} = \frac{\lambda^{-1}(u) - \lambda_1^{-1}}{\lambda_2^{-1} - \lambda^{-1}(u)} \cdot \frac{\lambda_2^{-1} - \lambda^{-1}(v)}{\lambda^{-1}(v) - \lambda_1^{-1}} \tag{1.6.5}
\end{gather*}
show that the ratios $r^2(u)/r^2(v)$ and $\Delta\lambda(u)/\Delta\lambda(v)$ remain invariant upon passing from the old definitions to the new ones. Hence both approaches lead to equivalent results with respect to the functionals $r^2(\cdot)$ and $\Delta\lambda(\cdot)$. The functional $\tan^2\varphi$ does not possess an invariance property of the form (1.6.4), (1.6.5), so passing to the new definitions also yields a new estimate: if for $M > 0$ the condition
\[
\tan^2\ang_M(u^0; u_1) < \infty \tag{1.6.6}
\]
holds, then we have the estimate
\[
\tan^2\ang_M(u^k; u_1) \leqslant \bar{\sigma}^k \tan^2\ang_M(u^0; u_1). \tag{1.6.7}
\]
Recall that in the analogous assertion (1.1.6) of Theorem~1.1.1, the angle is measured in $H_L$.

It is useful to know the relation between angles in $H_M$ and $H_L$:
\[
\tan^2\ang_M(u; u_1) / \tan^2\ang_L(u; u_1) \in [\lambda_{\min}/\lambda_1, \lambda_2/\lambda_1], \tag{1.6.8}
\]
which reveals the equivalence of conditions (1.1.5) and (1.6.6).

% === Orig. p. 26 [s15L] ===
Let us now proceed to the actual study of the algorithms of \S\S\,1.2--1.5 under substitution (1.6.3). We state a general rule for obtaining convergence rate estimates and then give examples.

If for one of the methods with $A = L^{-1}M$ listed above in \S\S\,1.2--1.5 a convergence rate estimate holds with $\bar{\sigma}^k$ of the type (1.6.2) and $\xi$ from (1.2.5), then for the counterpart of that method under substitution (1.6.3) the same estimate holds, but with $\xi$ given by
\[
\begin{aligned}
\xi &= (\lambda_2^{-1} - \lambda_1^{-1}) / (\lambda_{\min}^{-1} - \lambda_1^{-1}) \\
&= \lambda_{\min} / \lambda_2 \, (\lambda_1 - \lambda_2) / (\lambda_1 - \lambda_{\min}).
\end{aligned} \tag{1.6.9}
\]

As can be seen from (1.6.9), the new value of $\xi$ differs from the old $\xi$ in (1.2.5) by the extra factor $\lambda_{\min}/\lambda_2 < 1$. For mesh problems (0.1), this quantity is usually not merely small, but also decreases rapidly under mesh refinement because $\lambda_{\min} \to 0$. Consequently, using substitution (1.6.3) as compared with (1.6.1) leads to a deterioration of the convergence rate estimates, which is especially pronounced on fine meshes.

Let us now give several examples.

In the counterpart of simple iteration (1.2.1) of the form
\[
u^{n+1} = M^{-1}Lu^n - \gamma^n u^n, \quad n = 0, 1, \ldots, k - 1, \tag{1.6.10}
\]
with a constant parameter
\[
\gamma^n = \gamma > (\lambda_1^{-1} + \lambda_{\min}^{-1}) / 2, \tag{1.6.11}
\]
the optimal parameter is
\[
\gamma^n = \gamma = (\lambda_2^{-1} + \lambda_{\min}^{-1}) / 2. \tag{1.6.12}
\]

Under the corresponding conditions on $u^0$, all the estimates of Theorem~1.1.1, corollary (1.1.10), and estimate (1.6.7) hold with
% === Orig. p. 27 [s15R] ===
\[
\bar{\sigma}^k = [\max\{|\lambda_{\min}^{-1} - \gamma|; |\gamma - \lambda_2^{-1}|\} / |\gamma - \lambda_1^{-1}|]^{2k}. \tag{1.6.13}
\]
For the optimal choice of $\gamma^n$ from (1.6.12), the expression for $\bar{\sigma}^k$ in (1.6.13) reduces to the equality
\[
\bar{\sigma}^k = [(1 - \xi) / (1 + \xi)]^{2k}, \tag{1.6.14}
\]
where $\xi$ is defined in (1.6.9).

As we see, formulas (1.2.6) and (1.6.14) coincide in form, except that $\xi$ in (1.2.6) is defined by (1.2.5), whereas in (1.6.14) it is defined by (1.6.9).

In the case of variable shifts $\gamma^n$ in (1.6.10), the optimal a priori choice is provided by the Chebyshev parameters (cf.\ (1.2.7))
\[
\gamma^n = \{\lambda_{\min}^{-1} + \lambda_2^{-1} - \tau^n (\lambda_{\min}^{-1} - \lambda_2^{-1})\}, \tag{1.6.15} % [math note: the factor 1/2 in front of the braces is omitted in the original, cf. (1.2.7) and (1.6.12)]
\]
where $\tau^n \in [-1, 1]$ are all $k$ roots of the Chebyshev polynomial of the first kind $T_k(\tau)$ of degree $k$. Then in the estimates of Theorem~1.1.1, in corollary (1.1.10), and in estimate (1.6.7) we have
\[
\bar{\sigma}^k = |T_k\{(1 + \xi) / (1 - \xi)\}|^{-2}, \tag{1.6.16}
\]
just as in (1.2.8), but with $\xi$ from (1.6.9).

The counterpart of the method of steepest ascent (1.3.1) takes the form
\[
\begin{aligned}
& u^{n+1} = M^{-1}Lu^n - \alpha^n u^n, \quad n = 0, 1, \ldots, k - 1, \\
& \alpha^n \equiv \frac{1}{2}\left\{\lambda^{-1}(w^n) + \lambda^{-1}(u^n) + \sqrt{[\lambda^{-1}(w^n) - \lambda^{-1}(u^n)]^2 + 4\|w^n\|_M^2 \|u^n\|_M^2}\right\}, \\ % [math note: the original has \|u^n\|_M^2 instead of \|u^n\|_M^{-2}, cf. (1.3.1)]
& w^n \equiv M^{-1}Lu^n - \lambda^{-1}(u^n)u^n.
\end{aligned} \tag{1.6.17}
\]
% === Orig. p. 28 [s16L] ===
All the results of \S\,1.3 hold for this scheme as well, with $\xi$ from (1.2.5) replaced by $\xi$ from (1.6.9). Note that $\lambda^{-1}(\cdot)$ is the Rayleigh quotient for the operator $M^{-1}L$ in $H_M$, and method (1.6.17) actually performs steepest descent for this Rayleigh quotient.

In a similar manner, one can write down the schemes of the counterparts of the other algorithms from \S\S\,1.2--1.5 and determine their convergence rates.

New possibilities arise when $M > 0$ from using the substitution
\[
A = -M^{-1}L, \quad A = A^* \qquad \text{in} \qquad H_L. \tag{1.6.18}
\]

Compared with (1.6.3), the inner product of $H_L$ is used here instead of that of $H_M$. Naturally, methods with an a priori choice of parameters, such as $\{(1.6.10);\,(1.6.12)\}$, remain unchanged, since their schemes involve no inner products.

For variational-type methods under substitution (1.6.3)---steepest ascent (descent) (1.6.17), the Lanczos method, and the conjugate gradient method---their counterparts under substitution (1.6.18) converge at the same rate, although the resulting algorithm schemes have a distinctive form. For example, instead of the functional
\[
\lambda^{-1}(\cdot) = (M^{-1}L\,\cdot, \cdot)_M / (\cdot, \cdot)_M
\]
one encounters the somewhat unusual functional
\[
(M^{-1}L\,\cdot, \cdot)_L / (\cdot, \cdot)_L.
\]
Below, in \S\,2.2 and in Chapter~3, we investigate the effect of such a change of metric in the Rayleigh--Ritz method.

Let us briefly touch upon the computational aspect of reducing (0.1) to (0.2). At each iteration of the methods in this chapter, one has to solve linear systems of the form $Lu = f$ in the case of (1.6.1), or $Mu = f$ in the cases of (1.6.3) and (1.6.18). For many practical problems (0.1), these systems can only be solved approximately. How the accuracy of this inner approximation affects the convergence rate of the methods has been little studied in general. If iterations of the form
\[
u^{n+1} = u^n - \gamma^n B^{-1}(Lu^n - f)
\]
with some auxiliary operator $B$ are employed to solve the systems $Lu = f$, then in a number of cases the resulting method for finding the eigenvector $u_1$ coincides with one of the methods studied below in Chapter~4.

% === Orig. p. 29 [s16R] ===
\monosection{1.7}{Inverse iteration with shifts}

We consider only one class of variants of the method, in which the Rayleigh quotient is used as the shift:
\[
(M - \lambda^n L)v^n = Lu^n, \quad u^{n+1} = v^{n+1} - \gamma^n u^n, \quad n = 0, 1, \ldots \tag{1.7.1} % [math note: the first equation in the original has v^n instead of v^{n+1}, cf. (1.7.7)]
\]
For the traditional choice $\gamma^n = 0$ \citep{Parlett1983}, we obtain from Theorem~1.1.1 that for $\lambda^n > (\lambda_1 + \lambda_2) / 2$ the following estimates hold:
\[
\begin{aligned}
& \tan^2\varphi_{n+1} / \tan^2\varphi_n \leqslant \bar{\sigma}^n, \\
& r_{n+1}^2 / r_n^2 \leqslant \bar{\sigma}^n, \qquad \Delta\lambda_{n+1} / \Delta\lambda_n \leqslant \bar{\sigma}^n,
\end{aligned} \tag{1.7.2}
\]
where
\[
\bar{\sigma}^n \equiv [(\lambda_1 - \lambda^n) / (\lambda^n - \lambda_2)]^2. \tag{1.7.3}
\]
In the case of a nonzero parameter $\gamma^n$, by analogy with (1.2.1)--(1.2.6), when $\lambda^n > \lambda_2$ and $\gamma^n < \frac{1}{2}[(\lambda_1 - \lambda^n)^{-1} - (\lambda^n - \lambda_2)^{-1}]$, estimates (1.7.2) hold with
\[
\begin{aligned}
\bar{\sigma}^n &\equiv \max\left\{|(\lambda^n - \lambda_{\min})^{-1} + \gamma^n|^2; |(\lambda^n - \lambda_2)^{-1} + \gamma^n|^2\right\} \\
&\quad \cdot \{(\lambda_1 - \lambda^n)^{-1} - \gamma^n\}^{-2}.
\end{aligned} \tag{1.7.4}
\]

For the optimal choice of the parameter $\gamma^n$,
\[
\gamma^n = -[(\lambda^n - \lambda_{\min})^{-1} + (\lambda^n - \lambda_2)^{-1}] / 2, \tag{1.7.5}
\]
the expression for $\bar{\sigma}^n$ in (1.7.4) reduces to the equality
\[
\bar{\sigma}^n = [(1 - \xi_n) / (1 + \xi_n)]^2, \tag{1.7.6}
\]
% === Orig. p. 30 [s17L] ===
where $\xi_n \equiv (\lambda^n - \lambda_2)(\lambda_1 - \lambda_{\min})(\lambda^n - \lambda_{\min})^{-1}(\lambda_1 - \lambda_2)^{-1}$.

When $M > 0$, for the scheme analogous to (1.7.1),
\[
(M - \lambda^n L)v^{n+1} = Mu^n, \quad u^{n+1} = v^{n+1} - \gamma^n u^n, \quad n = 0, 1, \ldots, \tag{1.7.7}
\]
with $\gamma^n = 0$ and $\lambda^n > (\lambda_1 + \lambda_2) / 2$, estimates (1.7.2) also hold, with $\bar{\sigma}^n$ of the form
\[
\bar{\sigma}^n = \lambda_2^2 \lambda_1^{-2} [(\lambda_1 - \lambda^n) / (\lambda^n - \lambda_2)]^2. \tag{1.7.8}
\]
The value of $\bar{\sigma}^n$ in (1.7.8) differs from (1.7.3) by the extra factor $\lambda_2^2 / \lambda_1^2 < 1$. Consequently, when $\gamma^n = 0$, the convergence rate estimates for algorithm (1.7.7) are somewhat better than those for (1.7.1).

For the optimal choice of the parameter in (1.7.7),
\[
\gamma^n = -[(\lambda^n / \lambda_{\min} - 1)^{-1} + (\lambda^n / \lambda_2 - 1)^{-1}] / 2,
\]
method (1.7.7) naturally becomes equivalent to $\{(1.7.1),\,(1.7.5)\}$, and for $\lambda^n > \lambda_2$ estimates (1.7.2) with $\bar{\sigma}^n$ from (1.7.6) hold for it as well.

Summarizing the above, we conclude that inverse iterations with the Rayleigh quotient shift converge at least as fast as a geometric progression, and as $n \to \infty$ the convergence is cubic. Under mesh refinement in mesh problems (0.1), the convergence does not deteriorate---just as in the methods with $A = L^{-1}M$ analyzed earlier in \S\S\,1.1--1.5---since the quantity
\[
(\lambda_1 - \lambda_2) / (\lambda_1 - \lambda_{\min})
\]
in the problems of interest to us is bounded away from zero by a constant independent of the mesh size.

In practice, solving the linear systems $(M - \lambda^n L)u = f$ exactly often requires prohibitively high computational cost. When these systems are solved approximately, the difficult question arises of analyzing the effect of the resulting errors on the course of the iterations; for instance, cubic convergence may be lost.

% =============================================================================
% SOURCE MODULE: part_02_supp1.tex
% =============================================================================
% === Orig. p. 31 [s17R] ===
\monosuppchapter{Supplement to Chapter 1}{\texorpdfstring{A.\,V.~Knyazev, A.\,L.~Skorokhodov}{A. V. Knyazev, A. L. Skorokhodov}}{Derivation of Sharp Convergence Rate Estimates for the Method of Steepest Ascent}

For the reader's convenience, this supplement is organized as an independent, self-contained section. The necessary notation is reintroduced and the results from \S\,1.3 are restated.

In \S\,S.1, convergence rate estimates for the method in various terms and assertions on their attainability are presented.

In \S\,S.2, the results obtained are compared with known ones.

The concluding \S\,S.3 is devoted to the proofs.

Let us introduce the necessary notation. Let $H$ be a Euclidean space with inner product $(\cdot, *)$ and norm $\|\cdot\| = (\cdot, \cdot)^{1/2}$; let $A$ be a self-adjoint linear operator in $H$ (here, to simplify notation, we do not use the conventions adopted in Chapter~1 regarding the form $A = L^{-1}M$ in $H_L$ or $A = M^{-1}L$ in $H_M$ for $M > 0$, etc., and treat $A$ as an arbitrary operator: $A = A^*$ in $H$. This entails no loss of generality);

$\lambda_1 > \lambda_2 \geqslant \dots \geqslant \lambda_{\min}$ are the eigenvalues (without loss of generality, $\lambda_1$ may be assumed simple);

$u_1, u_2, \dots, u_{\min}$ are the corresponding orthonormal eigenvectors of the operator $A$, forming a basis of $H$.

% === Orig. p. 32 [s18L] ===
\monosection{S.1}{Convergence rate estimates}

In the method of steepest ascent \citep{FaddeevFaddeeva1963,KantorovichAkilov1984}, for a given initial vector $u^0 \neq 0$ we compute
\[
u^{n+1} = u^n + \gamma_n w^n, \qquad n = 0, 1, \dots, \tag{S.1.1}
\]
where
\begin{gather*}
w^n \equiv A u^n - \lambda(u^n) u^n, \tag{S.1.2} \\
\gamma_n = 2 \left\{ \lambda(u^n) - \lambda(w^n) + \sqrt{[\lambda(u^n) - \lambda(w^n)]^2 + 4 \|w^n\|^2 \|u^n\|^{-2}} \right\}^{-1} \tag{S.1.3}
\end{gather*}
and $\lambda(\cdot) \equiv (A\cdot, \cdot) / (\cdot, \cdot)$ is the Rayleigh quotient.

In a practical implementation of the method, normalization of the vectors $u^n$ should be added to formulas (S.1.1)--(S.1.3). The parameter $\gamma_n$ in (S.1.3) is chosen from the condition of maximizing $\lambda(u^{n+1})$.

As a measure of proximity of a vector $v \in H \setminus 0$ to $u_1$, we use the functionals (see \S\,1.1)
\begin{align*}
&r^2(v) \equiv (\lambda_1 - \lambda(v)) / \cos^2 \ang(v; u_1), \\
&\tan^2 \ang(v, u_1), \\
&\Delta\lambda(v) \equiv (\lambda_1 - \lambda(v)) / (\lambda(v) - \lambda_2).
\end{align*}

Let us denote $r_n^2 \equiv r^2(u^n)$, $\Delta\lambda_n \equiv \Delta\lambda(u^n)$, $\varphi_n \equiv \ang(u^n; u_1)$, $n = 0, k$, and formulate the convergence rate estimates in these terms.

\textbf{Theorem S.1.1.} The following inequalities hold (see Theorem~1.3.1):
\begin{gather*}
r_k^2 / r_0^2 \leqslant (1 - \xi)^{2k}, \tag{S.1.4} \\
% === Orig. p. 33 [s18R] ===
\tan^2 \varphi_k / \tan^2 \varphi_0 \leqslant (1 - \xi)^{2k}, \qquad k = 1, 2, \dots, \tag{S.1.5}
\end{gather*}
where $\xi \equiv (\lambda_1 - \lambda_2) / (\lambda_1 - \lambda_{\min})$.

Estimates (S.1.4) and (S.1.5) are attainable in the sense that for each $n = 0, 1, \dots$ the following assertions hold:
\[
\sup_{u^n} r_{n+1}^2 / r_n^2 = \sup_{u^n} \tan^2 \varphi_{n+1} / \tan^2 \varphi_n = (1 - \xi)^2. \tag{S.1.6}
\]

Asymptotically (as $u^0 \to u_1$), estimate (S.1.4) can be improved:
\[
r_k^2 \leqslant \sigma^k r_0^2, \qquad \sigma^k \equiv [(1 - \xi) / (1 + \xi)]^2 + o(1). \tag{S.1.7} % [math note: in original, \sigma^k \equiv [(1-\xi)/(1+\xi)]^2 + o(1); cf. the proof on orig. p. 44, where r_{n+1}^2 \leqslant \sigma_1 r_n^2 with \sigma_1 \equiv (1-\xi)^2/(1+\xi)^2 + o(1)]
\]

The following theorem holds (see Theorem~1.1.1 % [orig: I.I.I.]
and inequality~(1.3.3)).

\textbf{Theorem S.1.2.} If $\lambda(u^0) > \lambda_2$, then $\Delta\lambda_k \geqslant 0$ and
\[
\Delta\lambda_k / \Delta\lambda_0 \leqslant [(1 - \xi) / (1 + \xi)]^{2k}. \tag{S.1.8}
\]

Estimate (S.1.8) is attainable in the sense that
\[
\lim_{\varepsilon \to 0} \sup_{u^0 : \varphi_0 \leqslant \varepsilon} \Delta\lambda_k / \Delta\lambda_0 = [(1 - \xi) / (1 + \xi)]^{2k}, \tag{S.1.9}
\]
whence, for example, follows an attainability assertion of the type (S.1.6):
\[
\sup_{u^n : \Delta\lambda_n > 0} \Delta\lambda_{n+1} / \Delta\lambda_n = [(1 - \xi) / (1 + \xi)]. % [math note: in original, the exponent 2 is omitted on the right-hand side; cf. (S.1.8), (S.1.9), and orig. p. 46]
\]

% === Orig. p. 34 [s19L] ===
\monosection{S.2}{Comparison with known results}

The convergence of $u^k$ to $u_1$ for $\tan\varphi_0 < \infty$ is well known \citep{FaddeevFaddeeva1963,KantorovichAkilov1984}, but convergence rate estimates in such a general setting have apparently been obtained for the first time. Note that the condition $\tan\varphi_0 < \infty$, or the equivalent condition $r_0 < \infty$, is also necessary for $u^k \to u_1$ as $k \to \infty$.

A result close to (S.1.8) was published earlier in \citep{Prikazchikov1975}. Let us compare them.

In our notation, V.\,G.~Prikazchikov's estimate takes the form:
\[
\begin{aligned}
&\lambda_1 - \lambda(u^1) \leqslant \bar{q}(\lambda_1 - \lambda(u^0)), \\
&\bar{q} \equiv \left( \frac{\lambda_2 - \lambda_{\min} + \lambda_1 - \lambda(u^0)}{2\lambda(u^0) - \lambda_2 - \lambda_{\min} + \lambda_1 - \lambda(u^0)} \right)^2.
\end{aligned}
\tag{S.2.1}
\]
Like (S.1.8), it holds for $\lambda(u^0) > \lambda_2$. By equivalent transformations, inequality (S.1.8) can be brought to a similar form:
\[
\begin{aligned}
&\lambda_1 - \lambda(u^1) \leqslant q(\lambda_1 - \lambda(u^0)), \\
&q \equiv \frac{(\lambda_2 - \lambda_{\min})^2}{(2\lambda(u^0) - \lambda_2 - \lambda_{\min})^2 + 4(\lambda_1 - \lambda(u^0))(\lambda(u^0) - \lambda_2)}.
\end{aligned}
\tag{S.2.2}
\]
Now observe that
\[
q \leqslant \left( \frac{\lambda_2 - \lambda_{\min}}{2\lambda(u^0) - \lambda_2 - \lambda_{\min}} \right)^2 \leqslant \bar{q},
\]
where the left inequality is obvious, and the right one holds since the following function is increasing in $x \equiv \lambda_1 - \lambda(u^0) \geqslant 0$:
\[
(\lambda_2 - \lambda_{\min} + x) / (2\lambda(u^0) - \lambda_2 - \lambda_{\min} + x).
\]

Consequently, estimate (S.1.8) is better than (S.2.1), while as $u^0 \to u_1$ we have $\bar{q} \approx q$, i.e., asymptotically the two estimates coincide.

% === Orig. p. 35 [s19R] ===
Note that as $u^0 \to u_1$, estimates (S.2.1) and (S.2.2) agree with the classical result \citep{KantorovichAkilov1984} on the asymptotic convergence rate of the method of steepest ascent.

Finally, let us write out the estimate from \citep{CullumWilloughby1985}, specifying inequality (1.3.9):
\[
\|u^k - u_1\|^2 \leqslant c q^k \|u^0 - u_1\|^2, \tag{S.2.3}
\]
where
\[
c \equiv \frac{\lambda_1 - \lambda_{\min}}{\lambda_1 - \lambda_2}, \qquad q \equiv \left( \frac{1 - \xi}{1 + \xi} \right)^2
\]
and
\[
\|u^k\| = \|u^0\| = 1.
\]

As will be shown in the next \S\,S.3, this estimate is, generally speaking, invalid, just like any estimate of the form (1.3.8) or (1.3.9). To see this, at the end of \S\,S.3 we construct a family of initial vectors $\{u^0\}$ such that $u^0 \to u_2$, $\lambda(u^0) > \lambda_2$, and $u^k \to u_2$ for any fixed $k$. Leaving aside the vacuous case $c q^k \geqslant 1$ and assuming that $c q^k < 1$, for $\{u^k\}$ we have:
\[
\|u^k - u_1\|^2 \to 2 \quad \text{and} \quad \lambda_1 - \lambda(u^k) \to \lambda_1 - \lambda_2,
\]
which contradicts (S.2.3) and (1.3.8), respectively, since $c, q$ in (S.2.3) and $\rho_k$ in (1.3.8) do not depend on $u^n$, $n = 0, k$.

Under the stronger condition $u^0 \to u_1$ (stronger than $\lambda(u^0) > \lambda_2$), estimate (S.2.3) becomes valid with $c = 1 + o(1)$ and $q = (1 - \xi)^2$. This follows from inequality (S.1.5) and the fact that
\[
\|u^n - u_1\|^2 \sim \tan^2 \ang(u^n; u_1), \qquad n = 0, k,
\]
as $u^0 \to u_1$. Thus, the form of the domain of variation of $u^0$ plays a decisive role here.

% === Orig. p. 36 [s20L] ===
\monosection{S.3}{Proofs}

\textit{Proofs of Theorems S.1.1 and S.1.2.} Since $\xi$ does not depend on the index $k$, $k \geqslant 1$, it suffices to prove the theorems for the case $k = 1$.

Consider the subspaces
\[
\mathcal{K} = \myspan\{u^0, A u^0\}, \qquad H^3 = \myspan\{u_1, u^0, A u^0\}, \qquad \mathcal{K} \subseteq H^3.
\]

If $\dim \mathcal{K} = 1$, then $u^0$ is an eigenvector of the operator $A$. Then $w^1 = 0$, and formula (S.1.3) is meaningless since the value $\lambda(w^1)$ is undefined. % [math note: in original, w^1 and \lambda(w^1) instead of w^0 and \lambda(w^0)]
Therefore, we assume below that
\[
\dim \mathcal{K} = 2, \qquad \mathcal{K} = \myspan\{u^0, u^1\}.
\]

If $\dim H^3 = 2$, then $u_1 \in \mathcal{K}$. Then $u^1 = u_1$, and the left-hand sides of inequalities (S.1.4), (S.1.5), and (S.1.8) vanish, i.e., the estimates become trivial. Thus, the nontrivial case is
\[
\dim H^3 = 3, \qquad H^3 = \myspan\{u_1, u^0, u^1\}.
\]

Let us introduce the orthogonal projector $P^3$ onto $H^3$. Then the projection of the operator $A$ onto $H^3$---the operator $A^3 \equiv P^3 A$---is symmetric on $H^3$, and $H^3$ is its invariant subspace.

Let us state the following simple yet important lemma.

\textbf{Lemma S.3.1.} If in the formulas of the method (S.1.1)--(S.1.3) the operator $A$ is replaced by $A^3$, then the vector $u^1$ remains unchanged.

\textit{Proof of Lemma S.3.1.} Since
\[
\mathcal{K} \equiv \myspan\{u^0, A u^0\} = \myspan\{u^0, A^3 u^0\}
\]
and the Rayleigh quotients for the operators $A$ and $A^3$ coincide on the subspace $H^3$, the vector $u^1$ can be found by the method of steepest ascent with the operator $A^3$.

Lemma~S.3.1 states that for $k = 1$ we can in fact analyze the method (S.1.1)--(S.1.3) in a three-dimensional space. A similar technique of reduction to a ``low-dimensional'' case was proposed in \citep{Knyazev1985b,Knyazev1985c}.

% === Orig. p. 37 [s20R] ===
Let $\lambda_1^3 \geqslant \lambda_2^3 \geqslant \lambda_3^3$ denote the Ritz values obtained by the Rayleigh--Ritz method in the subspace $H^3$: $u_1^3, u_2^3, u_3^3$ are the corresponding orthonormal Ritz vectors. Then $\lambda_i^3$ are the eigenvalues and $u_i^3$ are the eigenvectors of the operator $A^3$.

Since $u_1 \in H^3$, we have $u_1^3 = u_1$ and $\lambda_1^3 = \lambda_1$. By the assumption $\lambda_1 > \lambda_2$ and the Courant--Fischer variational principle, the following inequalities hold:
\[
\lambda_1^3 = \lambda_1 > \lambda_2 \geqslant \lambda_2^3 \geqslant \lambda_3^3 \geqslant \lambda_{\min}. \tag{S.3.1}
\]

Let us show that the case $\lambda_2^3 = \lambda_3^3$ is trivial. Represent the vector $u^0$ in the form
\[
u^0 = \alpha u_1 + \beta v, \qquad \text{where } v \in \myspan\{u_2^3, u_3^3\}.
\]
Since $\dim \mathcal{K} = 2$, we have $\beta \neq 0$. If $\lambda_2^3 = \lambda_3^3$, then $v$ is an eigenvector of $A^3$, and hence $\alpha \neq 0$. Then $\mathcal{K} = \myspan\{u_1, v\}$, so that $u^1 = u_1$, and the estimates become trivial. Henceforth we assume that $\lambda_2^3 > \lambda_3^3$.

Let us show that to prove Theorems~S.1.1 and~S.1.2 for $k = 1$ it suffices to verify their analogues in the subspace $H^3$:
\begin{gather*}
r_1^2 / r_0^2 \leqslant \lambda^2, \tag{S.3.2} \\
\tan^2 \varphi_1 / \tan^2 \varphi_0 \leqslant \lambda^2, \tag{S.3.3} \\
\Delta\lambda_1^3 / \Delta\lambda_0^3 \leqslant \lambda^2 / (2 - \lambda)^2, \tag{S.3.4}
\end{gather*}
where
\begin{align*}
&r_n^2 \equiv \lambda_1^3 - \lambda(u^n) / \cos^2 \varphi_n, \qquad \varphi_n \equiv \ang(u^n, u_1^3), \\ % [math note: in original, parentheses around \lambda_1^3 - \lambda(u^n) are omitted]
&\Delta\lambda_n^3 \equiv (\lambda_1^3 - \lambda(u^n)) / (\lambda(u^n) - \lambda_2^3), \qquad n = 0, 1, \\
&\lambda \equiv (\lambda_2^3 - \lambda_3^3) / (\lambda_1^3 - \lambda_3^3),
\end{align*}
% === Orig. p. 38 [s21L] ===
and the vector $u^1$ is obtained from $u^0$ by a single iteration of the method of steepest ascent with the operator $A^3$. Indeed, we have $\lambda_1^3 = \lambda_1$, $u_1^3 = u_1$, and invoking Lemma~S.3.1, we conclude that the left-hand sides in inequalities (S.1.4), (S.1.5) and, respectively, (S.3.2), (S.3.3) coincide. From inequalities (S.3.1) and the fact that $\lambda(u^1) \geqslant \lambda(u^0)$, it follows that
\[
\Delta\lambda_1 / \Delta\lambda_0 \leqslant \Delta\lambda_1^3 / \Delta\lambda_0^3,
\]
since the function $\{(\lambda(u^0) - \lambda_2) / (\lambda(u^1) - \lambda_2)\}$ is nonincreasing in the parameter $\lambda_2 \geqslant \lambda_2^3$. Finally, by virtue of (S.3.1) the inequality $\lambda \leqslant 1 - \xi$ holds, and consequently the right-hand sides of inequalities (S.3.2)--(S.3.4) do not exceed the right-hand sides of the corresponding inequalities.

Thus, with our special choice of $H^3$, the inequalities for the ``three-dimensional'' case (S.3.2)--(S.3.4) are no weaker than the corresponding ones for the ``multidimensional'' case---(S.1.4), (S.1.5), and (S.1.8). In other words, if Theorems~1 and~2 % [orig: Theorems 1 and 2, i.e., Theorems S.1.1 and S.1.2]
hold for $\dim H = 3$, then they also hold for $\dim H \geqslant 3$.

Further in the proof we deal only with the ``three-dimensional'' case, so we drop the superscript $3$ for brevity.

The validity of inequalities (S.3.2)--(S.3.4) can be verified in several ways. The most direct one is to expand the vector $u_0$ % [math note: in original, u_0 instead of u^0]
in the basis of vectors $u_1, u_2, u_3$ and express the left-hand sides of inequalities (S.3.2)--(S.3.4) in terms of the coefficients of this expansion. Unfortunately, this leads to cumbersome algebraic expressions. We have preferred to write the left-hand sides of (S.3.2)--(S.3.4) in terms of the numbers $\lambda_1, \lambda_2, \lambda_3$ and the auxiliary quantities $\tilde{\lambda}_1, \tilde{\lambda}_2$, where $\tilde{\lambda}_1 \geqslant \tilde{\lambda}_2$ are the Ritz values obtained by the Rayleigh--Ritz method in the subspace $\mathcal{K}$. Note that $\lambda(u^1) = \tilde{\lambda}_1$, i.e., $u^1$ is a Ritz vector in $\mathcal{K}$ corresponding to $\tilde{\lambda}_1$.

Let us write down the inequalities following from the variational principle:
\[
\lambda_1 \geqslant \tilde{\lambda}_1 \geqslant \lambda_2 \geqslant \tilde{\lambda}_2 \geqslant \lambda_3
\]
and show that only the case where all inequalities are strict is nontrivial. The following lemma holds.

\textbf{Lemma S.3.2.} Represent $u^0$ in the form
\[
u^0 = \sum_{i=1}^3 a_i u_i, \qquad \sum_{i=1}^3 a_i^2 \neq 0.
\]
% === Orig. p. 39 [s21R] ===
Then, if at least one $\tilde{\lambda}_i$, $i = 1, 2$, coincides with some $\lambda_j$, $j = 1, 2, 3$, then $a_1 a_2 a_3 = 0$.

\textbf{Corollary.} If the hypotheses of the lemma hold, then estimates (S.3.2)--(S.3.4) become trivial.

\textit{Proof of Lemma S.3.2.} Let us introduce the moments
\[
m_l = (A^l u^0, u^0); \qquad l = 0, \dots, 3, \qquad m_l = \sum_{i=1}^3 (\lambda_i)^l a_i^2.
\]
Then it follows from Lemma~S.3.1 that the equation for determining $\tilde{\lambda}_1$ and $\tilde{\lambda}_2$ by the Rayleigh--Ritz method has the form:
\[
\det \begin{bmatrix}
m_1 - \tilde{\lambda} m_0 & m_2 - \tilde{\lambda} m_1 \\
m_2 - \tilde{\lambda} m_1 & m_3 - \lambda m_2
\end{bmatrix} = 0. % [math note: in original, the bottom-right matrix entry is m_3 - \lambda m_2 instead of m_3 - \tilde{\lambda} m_2]
\]
Performing the necessary manipulations, we obtain the expressions for $\tilde{\lambda}_1 + \tilde{\lambda}_2$ and $\tilde{\lambda}_1 \cdot \tilde{\lambda}_2$:
\begin{align*}
\tilde{\lambda}_1 + \tilde{\lambda}_2 &= \frac{\sum_{i,j=1}^3 (\lambda_i - \lambda_j)^2 (\lambda_i + \lambda_j) a_i^2 a_j^2}{\sum_{i,j=1}^3 (\lambda_i - \lambda_j)^2 a_i^2 a_j^2}, \\[1ex]
\tilde{\lambda}_1 \cdot \tilde{\lambda}_2 &= \frac{\sum_{i,j=1}^3 (\lambda_i - \lambda_j)^2 \lambda_i \lambda_j a_i^2 a_j^2}{\sum_{i,j=1}^3 (\lambda_i - \lambda_j)^2 a_i^2 a_j^2}.
\end{align*}

Consider the following cases:

1. Let $\tilde{\lambda}_2 = \lambda_3$. Then
\[
(\lambda_3^2 - (\lambda_1 + \lambda_2)\lambda_3 + \lambda_1 \lambda_2) a_1^2 a_2^2 = 0,
\]
and since $\lambda_2 > \lambda_3$, we have $a_1 a_2 = 0$.

2. Similarly: if $\tilde{\lambda}_2 = \lambda_2$, then $a_1 a_3 = 0$; if $\tilde{\lambda}_1 = \lambda_1$, then $a_2 a_3 = 0$; if $\tilde{\lambda}_1 = \lambda_2$, then $a_1 a_3 = 0$. The lemma is proved.

\textit{Proof of the Corollary.} If $a_1 = 0$, then $u^0 \perp u_1$, and the estimates are vacuous; if $a_2 a_3 = 0$, then $u_1 \in \mathcal{K}$, $u^1 = u_1$, and the estimates are trivial.

% === Orig. p. 40 [s22L] ===
\textbf{Remark.} Lemma~S.3.2 is a special case of the following assertion.

Suppose that the dimension of the Krylov subspace $\mathcal{K}(A, u)$ \citep{Parlett1983} is one less than the dimension of the entire space. Carry out the Rayleigh--Ritz procedure in $\mathcal{K}(A, u)$ and determine the Ritz values. Then, if at least one of them coincides with some eigenvalue of the operator $A$, the subspace $\mathcal{K}(A, u)$ is invariant under $A$, i.e., in the expansion of the vector $u$ in the basis of eigenvectors of the operator $A$, exactly one term is absent (see Lemma~2.2.3 below).

Consider the nontrivial case: $\lambda_1 > \tilde{\lambda}_1 > \lambda_2 > \tilde{\lambda}_2 > \lambda_3$. Represent the vectors $u^0$ and $u^1$ in the form:
\begin{align*}
u^0 &= u_1 + \alpha_0 u_2 + \beta_0 u_3, \\
u^1 &= u_1 + \alpha_1 u_2 + \beta_1 u_3.
\end{align*}
Then
\begin{align*}
&r_i^2 = (\lambda_1 - \lambda_2)\alpha_i^2 + (\lambda_1 - \lambda_3)\beta_i^2, \\
&\tan^2 \varphi_i = \alpha_i^2 + \beta_i^2, \qquad i = 0, 1, \\[1ex]
&\Delta\lambda_0 = \frac{(\lambda_1 - \lambda_2)\alpha_0^2 + (\lambda_1 - \lambda_3)\beta_0^2}{\lambda_1 - \lambda_2 - (\lambda_2 - \lambda_3)\beta_0^2}, \\[1ex]
&\Delta\lambda_1 = (\lambda_1 - \tilde{\lambda}_1) / (\tilde{\lambda}_1 - \lambda_2).
\end{align*}

To pass completely to the variables $\lambda_i, \tilde{\lambda}_i$, it remains to express the quantities $\alpha_i^2, \beta_i^2$ in terms of them. Expressions for $\alpha_1^2$ and $\beta_1^2$ are easily obtained from known formulas for the squared components of Ritz vectors from a two-dimensional subspace in a three-dimensional space. For the second component of the unit Ritz vector corresponding to $\tilde{\lambda}_1$, this formula is given and proved below as (2.2.24). For the remaining components, the required formula is easily obtained from (2.2.24) by an appropriate permutation of indices. One may also use the general statement of Lemma~2.2.2, which is proved in \S\,2.2.

% === Orig. p. 41 [s22R] ===
Thus:
\begin{align*}
\alpha_1^2 &= \frac{(\lambda_1 - \tilde{\lambda}_1)(\lambda_2 - \tilde{\lambda}_2)(\lambda_1 - \lambda_3)}{(\tilde{\lambda}_1 - \lambda_2)(\lambda_1 - \tilde{\lambda}_2)(\lambda_2 - \lambda_3)}, \\[1ex]
\beta_1^2 &= \frac{(\lambda_1 - \tilde{\lambda}_1)(\tilde{\lambda}_2 - \lambda_3)(\lambda_1 - \lambda_2)}{(\tilde{\lambda}_1 - \lambda_3)(\lambda_1 - \tilde{\lambda}_2)(\lambda_2 - \lambda_3)}.
\end{align*}
From a general result for the Lanczos method \citep[p.~259]{Parlett1983}, it follows that (up to normalization) $u^1 = (A - \tilde{\lambda}_2 E) u^0$. Hence,
\[
\alpha_0^2 = \left( \frac{\lambda_1 - \tilde{\lambda}_2}{\lambda_2 - \tilde{\lambda}_2} \right)^2 \alpha_1^2, \qquad
\beta_0^2 = \left( \frac{\lambda_1 - \tilde{\lambda}_2}{\tilde{\lambda}_2 - \lambda_3} \right)^2 \beta_1^2.
\]
For convenience, let us introduce the new notation
\begin{align*}
\mu_1 &\equiv (\tilde{\lambda}_1 - \lambda_3) / (\lambda_1 - \lambda_3) \in (\lambda, 1), \\
\mu_2 &\equiv (\tilde{\lambda}_2 - \lambda_3) / (\lambda_1 - \lambda_3) \in (0, \lambda),
\end{align*}
where, recall, $\lambda = (\lambda_2 - \lambda_3) / (\lambda_2 - \lambda_3) \in (0, 1)$. % [math note: typo in original: denominator has \lambda_2 - \lambda_3 instead of \lambda_1 - \lambda_3; cf. orig. p. 37]
Then in these terms we have:
\begin{gather*}
\alpha_1^2 = \frac{(1 - \mu_1)(\lambda - \mu_2)}{\lambda(\mu_1 - \lambda)(1 - \mu_2)}, \qquad
\beta_1^2 = \frac{(1 - \lambda)(1 - \mu_1)\mu_2}{\lambda \mu_1 (1 - \mu_2)}; \\[1ex]
\alpha_0^2 = \frac{(1 - \mu_1)(1 - \mu_2)}{\lambda(\mu_1 - \lambda)(\lambda - \mu_2)}, \qquad
\beta_0^2 = \frac{(1 - \lambda)(1 - \mu_1)(1 - \mu_2)}{\lambda \mu_1 \mu_2}.
\end{gather*}

Now we are ready to prove each of the inequalities (S.3.2)--(S.3.4) in turn.

Let us prove (S.3.2). Write:
\[
\frac{r_1^2}{r_0^2} = \frac{(\mu_1 - \mu_2)\mu_2(\lambda - \mu_2)}{(\mu_1 + \mu_2 - \lambda)(1 - \mu_2)^2} \equiv f(\mu_1, \mu_2).
\]
% === Orig. p. 42 [s23L] ===
By direct computation one can verify that
\[
\operatorname{sgn}(\partial f / \partial \mu_1) = \operatorname{sgn}(2\mu_2 - \lambda),
\]
i.e., when examining $f(\mu_1, \mu_2)$ for an extremum with respect to $\mu_1$, it suffices to consider two cases:

1. $0 < \mu_2 \leqslant \lambda/2$. The function $f(\mu_1, \mu_2)$ is nonincreasing in $\mu_1$; hence,
\[
\sup_{\mu_1 \in (\lambda, 1)} f(\mu_1, \mu_2) = f(\lambda, \mu_2) = \left( \frac{\lambda - \mu_2}{1 - \mu_2} \right)^2
\]
and, obviously,
\[
\sup_{\mu_2 \in (0, \lambda/2]} f(\lambda, \mu_2) = f(\lambda, 0) = \lambda^2. \tag{S.3.5}
\]

2. $\lambda/2 \leqslant \mu_2 < 0$. % [math note: typo in original: < 0 instead of < \lambda; cf. below and orig. p. 43]
The function $f(\mu_1, \mu_2)$ is nondecreasing in $\mu_1$, and
\[
\sup_{\mu_1 \in (\lambda, 1)} f(\mu_1, \mu_2) = f(1, \mu_2) = \frac{\mu_2(\lambda - \mu_2)}{(1 - \mu_2)(1 - \lambda + \mu_2)}.
\]
It is easy to verify that
\[
\begin{aligned}
\max_{\mu_2 \in [\lambda/2, \lambda)} f(1, \mu_2) &= \max_{\mu_2 \in (0, \lambda)} f(1, \mu_2) \\
&= f(1, \lambda/2) = \left( \frac{\lambda/2}{1 - \lambda/2} \right)^2 < \lambda^2.
\end{aligned}
\tag{S.3.6}
\]
Indeed, the inequality
\[
\frac{\mu_2(\lambda - \mu_2)}{(1 - \mu_2)(1 - \lambda + \mu_2)} \leqslant \left( \frac{\lambda/2}{1 - \lambda/2} \right)^2
\]
is equivalent to the inequality $(\mu_2 - \lambda/2)^2 \geqslant 0$. Note that assertion (S.3.6) will be used in part~2 to verify inequality (S.1.7), and (S.3.5) in part~3 to verify (S.1.6).

% === Orig. p. 43 [s23R] ===
Let us prove inequality (S.3.3):
\[
\frac{\tan^2 \varphi_1}{\tan^2 \varphi_0} = \frac{(\lambda - \mu_2)\mu_2}{(1 - \mu_2)^2} \cdot \frac{\mu_1(1 - \mu_2) - (1 - \lambda)\mu_2}{\mu_1(1 - \lambda + \mu_2) - (1 - \lambda)(\lambda - \mu_2)} \equiv g(\mu_1, \mu_2).
\]
By analogy with the preceding argument, one can show that
\[
\operatorname{sgn}(\partial g / \partial \mu_1) = \operatorname{sgn}(2\mu_2 - \lambda),
\]
and the same two cases as before should be considered.

1. $0 < \mu_2 \leqslant \lambda/2$. For such $\mu_2$ we have:
\[
\begin{aligned}
&\sup_{\mu_1 \in (\lambda, 1)} g(\mu_1, \mu_2) = g(\lambda, \mu_2) = \left( \frac{\lambda - \mu_2}{1 - \mu_2} \right)^2, \\
&\sup_{\mu_2 \in (0, \lambda/2]} g(\lambda_1, \mu_2) = g(\lambda, 0) = \lambda^2. % [math note: typo in original: g(\lambda_1, \mu_2) instead of g(\lambda, \mu_2)]
\end{aligned}
\tag{S.3.7}
\]

2. $\lambda/2 \leqslant \mu_2 < \lambda$. For such $\mu_2$ we have:
\[
\begin{aligned}
\sup_{\mu_1 \in (\lambda, 1)} g(\mu_1, \mu_2) &= g(1, \mu_2) \\
&= \frac{(\lambda - \mu_2)\mu_2}{(1 - \mu_2)^2} \cdot \frac{1 - (2 - \lambda)\mu_2}{1 - (2 - \lambda)(\lambda - \mu_2)}.
\end{aligned}
\]
Let us show that in this case
\[
\max_{\mu_2 \in [\lambda/2, \lambda)} g(1, \mu_2) = g(1, \lambda/2) = \left( \frac{\lambda/2}{1 - \lambda/2} \right)^2.
\]
Indeed, the inequality
\[
\frac{(\lambda - \mu_2)\mu_2}{(1 - \mu_2)^2} \cdot \frac{1 - (2 - \lambda)\mu_2}{1 - (2 - \lambda)(\lambda - \mu_2)} \leqslant \left( \frac{\lambda/2}{1 - \lambda/2} \right)^2
\]
% === Orig. p. 44 [s24L] ===
is equivalent to the inequality
\[
(\mu_2 - \lambda/2) \left[ (\mu_2 - \lambda/2)^2 - \frac{1 - \lambda}{2 - \lambda} (\mu_2 - \lambda/2) - \lambda^2 / 4 \right] \leqslant 0,
\]
which is verified directly. Theorem~S.1.1 is proved.

Let us prove inequality (S.3.4). Write:
\[
\frac{\Delta\lambda_1}{\Delta\lambda_0} = \frac{(\lambda - \mu_2)(\mu_1 + \mu_2 - 1)}{(1 - \mu_2)(\mu_1 + \mu_2 - \lambda)} \equiv h(\mu_1, \mu_2).
\]
Next, we have:
\[
\sup_{\mu_1 \in (\lambda, 1)} h(\mu_1, \mu_2) = h(1, \mu_2) = \frac{(\lambda - \mu_2)\mu_2}{(1 - \mu_2)(1 + \mu_2 - \lambda)}
\]
and
\[
\max_{\mu_2 \in (0, \lambda)} h(1, \mu_2) = h(1, \lambda/2) = \left( \frac{\lambda/2}{1 - \lambda/2} \right)^2. \tag{S.3.8}
\]
Thus Theorem~S.1.2 is proved.

\textit{Proof of inequality (S.1.7).} From assertion (S.3.6) it follows that
\[
\max_{\mu_2 \in (0, \lambda)} f(1, \mu_2) \leqslant \left( \frac{1 - \xi}{1 + \xi} \right)^2.
\]
Consequently, as $u^0 \to u_1$ (or $\mu_1 \to 1$),
\[
f(\mu_1, \mu_2) \leqslant \sigma_1,
\]
where $\sigma_1 \equiv (1 - \xi)^2 / (1 + \xi)^2 + o(1)$. Similarly to the case $n = 0$, the inequality $r_{n+1}^2 \leqslant \sigma_1 r_n^2$ can be established for all $n = 1, \dots, k - 1$, which leads to (S.1.7).

% === Orig. p. 45 [s24R] ===
Note here that alongside (S.1.7) there also holds an inequality improving estimate (S.1.5) as $u^0 \to u_1$:
\[
\tan^2 \varphi_k \leqslant \rho_k \tan^2 \varphi_0,
\]
where $\rho_k \equiv \rho_k(u^0) < (1 - \xi)^{2k}$.

Deriving an explicit formula for $\rho_k$ (namely, finding $\max g(1, \mu_2)$ over $\mu_2 \in (0, \lambda)$) involves cumbersome calculations, and the authors deemed it inexpedient.

\textit{Proofs of assertions (S.1.6) and (S.1.9).}

Let $\tilde{\lambda}_1^n > \tilde{\lambda}_2^n$ denote the Ritz values obtained by the Rayleigh--Ritz method in the subspace $\mathcal{K}_n = \myspan\{u^n, A u^n\}$, $n = 0, 1, \dots$, and introduce the quantities:
\[
\mu_1^n \equiv \frac{\tilde{\lambda}_1^n - \lambda_3}{\lambda_1 - \lambda_3}, \qquad \mu_2^n \equiv \frac{\tilde{\lambda}_2^n - \lambda_3}{\lambda_1 - \lambda_3}.
\]
We choose initial approximations from the subspace $H^3 = \myspan\{u_1, u_2, u_{\min}\}$. Then $\lambda_3 = \lambda_{\min}$, $1 - \xi = \lambda$, and $u^n \in H^3$.

To prove assertion (S.1.6), it suffices to check the case $n = 0$. As follows from (S.3.5) and (S.3.7), we need to give an example of a family of vectors $\{u^0\}$ such that $\mu_1^0 \to \lambda$ and $\mu_2^0 \to 0$. Take
\[
u^0 = \varepsilon u_1 + u_2 + u_{\min},
\]
where $\varepsilon \to 0$. Then from the formulas for $\tilde{\lambda}_1 + \tilde{\lambda}_2$ and $\tilde{\lambda}_1 \cdot \tilde{\lambda}_2$ in the proof of Lemma~S.2.2 % [math note: typo in original: Lemma S.3.2 on orig. p. 39 is meant]
it is easy to see that $\tilde{\lambda}_1^0 \to \lambda_2$ and $\tilde{\lambda}_2^0 \to \lambda_{\min}$, and hence $\mu_1^0 \to \lambda$, $\mu_2^0 \to 0$.

Let us prove (S.1.9). Choose a family of vectors $\{u^0\}$ such that $\mu_1^0 \to 1 - 0$, $\mu_2^0 \to \lambda/2$. It is then easy to see that $\varphi_0 \to 0$, and according to (S.3.8) we have:
\[
\Delta\lambda_1 / \Delta\lambda_0 \to (1 - \xi)^2 / (1 + \xi)^2.
\]
% === Orig. p. 46 [s25L] ===
Perform one step of the method (S.1.1)--(S.1.3) to obtain the family $\{u^1\}$. Let us show that
\[
\mu_1^1 \to 1 - 0, \qquad \mu_2^1 \to \lambda/2.
\]
Then we can assert that
\[
\Delta\lambda_2 / \Delta\lambda_1 \to (1 - \xi)^2 / (1 + \xi)^2,
\]
and by induction
\[
\Delta\lambda_{n+1} / \Delta\lambda_n \to (1 - \xi)^2 / (1 + \xi)^2
\]
for any $n = 2, \dots, k - 1$, which leads to (S.1.9).

Choose the vectors $u^1 \in \{u^1\}$ and $w^1$ from (S.1.2) as a basis of the subspace $\mathcal{K}_1$ and carry out the Rayleigh--Ritz procedure to determine $\tilde{\lambda}_1^1$ and $\tilde{\lambda}_2^1$. Then we have:
\begin{align*}
\tilde{\lambda}_1^1 &= \frac{1}{2}(\tilde{\lambda}_1^0 + \lambda(w^1) + |G|), \\
\tilde{\lambda}_1^2 &= \frac{1}{2}(\tilde{\lambda}_1^0 + \lambda(w^1) - |G|), % [math note: typo in original: \tilde{\lambda}_1^2 instead of \tilde{\lambda}_2^1]
\end{align*}
where
\[
G \equiv \left[ (\tilde{\lambda}_1^0 - \lambda(w^1))^2 + 4 \|w^1\|^2 \|u^1\|^{-2} \right]^{1/2}.
\]
Taking into account that $H^3$ is an invariant subspace of the operator $A$ and performing the necessary calculations, we obtain
\begin{align*}
\frac{\lambda(w^1) - \lambda_3}{\lambda_1 - \lambda_3} &= 1 - \mu_1^0 - \mu_2^0 + \lambda, \\[1ex]
\|w^1\|^2 \|u^1\|^{-2} &= \frac{(1 - \mu_1^0)(\mu_1^0 - \lambda)\mu_1^0}{\mu_1^0 - \mu_2^0}. % [math note: in original, the factor (\lambda_1 - \lambda_3)^{-2} is omitted on the left-hand side]
\end{align*}
Then for $\mu_1^1, \mu_2^1$ we obtain the equalities:
\begin{gather*}
\mu_1^1 = \frac{1}{2}(1 - \mu_2^0 + \lambda + |F|), \tag{S.3.9} \\
% === Orig. p. 47 [s25R] ===
\mu_2^1 = \frac{1}{2}(1 - \mu_2^0 + \lambda - |F|), \tag{S.3.10} \\[1ex]
F \equiv \left[ (2\mu_2^0 + \mu_1^0 - 1 - \lambda)^2 + 4 \frac{(1 - \mu_1^0)(\mu_1^0 - \lambda)\mu_1^0}{\mu_1^0 - \mu_2^0} \right]^{1/2}. \tag{S.3.11} % [math note: typo in original: 2\mu_2^0 + \mu_1^0 instead of 2\mu_1^0 + \mu_2^0]
\end{gather*}
Since $\mu_1^1 > \mu_1^0$ and $|F| \to 1 - \lambda/2$, we obviously have
\[
\mu_1^1 \to 1 - 0, \qquad \mu_2^1 \to \lambda/2.
\]

Assertions (S.1.6) and (S.1.9) are proved.

Let us give an example of a family $\{u^0\}$ of vectors of the form
\[
u^0 = u_1 + \alpha_0 u_2 + \beta_0 u_{\min}
\]
such that $u^0 \to u_2$, $\lambda(u^0) > \lambda_2$, and for any fixed $k$ we also have $u^k \to u_2$.

Suppose that $2\lambda_2 > \lambda_1 + \lambda_{\min}$. Then $\lambda > 1/2$. Choose $\{u^0\}$ so that $\mu_1^0 \to \lambda + 0$, $\mu_2^0 = 1/2$. Then
\begin{align*}
\alpha_0^2 &= \frac{(1 - \mu_1^0)(1 - \mu_2^0)}{\lambda(\mu_1^0 - \lambda)(\lambda - \mu_2^0)} \to \infty, \\[1ex]
\beta_0^2 &= \frac{(1 - \lambda)(1 - \mu_1^0)(1 - \mu_2^0)}{\lambda \mu_1^0 \mu_2^0} \to \left( \frac{1 - \lambda}{\lambda} \right)^2,
\end{align*}
and consequently $\ang(u^0; u_2) \to 0$.

Let us show that $\lambda(u^0) > \lambda_2$. Indeed, with our choice of $\mu_2^0$ this inequality is equivalent to the inequality $\mu_1^0 > 1/2$, which holds since $\mu_1^0 > \lambda > 1/2$. From the inequality $\lambda(u^0) > \lambda_2$ it follows that $\lambda(u^k) > \lambda_2$ as well.

% === Orig. p. 48 [s26L] ===
Let us prove that if $\mu_1^0 \to \lambda + 0$ and $\mu_2^0 \to 1/2$, then $\mu_1^k \to \lambda + 0$, $\mu_2^k \to 1/2$, whence one can conclude that $\ang(u^k, u_2) \to 0$.

It suffices to check the case $k = 1$. To this end, observe that for our $\mu_1^0$ and $\mu_2^0$ in formula (S.3.11) we have $|F| \to \lambda - 1/2$, and from (S.3.9) and (S.3.10) it is easy to see that $\mu_1^1 \to \lambda + 0$ and $\mu_2^1 \to 1/2$.

In conclusion, based on the preceding arguments of \S\,3, % [orig: \S 3, i.e., \S S.3]
we give a natural and exhaustive proof of the following fact (cf.~\citep{Zhuk1981}).

\textbf{Lemma S.3.3.} If the expansion of the initial vector $u^0$ in the basis of eigenvectors of the operator $A$ contains three or more nonzero terms, then the method of steepest ascent (S.1.1)--(S.1.3) does not stabilize, i.e., $w^n \neq 0$, $n = 0, 1, \ldots$. % [orig: double punctuation after verb in 1986 print]

\textit{Proof of Lemma S.3.3.} Let $d_i^n = (u^n, u_i)$, $n = 0, 1, \ldots$. Without loss of generality, we assume that $d_1^0 \cdot d_2^0 \cdot d_{\min}^0 \neq 0$ (unlike \citep{Zhuk1981}, here it is not assumed that all remaining coefficients are zero). Let us show that $d_1^1 \cdot d_2^1 \cdot d_{\min}^1 \neq 0$; then by induction we can conclude that $d_1^n \cdot d_2^n \cdot d_{\min}^n \neq 0$, $n = 0, 1, \ldots$, and hence the method does not stabilize. We use the fact that, up to normalization, $u^1 = (A - \tilde{\lambda}_2 E) u^0$ (see the proofs of Theorems~S.1.1 and~S.1.2). Based on this representation, it obviously suffices to prove that $\tilde{\lambda}_2 \notin \{\lambda_1, \lambda_2, \lambda_{\min}\}$. To this end, write:
\[
\lambda_1^3 = \lambda_1 > \lambda_2 \geqslant \lambda_2^3 \geqslant \tilde{\lambda}_2 \geqslant \lambda_3^3 \geqslant \lambda_{\min}, \qquad a_1 = d_1^0 \neq 0.
\]
Suppose that $\lambda_2 = \lambda_2^3 = \tilde{\lambda}_2$. Then $a_2 = d_2^0 \neq 0$, and from Lemma~S.3.2 we conclude that $a_3 = 0$. Consequently,
\[
u^0 \in \myspan\{u_1, u_2\}, \qquad d_{\min}^0 = 0,
\]
which contradicts our assumption. Similarly, the case $\tilde{\lambda}_2 = \lambda_3^3 = \lambda_{\min}$ also leads to a contradiction. Thus, Lemma~S.3.3 is completely proved.

% =============================================================================
% SOURCE MODULE: part_03_ch2_a.tex
% =============================================================================
% === Orig. p. 49 [s26R] ===
\monochapter{2}{The Rayleigh--Ritz Method in Subspace Iterations and Block Algorithms}

We consider problem (0.1) with $p > 1$. Various error estimates for the Rayleigh--Ritz method are discussed, and generalizations of several estimates from Chapter~1 to subspace iterations are presented.

In \S\,2.1, the concepts of the gap between subspaces and the principal angle between subspaces are introduced, and their basic properties are stated.

In \S\,2.2, error estimates for Ritz vectors and Ritz values are proved, and the asymptotic behavior of Ritz vectors is studied.

In \S\,2.3, the abstract subspace iteration method is defined, and estimates for its convergence rate are formulated and proved.

In \S\,2.4, subspace iterations with a constant shift and with Chebyshev parameters are considered. A natural notion of optimality for subspace iterations is defined, and it is proved that iterations with Chebyshev parameters are optimal.

In \S\S\,2.5 and 2.6, the theory of block analogues of variational methods---steepest ascent, Lanczos, and conjugate gradients---is discussed.

Section~2.7 is devoted to the error analysis of orthogonalization against approximately computed eigenvectors.

In the concluding \S\,2.8, we point out the possibility of using, for $M > 0$, the operator $-M^{-1}L$ and the metric of the space $H_M$ instead of $H_L$ in the results of \S\S\,2.2--2.5 (cf.\ \S\,1.6). The relationship between the angles between subspaces in the metrics of $H_M$ and $H_L$ is studied.

Concluding this overview of the chapter, let us note that Chapter~2 presents block analogues for all the methods of Chapter~1 except inverse iteration with shifts from \S\,1.7. This is because, in the theoretically justified block analogues of inverse iteration with shifts known to the author, cubic convergence is observed for only one Ritz vector. Therefore, in passing from the vector version to the block version, the increase in computational cost does not pay off.

% === Orig. p. 50 [s27L] ===
\monosection{2.1}{Gaps between subspaces}

To measure the proximity in $H_L$ of two subspaces $V$ and $W$ with $\dim V = \dim W$, Chapter~2 uses the concept of the gap between subspaces (e.g., \citep{GlazmanLyubich1969}):
\[
\theta_L(V; W) \equiv \| P_V - P_W \|_L ,
\]
where $P_V$ and $P_W$ are the $H_L$-orthogonal projectors onto $V$ and $W$, and $\|\cdot\|_L$ is the induced operator norm in $H_L$. It is known that

1) $\theta_L(V; W) \in [0, 1]$;

2) if $\theta_L(V; W) = 0$, then $V = W$; the converse also holds;

3) if $\theta_L(V; W) = 1$, then there exists a vector $v \in V \setminus 0$ such that $P_W v = 0$; the converse also holds;

4) $\displaystyle \theta_L(V; W) = \max_{v \in V \setminus 0} \| v - P_W v \|_L \| v \|_L^{-1}$;

5) $\displaystyle \theta_L^2(V; W) = 1 - \min_{v \in V \setminus 0} \| P_W v \|_L^2 \| v \|_L^{-2}$.

It is obvious from the definition that $\theta_L(V; W) = \theta_L(W; V)$, and alongside statements 3)--5) their symmetric counterparts also hold.

If the $p$-dimensional subspaces $V$ and $W$ are specified by $H_L$-orthonormal bases $v_1, \ldots, v_p$ and $w_1, \ldots, w_p$, then property 5) implies $\theta_L^2(V; W) = 1 - \sigma_{\min}$, where $\sigma_{\min}$ is the smallest eigenvalue of the symmetric positive semidefinite $p \cdot p$ matrix % [orig: positive (in the operator sense >= 0, i.e., positive semidefinite)]
with entries
\[
\sigma_{ij} = \sum_{k=1}^p (v_i; w_k)_L \cdot (v_j; w_k)_L .
\]
% === Orig. p. 51 [s27R] ===
In \citep{Parlett1983}, the concept of the principal angle between subspaces in $H_L$ is used:
\[
\ang_L(V; W) \equiv \arccos \sqrt{\sigma_{\min}} .
\]
The angle between subspaces and the gap between subspaces are related by the simple identity $\theta_L(V; W) = \sin \ang_L(V; W)$, which allows one to compare results formulated in these terms.

The concept of the gap between subspaces has been known for a relatively long time \citep{KreinKrasnoselskii1947} and is widely used in works on operator spectral theory and perturbation theory, e.g., \citep{Kato1970ru,Chatelin1983}. The concept of the principal angle between subspaces is secondary, yet more intuitive \citep{Parlett1983,Wedin1983}.

In the present work, angles between subspaces are used in formulating the results, while both concepts are employed in the proofs. Most frequently, one encounters the quantity
\[
\tan^2 \ang_L(V; W) = \max_{v \in V \setminus 0} \| v - P_W v \|_L^2 \| P_W v \|_L^{-2} .
\]

Gaps and angles are defined in a completely analogous manner for other inner products, for instance, in $H_M$ when $M > 0$. Later, in \S\,2.8, the dependence of gaps and angles on the choice of metric is studied.

It is obvious from the definition that the gap between the subspaces $V$ and $W$ coincides with the largest singular value of the operator $P_V - P_W$. One can also consider the remaining singular values of this operator, to which angles between the subspaces will likewise correspond. In particular, the smallest singular value corresponds to the smallest angle; if they are equal to zero, this means that the two given subspaces have at least a one-dimensional intersection. Details can be found in the book \citep{GolubVanLoan1983}, which also presents the algorithm from \citep{BjorckGolub1973} for computing the angles via the singular value decomposition.

% === Orig. p. 52 [s28L] ===
\monosection{2.2}{Error estimates and asymptotic properties of the Rayleigh--Ritz method}

For a subspace $\widetilde{U}$ ($\dim \widetilde{U} = p$), let us define the projection of problem (0.1) onto $\widetilde{U}$ by the Rayleigh--Ritz method \citep{Parlett1983}:
\[
\widetilde{P} M \big|_{\widetilde{U}} \tilde{u} = \tilde{\lambda} \widetilde{P} L \big|_{\widetilde{U}} \tilde{u} , \quad \tilde{u} \in \widetilde{U}, \tag{2.2.1}
\]
where $\widetilde{P}$ is the orthogonal projector in $H$ onto $\widetilde{U}$, and the restriction of the operators to $\widetilde{U}$ is taken. The eigenvalues $\tilde{\lambda}_1 \geqslant \ldots \geqslant \tilde{\lambda}_p$ of (2.2.1) and the corresponding $H_L$-orthogonal eigenvectors $\tilde{u}_1, \ldots, \tilde{u}_p$ will be called the Ritz values and Ritz vectors, and will be regarded as approximations to $\lambda_i$ and $u_i$ ($i = 1, \ldots, p$). Let us present several error estimates.

\textbf{Theorem 2.2.1.} For any $i \in [1, p]$, the following inequality holds:
\[
\begin{aligned}
&\sin^2 \ang_L(u_i; \tilde{u}_i) \leqslant (1 + \gamma^2 / d^2) \sin^2 \ang_L(u_i; \widetilde{U}), \\
&\gamma \equiv \| (E - \widetilde{P}_L) L^{-1} M \widetilde{P}_L \|_L , \\
&d \equiv \min_{j \neq i} | \tilde{\lambda}_j - \lambda_i | ,
\end{aligned}
\tag{2.2.2}
\]
where $\widetilde{P}_L$ is the orthogonal projector in $H_L$ onto $\widetilde{U}$.

Theorem~2.2.1 in a slightly modified formulation is due to Saad \citep{Saad1983}; in \citep{Parlett1983}, an elegant version of Saad's original proof \citep{Saad1980} is given.

Inequality (2.2.2) turns into an equality in a number of cases:
\[
\max_{\widetilde{U}} \frac{\sin^2 \ang_L(u_i; \tilde{u}_i)}{\sin^2 \ang_L(u_i; \widetilde{U})} = 1 + \gamma^2 / d^2 \tag{2.2.3}
\]
and is asymptotically sharp:
\[
\sup_{\widetilde{U} : \, \sin^2 \ang_L(u_i; \widetilde{U}) \approx 0} \frac{\sin^2 \ang_L(u_i; \tilde{u}_i)}{\sin^2 \ang_L(u_i; \widetilde{U})} = 1 + \gamma^2 / d^2 . \tag{2.2.4}
\]
Equalities (2.2.3) and (2.2.4) are probably valid in the general case as well (for arbitrary operators $L^{-1}M$).

% === Orig. p. 53 [s28R] ===
Estimate (2.2.2) is not maximally sharp, that is, one cannot impose in (2.2.3) a constraint of the form $\sin^2 \ang_L(u_i; \widetilde{U}) \equiv c \in [0, 1]$. Preliminary investigations for a three-dimensional space $H$ have shown that a correspondingly refined version promises to be quite cumbersome, so such an improvement of estimate (2.2.2) is of little interest.

A comparison was carried out between (2.2.2) and a special case of one of Vainikko's classical results \citep{Krasnoselskii1969}, in which the quantity
\[
\sin^2 \ang_L(u_i; \tilde{u}_i) / \sin^2 \ang_L(u_i; \widetilde{U})
\]
is in effect also bounded from above, the role of $\gamma$ on the right-hand side is played by the norm $\| (E - \widetilde{P}_L) L^{-1} M \|_L$, which is no smaller than $\gamma$, and $\tilde{\lambda}_i$ enters additionally. Estimate (2.2.2) turned out to be sharper, even though it uses less information (the necessary derivations are given in the concluding part of the present section).

For certain problems (0.1) with closely spaced eigenvalues, the problem of estimating the error of the corresponding invariant subspaces is of prime importance. % [orig: typo in Russian scan for "subspaces"]
Possibly, an estimate of the form\begingroup\renewcommand{\thefootnote}{*)}\footnote{\,For $\lambda_i = \ldots = \lambda_{i+l}$, the estimate was proved in \citep{Saad1980}.}\endgroup
\[
\begin{aligned}
&\sin^2 \ang_L(\myspan\{u_i, \ldots, u_{i+l}\}; \myspan\{\tilde{u}_i, \ldots, \tilde{u}_{i+l}\}) \\
&\qquad \leqslant (1 + \gamma^2 / d^2) \sin^2 \ang_L(\myspan\{u_i, \ldots, u_{i+l}\}; \widetilde{U}), \\
&d \equiv \min_{j \neq i, \ldots, i+l; \, k = i, \ldots, i+l} | \tilde{\lambda}_j - \lambda_k |
\end{aligned}
\tag{2.2.5}
\]
holds with $\gamma$ from (2.2.2); however, the known methods of proving (2.2.2) do not readily admit such a generalization. In (2.2.5) and hereafter, the distance from one subspace to another, say from $V$ to $W$ ($\dim V \leqslant \dim W$) in the space $H_L$, is measured by the quantity
\[
\sin \ang_L(V; W) \equiv \max_{v \in V, \, \|v\|_L = 1} \operatorname{dist}_L(v; W) \tag{2.2.6}
\]
% === Orig. p. 54 [s29L] ===
and various functions thereof; note that for $\dim V = \dim W$, (2.2.6) coincides with the definition of the gap between the subspaces $V$ and $W$ in $H_L$ (see \S\,2.1). % [orig: missing period after "(see \S 2.1)" in scan]

Concluding our discussion of Saad's estimate, let us note that it is easy to generalize to the case of a Hilbert space $H_L$ and a bounded self-adjoint operator $L^{-1}M$ whose relevant part of the spectrum is discrete. This generalization can be used to study the Rayleigh--Ritz method as a discretization procedure (cf.\ \citep{Krasnoselskii1969}).

An error estimate for an eigenvector in terms of the eigenvalue error is given in the following theorem.

\textbf{Theorem 2.2.2.} If for some $i \in [2, p]$ we have $\tilde{\lambda}_{i-1} \geqslant \lambda_i$ and $\tilde{\lambda}_i > \lambda_{i+1}$, then
\[
\sin^2 \ang_L(u_i; \tilde{u}_i) \leqslant 1 - \frac{(\lambda_1 - \tilde{\lambda}_i)(\tilde{\lambda}_i - \lambda_{i+1})(\tilde{\lambda}_{i-1} - \lambda_i)}{(\lambda_1 - \lambda_i)(\lambda_i - \lambda_{i+1})(\tilde{\lambda}_{i-1} - \tilde{\lambda}_i)} . \tag{2.2.7}
\]

Inequality (2.2.7) is maximally sharp: for any $c \in [0, 1]$, there exists a subspace $\widetilde{U}$ such that both sides of (2.2.7) are equal to $c$. % [orig: missing period after handwritten c]

Estimate (2.2.7) (directly in its most general form) was proved by Weinberger \citep{Weinberger1960,Weinberger1974}. At the end of the present section, inequality (2.2.7) is proved anew in a simpler way. Note that $\lambda_1$ in (2.2.7) can be eliminated, which leads to a slightly weaker, yet more compact estimate:
\[
\sin^2 \ang_L(u_i; \tilde{u}_i) \leqslant \frac{(\lambda_i - \tilde{\lambda}_i)(\tilde{\lambda}_{i-1} - \lambda_{i+1})}{(\lambda_i - \lambda_{i+1})(\tilde{\lambda}_{i-1} - \tilde{\lambda}_i)} . \tag{2.2.8}
\]
Inequality (2.2.8) is sharp in a weaker sense than (2.2.7)---over the class of operators $A$ with no restriction on the magnitude of $\lambda_1$.

For $i = 1$, when estimates (2.2.7) and (2.2.8) are inapplicable, there is the classical result
\[
\sin^2 \ang_L(u_1; \tilde{u}_1) \leqslant (\lambda_1 - \tilde{\lambda}_1) / (\lambda_1 - \lambda_2), \tag{2.2.9}
\]
which essentially coincides with the left-hand inequality in (1.1.3).

% === Orig. p. 55 [s29R] ===
The validity of generalizations of inequalities (2.2.9) and (2.2.8) to invariant subspaces has been proved in \citep{Weinberger1960}:
\[
\begin{aligned}
&\sin^2 \ang_L(\myspan\{u_1, \ldots, u_l\}; \myspan\{\tilde{u}_1, \ldots, \tilde{u}_l\}) \leqslant (\lambda_1 - \tilde{\lambda}_l) / (\lambda_1 - \lambda_{l+1}), \\
&\sin^2 \ang_L(\myspan\{u_i, \ldots, u_{i+l}\}; \myspan\{\tilde{u}_i, \ldots, \tilde{u}_{i+l}\}) \\
&\qquad \leqslant (\lambda_i - \tilde{\lambda}_{i+l})(\tilde{\lambda}_{i-1} - \lambda_{i+l+1})(\lambda_i - \lambda_{i+l+1})^{-1}(\tilde{\lambda}_{i-1} - \tilde{\lambda}_{i+l})^{-1}
\end{aligned}
\]
(cf.\ (2.2.5)); the latter inequality holds under the conditions
\[
2 \leqslant i \leqslant i + l \leqslant p, \quad \tilde{\lambda}_{i-1} \geqslant \lambda_i, \quad \tilde{\lambda}_{i+l} \geqslant \lambda_{i+l+1} .
\]

Now let us state a simple error estimate for the eigenvalues.

\textbf{Theorem 2.2.3.} For $i \in [1, p]$, the following two-sided inequality holds:
\[
0 \leqslant (\lambda_i - \tilde{\lambda}_i) / (\tilde{\lambda}_i - \lambda_{\min}) \leqslant \tan^2 \ang_L(\widetilde{U}; U) . \tag{2.2.10}
\]

Here and below, $U = \myspan\{u_1, \ldots, u_p\}$.

Note that (2.2.10) generalizes the right-hand inequality in (1.1.9) and is maximally sharp, that is, for any $c > 0$ there exists a subspace $\widetilde{U}$ such that both sides of the right-hand inequality in (2.2.10) are equal to $c$.

Estimates of the type (2.2.10) are also proved in \citep{StrangFix1973}.

The difficulties in deriving error estimates for the Rayleigh--Ritz method stem from the intricate behavior of the Ritz vectors $\tilde{u}_i$. At the same time, under the condition $\tan \ang_L(\widetilde{U}; U) < \infty$, the subspace $\widetilde{U}$ contains two sets of vectors with simpler properties: $u_i^*$, the $H_L$-orthogonal projections of the vectors $u_i$ onto $\widetilde{U}$, and $\bar{u}_i$, the vectors defined by the conditions
\[
(\bar{u}_i, u_j)_L = 0, \quad \bar{u}_i \in \widetilde{U} \setminus 0, \quad i \neq j; \quad i, j = 1, \ldots, p. \tag{2.2.11}
\]
% === Orig. p. 56 [s30L] ===
A proof of the existence and uniqueness up to normalization of the vectors $\bar{u}_i$ is given in \citep{FaddeevFaddeeva1963}, where the vectors $\bar{u}_i$ are applied to analyze the convergence of simultaneous iterations; they are used for the same purpose in \citep{Parlett1983} and in \S\,2.3 of the present work.

The following theorem shows that the generally false assertion that the vectors $\tilde{u}_i$, $u_i^*$, and $\bar{u}_i$ coincide turns out to be justified to some extent asymptotically (as $\widetilde{U} \to U$).

\textbf{Theorem 2.2.4.} Let
\[
\varepsilon_i \equiv \ang(u_i; \widetilde{U}), \quad i = 1, \ldots, p, \quad \varepsilon \equiv \max \varepsilon_i .
\]
Suppose that $\varepsilon \to 0$ as a function of the subspace $\widetilde{U}$ and that $\varepsilon^3 = O(\varepsilon_i)$. Then
\begin{gather*}
\ang(\tilde{u}_i; u_i) \asymp \ang(u_i^*; u_i) \asymp \ang(\bar{u}_i; u_i) \asymp \varepsilon_i , \tag{2.2.12} \\
\begin{aligned}
\ang(\tilde{u}_i; \bar{u}_i) &= O(\varepsilon \varepsilon_i) , \\
\ang(\bar{u}_i; u_i^*) &= O(\varepsilon \varepsilon_i) , \\
\ang(u_i^*; \tilde{u}_i) &= O(\varepsilon \varepsilon_i) .
\end{aligned}
\tag{2.2.13}
\end{gather*}

If greater freedom in choosing transformations of the problem is allowed when defining the Ritz vectors, then it turns out that the asymptotic properties of such vectors are preserved.

\textbf{Theorem 2.2.5.} Let $F$ and $G$ be $H_L$-self-adjoint operators commuting with $L^{-1}M$ such that $\alpha F + \beta G > 0$ in $H_L$ for some numbers $\alpha, \beta$ and the eigenvectors $u_1, \ldots, u_p$ correspond to simple eigenvalues in the problem $Fu = \eta Gu$. Suppose also that the hypotheses of Theorem~2.2.4 hold. Define the vectors $\tilde{u}_i^\eta$, $i = 1, \ldots, p$, as the eigenvectors of the problem
\[
\widetilde{P}_L F \big|_{\widetilde{U}} \tilde{u}^\eta = \tilde{\eta} \widetilde{P}_L G \big|_{\widetilde{U}} \tilde{u}^\eta . \tag{2.2.14}
\]
% === Orig. p. 57 [s30R] ===
Then, under a suitable numbering of the vectors $\tilde{u}_i^\eta$, the following relations hold:
\[
\ang(\tilde{u}_i^\eta; u_i) \asymp \varepsilon_i , \quad \ang(\tilde{u}_i^\eta; \tilde{u}_i) = O(\varepsilon \varepsilon_i), \quad i = 1, \ldots, p. \tag{2.2.15}
\]

Let us consider two important examples with $L = E$.

Set $F = A^{-2}$ and $G = E$. Then the hypotheses of Theorem~2.2.5 are satisfied and relations (2.2.15) hold. One may say that the vectors $\tilde{u}_i$ and $\tilde{u}_i^\eta$ coincide asymptotically, although the exact assertion $\ang(\tilde{u}_i^\eta; \tilde{u}_i) = 0$ \citep{Pissanetzky1984} is, generally speaking, false.

In the second example, set $F = A(A - \alpha E)$ and $G = A - \alpha E$, in which case formula (2.2.14) describes the Temple--Lehmann method. We have $F - \alpha G > 0$ provided $\alpha$ is not an eigenvalue of problem (0.1), and the hypotheses of Theorem~2.2.5 are satisfied.

Let us proceed to the proofs of the statements formulated above.

\textit{Proof of the attainability relation (2.2.3).} For the case $\dim H = 3$, given any $\gamma$ and $d > 0$, we construct an example of an operator $A$ and a subspace $\widetilde{U}$ such that the nonstrict inequality (2.2.2) becomes an equality (assuming $L = E$ and $A = M$).

Let $\tilde{u}_1 = (1, 0, 0)^T$, $\tilde{u}_2 = (0, 1, 0)^T$, $\widetilde{U} = \myspan(\tilde{u}_1, \tilde{u}_2)$,
\[
A = \begin{pmatrix} \tilde{\lambda}_1 & 0 & a \\ 0 & \tilde{\lambda}_2 & 0 \\ a & 0 & c \end{pmatrix} , \quad \tilde{\lambda}_1 > 0 > \tilde{\lambda}_2 , \quad c = a^2 / \tilde{\lambda}_1 .
\]
Then the eigenvalue $\lambda_2 = 0$ has the corresponding eigenvector $u_2 = (+a, 0, -\tilde{\lambda}_1)^T$. We have
\[
\begin{aligned}
\sin^2 \ang(u_2, \tilde{u}_2) &= 1 , \\
\sin^2 \ang(u_2, \widetilde{U}) &= \frac{\tilde{\lambda}_1^2}{\tilde{\lambda}_1^2 + a^2} .
\end{aligned}
\]
On the other hand, $d = \tilde{\lambda}_1 - \lambda_2 = \tilde{\lambda}_1$ and $\gamma^2 = a^2$. Consequently, in this case we indeed have the equality
\[
\frac{\sin^2 \ang(u_2, \tilde{u}_2)}{\sin^2 \ang(u_2, \widetilde{U})} = 1 + \frac{\gamma^2}{d^2} .
\]
% === Orig. p. 58 [s31L] ===
The condition $\dim H = 3$ here apparently involves no loss of generality.

\textit{Proof of relation (2.2.4).} As in the proof of relation (2.2.3), we carry out the argument only for the case $\dim H = 3$ and assume $L = E$, $A = M$. For given numbers $\gamma$ and $d > 0$, we construct an example of operators $A$ and a subspace $\widetilde{U}$ such that the left-hand and right-hand sides in the nonstrict inequality (2.2.2) tend to zero and the inequality turns into an asymptotic equality.

Let $\tilde{u}_1 = (1, 0, 0)^T$, $\tilde{u}_2 = (0, 1, 0)^T$, $\widetilde{U} = \myspan(\tilde{u}_1, \tilde{u}_2)$,
\[
A = \begin{pmatrix} \tilde{\lambda}_1 & 0 & a \\ 0 & \tilde{\lambda}_2 & b \\ a & b & c \end{pmatrix} , \quad \tilde{\lambda}_1 > 0 > \tilde{\lambda}_2 , \quad c = a^2 / \tilde{\lambda}_1 + b^2 / \tilde{\lambda}_2 .
\]
Then the eigenvalue $\lambda_2 = 0$ has the corresponding eigenvector
\[
u_2 = (\tilde{\lambda}_2 a, \tilde{\lambda}_1 b, -\tilde{\lambda}_1 \tilde{\lambda}_2)^T .
\]
We have
\[
\begin{aligned}
\sin^2 \ang(u_2, \tilde{u}_2) &= \frac{\tilde{\lambda}_2^2 a^2 + \tilde{\lambda}_1^2 \tilde{\lambda}_2^2}{\tilde{\lambda}_2^2 a^2 + \tilde{\lambda}_1^2 \tilde{\lambda}_2^2 + \tilde{\lambda}_1^2 b^2} , \\
\sin^2 \ang(u_2, \widetilde{U}) &= \frac{\tilde{\lambda}_1^2 \tilde{\lambda}_2^2}{\tilde{\lambda}_2^2 a^2 + \tilde{\lambda}_1^2 \tilde{\lambda}_2^2 + \tilde{\lambda}_1^2 b^2} ,
\end{aligned}
\]
consequently,
\[
\frac{\sin^2 \ang(u_2, \tilde{u}_2)}{\sin^2 \ang(u_2, \widetilde{U})} = 1 + \frac{a^2}{\tilde{\lambda}_1^2} .
\]
On the other hand, $\gamma^2 = a^2 + b^2$ and, obviously, $d^2 = \lambda_1^2$. % [math note: author slip in scan; meant $d^2 = \tilde{\lambda}_1^2$]

Thus, taking $b \to 0$ and $\tilde{\lambda}_2 = o(b)$, we obtain
\begin{gather*}
\sin^2 \ang(u_2, \widetilde{U}) \sim \tilde{\lambda}_2^2 / b^2 \to 0 , \\
% === Orig. p. 59 [s31R] ===
\frac{\sin^2 \ang(u_2, \tilde{u}_2)}{\sin \ang(u_2, \widetilde{U})} = 1 + \frac{\gamma^2}{d^2} - \frac{b^2}{d^2} \longrightarrow 1 + \frac{\gamma^2}{d^2} . % [math note: author slip in scan; meant $\sin^2 \ang(u_2, \widetilde{U})$ in denominator, and period omitted at end of display]
\end{gather*}

\textit{Comparison of estimate (2.2.2) with a special case of one of Vainikko's results from \citep{Krasnoselskii1969}.} To simplify the notation, assume $L = E$ and $A = M$. Let us write out the estimate from \citep{Krasnoselskii1969} for the error of an eigenvector $u_i$ corresponding to a simple eigenvalue $\lambda_i$ ($1 < i < p$) of an operator $A > 0$:
\begin{gather*}
\| u_i - \widetilde{P}_i u_i \| \leqslant \tilde{b}_i' \| (E - \widetilde{P}) A u_i \| , \tag{2.2.16} \\
\tilde{b}_i' \equiv \frac{\lambda_i^{-1} \tilde{\lambda}_i^{-1} (1 + \tilde{\gamma}_i \lambda_i^{-1} \| (E - \widetilde{P}) A \|)}{\lambda_i^{-1} - \tilde{\gamma}_i (\tilde{\lambda}_i^{-1} - \lambda_i^{-1})} , \tag{2.2.17} \\
\tilde{\gamma}_i \equiv \max \Bigl\{ 1, \, \max_{j \neq i} \frac{\tilde{\lambda}_i^{-1}}{\tilde{\lambda}_j^{-1} - \tilde{\lambda}_i^{-1}} \Bigr\} , \tag{2.2.18} % [math note: scan has $\tilde{\lambda}_i^{-1} / (\tilde{\lambda}_j^{-1} - \tilde{\lambda}_i^{-1})$; cf. (2.2.21) which uses $\tilde{\lambda}_j^{-1} / |\tilde{\lambda}_j^{-1} - \tilde{\lambda}_i^{-1}| = \tilde{\lambda}_i |\tilde{\lambda}_j - \tilde{\lambda}_i|^{-1}$]
\end{gather*}
where it is assumed that $\tilde{b}_i' \geqslant 0$, and $\widetilde{P}_i$ denotes the orthogonal projector onto $\myspan\{\tilde{u}_i\}$. In its original version, estimate (2.2.16)--(2.2.18) is written in terms of the characteristic numbers of the operators $A$ and $\widetilde{P}A$, which is why the reciprocals of $\lambda$ and $\tilde{\lambda}$ appear in (2.2.17) and (2.2.18). Let us rewrite (2.2.16) in the form
\[
\frac{\| u_i - \widetilde{P}_i u_i \|}{\| u_i \|} \leqslant \lambda_i \tilde{b}_i' \frac{\| u_i - \widetilde{P} u_i \|}{\| u_i \|} ,
\]
which is easily expressed in terms of the corresponding angles:
\[
\sin \ang(u_i, \tilde{u}_i) \leqslant c_i \sin \ang(u_i, \widetilde{U}), \quad c_i \equiv \lambda_i \tilde{b}_i' . \tag{2.2.19}
\]
% === Orig. p. 60 [s32L] ===
Next we have (setting $\beta \equiv \| (E - \widetilde{P}) A \|$):
\begin{gather*}
c_i \equiv \lambda_i \tilde{b}_i' = \frac{\lambda_i + \tilde{\gamma}_i \beta}{\tilde{\lambda}_i - \tilde{\gamma}_i (\lambda_i - \tilde{\lambda}_i)} , \tag{2.2.20} \\
\begin{aligned}
\tilde{\gamma}_i &= \max \bigl\{ 1, \, \tilde{\lambda}_i \max_{j \neq i} | \tilde{\lambda}_j - \tilde{\lambda}_i |^{-1} \bigr\} \\
&= \tilde{\lambda}_i / \min_{j = i \pm 1} | \tilde{\lambda}_j - \tilde{\lambda}_i | .
\end{aligned}
\tag{2.2.21}
\end{gather*}

Substituting the value of $\tilde{\gamma}_i$ from (2.2.21) into (2.2.20), we obtain
\[
\begin{aligned}
c_i &= \frac{\lambda_i \min\limits_{j = i \pm 1} | \tilde{\lambda}_j / \tilde{\lambda}_i - 1 | + \beta}{\min\limits_{j = i \pm 1} | \tilde{\lambda}_j - \tilde{\lambda}_i | - (\lambda_i - \tilde{\lambda}_i)} \\
&= \max \bigl\{ [ \lambda_i (\tilde{\lambda}_{i-1} / \tilde{\lambda}_i - 1) + \beta ] / [ \tilde{\lambda}_{i-1} - \lambda_i ] ; \\
&\phantom{= \max \bigl\{ } [ \lambda_i (1 - \tilde{\lambda}_{i+1} / \tilde{\lambda}_i) + \beta ] / [ 2\tilde{\lambda}_i - \tilde{\lambda}_{i+1} - \lambda_i ] \bigr\} .
\end{aligned}
\]
Taking into account that $\tilde{\lambda}_{i+1} \leqslant \tilde{\lambda}_i \leqslant \lambda_i$ and that both denominators are positive (by assumption), it is easy to verify that each of the fractions in the last expression for $c_i$ is decreasing in the parameter $\tilde{\lambda}_i \leqslant \lambda_i$. Then, under the additional and quite realistic assumption $\tilde{\lambda}_{i+1} < \lambda_i$, we deduce the inequality
\[
c_i \geqslant 1 + \beta / \min(\tilde{\lambda}_{i-1} - \lambda_i; \, \lambda_i - \tilde{\lambda}_{i+1}). \tag{2.2.22}
\]
By virtue of the assumptions made above,
\[
\begin{aligned}
0 &< \min(\tilde{\lambda}_{i-1} - \lambda_i, \, \lambda_i - \tilde{\lambda}_{i+1}) \\
&= \min_{j \neq i} | \tilde{\lambda}_j - \lambda_i | = d .
\end{aligned}
\]
% === Orig. p. 61 [s32R] ===
For the quantity $\beta$ in (2.2.22), the following bound is obvious:
\[
\beta \equiv \| (E - \widetilde{P}) A \| \geqslant \| (E - \widetilde{P}) A \widetilde{P} \| \equiv \gamma .
\]
Summing up what has been said, we conclude that
\[
c_i \geqslant 1 + \frac{\gamma}{d} \qquad \text{and} \qquad c_i \geqslant 1 + \frac{\gamma^2}{d^2} . % [math note: author slip in scan; meant $c_i^2 \geqslant 1 + \gamma^2/d^2$ since (2.2.19) involves $c_i$ unsquared]
\]
Comparing now estimates (2.2.2) and (2.2.19), it is easy to see that (2.2.2) is sharper than (2.2.19), and hence sharper than (2.2.16)--(2.2.18).

\textit{Proof of Theorem 2.2.2.} Following the technique of reduction to a ``low-dimensional'' case, we begin the proof of Theorem~2.2.2 by verifying the simplest special case in which estimate (2.2.7) remains nontrivial. Without loss of generality, we may assume $L = E$ and $A = M$ for notational convenience.

\textbf{Lemma 2.2.1.} Let $\dim H^3 = 3$, let $\lambda_1^3 > \lambda_2^3 > \lambda_3^3$ be the eigenvalues of an operator $A^3 = (A^3)^*$ corresponding to the eigenvectors $u_1^3, u_2^3, u_3^3$, let $\dim \widetilde{U}^3 = 2$, and let $\tilde{\lambda}_1^2 > \tilde{\lambda}_2^2$ be the Ritz values corresponding to the Ritz vectors $\tilde{u}_1^2, \tilde{u}_2^2$ of the operator $A^3$ in the subspace $\widetilde{U}^3$ (the superscript $3$ is used here to distinguish this three-dimensional case from the general one). % [math note: author slip in scan; superscript 2 written instead of 3 on $\tilde{\lambda}_1, \tilde{\lambda}_2, \tilde{u}_1, \tilde{u}_2$]
Then
\[
\sin^2 \ang(u_2^3, \tilde{u}_2^3) \leqslant 1 - \frac{(\lambda_1^3 - \tilde{\lambda}_2^3)(\tilde{\lambda}_2^3 - \lambda_3^3)(\tilde{\lambda}_1^3 - \lambda_2^3)}{(\lambda_1^3 - \lambda_2^3)(\lambda_2^3 - \lambda_3^3)(\tilde{\lambda}_1^3 - \tilde{\lambda}_2^3)} . \tag{2.2.23}
\]
Moreover, equality actually holds in (2.2.23), i.e., the nonstrict inequality (2.2.7) turns into an equality in this simplest special case.

\textit{Proof of Lemma 2.2.1.} For brevity, we omit the superscript $3$ within this proof. Let $\tilde{u}_2$ be a unit vector, and denote by $b_1, b_2, b_3$ the coefficients of its expansion in the orthonormal basis $u_1, u_2, u_3$; similarly, denote by $(a_1, a_2, a_3)$ the expansion coefficients of the unit vector $\tilde{u}_1$. Obviously,
% === Orig. p. 62 [s33L] ===
\[
\sin^2 \ang(u_2, \tilde{u}_2) = 1 - b_2^2 ,
\]
so to establish the validity of (2.2.23) it suffices to prove that
\[
b_2^2 = \frac{(\lambda_1 - \tilde{\lambda}_2)(\tilde{\lambda}_2 - \lambda_3)(\tilde{\lambda}_1 - \lambda_2)}{(\lambda_1 - \lambda_2)(\lambda_2 - \lambda_3)(\tilde{\lambda}_1 - \tilde{\lambda}_2)} . \tag{2.2.24}
\]
Two ways of proving (2.2.24) can be proposed---direct verification and the use of a more general fact for an inverse eigenvalue problem of arbitrary dimension. We first demonstrate the former approach as the simpler one, and then state as Lemma~2.2.2 and prove the general fact.

Let us write down the system of six equations:
\[
\begin{aligned}
&a_1 b_1 + a_2 b_2 + a_3 b_3 = 0 , \\
&\lambda_1 a_1 b_1 + \lambda_2 a_2 b_2 + \lambda_3 a_3 b_3 = 0 , \\
&\lambda_1 a_1^2 + \lambda_2 a_2^2 + \lambda_3 a_3^2 = \tilde{\lambda}_1 , \\
&\lambda_1 b_1^2 + \lambda_2 b_2^2 + \lambda_3 b_3^2 = \tilde{\lambda}_2 , \\
&a_1^2 + a_2^2 + a_3^2 = 1 , \\
&b_1^2 + b_2^2 + b_3^2 = 1 .
\end{aligned}
\]
These equations express in coordinates with respect to the eigenvector basis $u_1, u_2, u_3$ the following assertions: the first pair---that the vectors $\tilde{u}_1$ and $\tilde{u}_2$ are Ritz vectors in the subspace $\widetilde{U} = \myspan\{\tilde{u}_1, \tilde{u}_2\}$; the second pair---that $\tilde{\lambda}_j = \lambda(\tilde{u}_j)$, $j = 1, 2$; and the third pair---that $\tilde{u}_1$ and $\tilde{u}_2$ are unit vectors.

% =============================================================================
% SOURCE MODULE: part_04_ch2_b.tex
% =============================================================================
% === Orig. p. 63 [s33R] ===
Let us transform the system into the form
\[
\begin{aligned}
(\lambda_1 - \lambda_3) a_1 b_1 &= -(\lambda_2 - \lambda_3) a_2 b_2, \\
(\lambda_1 - \lambda_3) a_3 b_3 &= -(\lambda_1 - \lambda_2) a_2 b_2, \\
(\lambda_1 - \lambda_3) a_1^2 &= -(\lambda_2 - \lambda_3) a_2^2 + \tilde{\lambda}_1 - \lambda_3, \\
(\lambda_1 - \lambda_3) a_3^2 &= -(\lambda_1 - \lambda_2) a_2^2 + \lambda_1 - \tilde{\lambda}_1, \\
(\lambda_1 - \lambda_3) b_1^2 &= -(\lambda_2 - \lambda_3) b_2^2 + \tilde{\lambda}_2 - \lambda_3, \\
(\lambda_1 - \lambda_3) b_3^2 &= -(\lambda_1 - \lambda_2) b_2^2 + \lambda_1 - \tilde{\lambda}_2.
\end{aligned}
\]

Squaring the first two equations and substituting the expressions for $a_1^2$, $a_3^2$, $b_1^2$, and $b_3^2$ from the remaining equations, we thereby obtain two equations for $a_2^2$ and $b_2^2$:
\begin{gather*}
\bigl\{\tilde{\lambda}_1 - \lambda_3 - (\lambda_2 - \lambda_3) a_2^2\bigr\}\bigl\{\tilde{\lambda}_2 - \lambda_3 - (\lambda_2 - \lambda_3) b_2^2\bigr\} = (\lambda_2 - \lambda_3)^2 a_2^2 b_2^2, \\
\bigl\{\lambda_1 - \tilde{\lambda}_1 - (\lambda_1 - \lambda_2) a_2^2\bigr\}\bigl\{\lambda_1 - \tilde{\lambda}_2 - (\lambda_1 - \lambda_2) b_2^2\bigr\} = (\lambda_1 - \lambda_2)^2 a_2^2 b_2^2.
\end{gather*}

Expanding the braces and canceling like terms, we arrive at a system of two linear equations in the unknowns $a_2^2$ and $b_2^2$:
\begin{gather*}
(\tilde{\lambda}_2 - \lambda_3) a_2^2 + (\tilde{\lambda}_1 - \lambda_3) b_2^2 = (\tilde{\lambda}_2 - \lambda_3)(\tilde{\lambda}_1 - \lambda_3) / (\lambda_2 - \lambda_3), \\
(\lambda_1 - \tilde{\lambda}_2) a_2^2 + (\lambda_1 - \tilde{\lambda}_1) b_2^2 = (\lambda_1 - \tilde{\lambda}_2)(\lambda_1 - \tilde{\lambda}_1) / (\lambda_1 - \lambda_2).
\end{gather*}

Solving the latter system for $b_2^2$, we obtain (2.2.24). Since the above transformations of the original system can only introduce additional solutions, while the latter system has a unique solution, $b_2^2$ is uniquely expressed in terms of $\lambda_1 \geqslant \tilde{\lambda}_1 > \lambda_2 \geqslant \lambda_2 > \lambda_3$ % [math note: tilde omitted over the second \lambda_2 in original; should be \tilde{\lambda}_2]
by formula (2.2.24). The existence of a solution to the original system
% === Orig. p. 64 [s34L] ===
for $\lambda_1 \geqslant \tilde{\lambda}_1 \geqslant \lambda_2 \geqslant \lambda_2 \geqslant \lambda_3$ is known; % [math note: tilde omitted over the second \lambda_2 in original; should be \tilde{\lambda}_2]
see Exercise~10.1.2 in~\citep{Parlett1983}. The direct verification of statement (2.2.24) is complete.

Formula (2.2.24) can also be obtained as a consequence of a more general result.

\textbf{Lemma 2.2.2.} Let $\dim H = p + 1$, let $\lambda_1 > \ldots > \lambda_{p+1}$ be the eigenvalues of the operator $A$ corresponding to the orthonormal eigenvectors $u_1, \ldots, u_{p+1}$, and let $\dim \widetilde{U} = p$ and $\tilde{\lambda}_1 > \ldots > \tilde{\lambda}_p$ be the Ritz values corresponding to the Ritz vectors $\tilde{u}_1, \ldots, \tilde{u}_p$ (not normalized) of the operator $A$ in the subspace $\widetilde{U}$, with
\[
\lambda_1 > \tilde{\lambda}_1 > \lambda_2 > \tilde{\lambda}_2 > \ldots > \lambda_p > \tilde{\lambda}_p > \lambda_{p+1}. \tag{2.2.25}
\]
Then
\[
(\tilde{u}_k, u_i)^2 = \frac{1}{(\lambda_i - \tilde{\lambda}_k)^2} \, \frac{\prod_{j=1}^p (\lambda_i - \tilde{\lambda}_j)}{\prod_{\substack{j=1 \\ j \neq i}}^{p+1} (\lambda_i - \lambda_j)}. \tag{2.2.26}
\]

\textit{Proof of Lemma 2.2.2.} Let $w$ denote a unit vector in the orthogonal complement to $\widetilde{U}$ in $H$. Then
\[
(w, u_i)^2 = \frac{\prod_{j=1}^p (\lambda_i - \tilde{\lambda}_j)}{\prod_{\substack{j=1 \\ j \neq i}}^{p+1} (\lambda_i - \lambda_j)} > 0. \tag{2.2.27}
\]
This relation is proved in~\citep[pp.~432--433]{GolubVanLoan1983}. Implicitly, it is also contained in~\citep[formula~(7.9.4c)]{Parlett1983}, where the matrix $A$ is assumed to be reduced to tridiagonal form by the Lanczos method, the vector $w$ is the last coordinate vector and simultaneously the last vector of the Lanczos basis, and $\tilde{\lambda}_1, \ldots, \tilde{\lambda}_p$ is the spectrum of the leading principal submatrix of order $p$ of the tridiagonal matrix.

The key assertion is that under the interlacing conditions (2.2.25), the subspace $\widetilde{U}$ can be interpreted as a Krylov subspace of the operator $A$ with some, as yet unknown, initial
% === Orig. p. 65 [s34R] ===
approximation $u^0$. Next, one should make use of the remarkable representations~\citep[formula~(12.3.7)]{Parlett1983}:
\begin{gather*}
\tilde{u}_k = \prod_{\substack{j=1 \\ j \neq k}}^p (A - \tilde{\lambda}_j E) u^0, \\
w = \prod_{j=1}^p (A - \tilde{\lambda}_j E) u^0,
\end{gather*}
from which it follows that $\tilde{u}_k = (A - \tilde{\lambda}_k E)^{-1} w$. The latter equality together with (2.2.27) leads to (2.2.26).

Let us also indicate the form of the initial approximation:
\[
(u^0, u_i)^{-2} = \prod_{j=1}^p (\lambda_i - \tilde{\lambda}_j) \prod_{\substack{j=1 \\ j \neq i}}^{p+1} (\lambda_i - \lambda_j).
\]

We conclude the proof of Lemma~2.2.2 with a reference to~\citep{Weinberger1960}, where formulas of the type (2.2.26) are also used.

To obtain relation (2.2.24) from (2.2.26), one should set $p = 2$ in Lemma~2.2.2 and normalize the vector $\tilde{u}_2$ (in Lemma~2.2.2 the vectors $\tilde{u}_k$ are not normalized).

Let us consider the question of how essential the strict inequalities (2.2.25) are. The following lemma holds.

\textbf{Lemma 2.2.3.} Suppose that under the hypotheses of Lemma~2.2.2 at least one equality holds among the strict inequalities (2.2.25), for instance, $\lambda_i = \tilde{\lambda}_i$. Then $u_i = \tilde{u}_i$. If, in addition, it is known that the subspace $\widetilde{U}$ is a Krylov subspace, then all the remaining Ritz vectors also coincide with certain eigenvectors.

\textit{Proof of Lemma 2.2.3.} We use the well-known characterization of Ritz vectors via the properties of residuals~\citep[Exercise~11.4.6]{Parlett1983}:
\[
w \equiv A\tilde{u}_i - \tilde{\lambda}_i \tilde{u}_i \perp \widetilde{U}.
\]
On the other hand, $\tilde{\lambda}_i = \lambda_i$ implies $(w, u_i) = 0$. Since $\dim H = \dim \widetilde{U} + 1$ under the hypotheses of Lemma~2.2.2, we then have $u_i \in \widetilde{U}$ and, obviously, $u_i = \tilde{u}_i$.

Suppose, in addition, it is known that
\[
\widetilde{U} = \myspan\{u^0, Au^0, \ldots, A^{p-1} u^0\}.
\]
% === Orig. p. 66 [s35L] ===
Since $u_i = \widetilde{U}$, % [math note: u_i = \widetilde{U} in original instead of u_i \in \widetilde{U}]
there exists a polynomial $Q_{p-2}(\cdot)$ of degree at most $p - 2$ such that
\[
A^{p-1} u^0 \in \myspan\{u_i, Q_{p-2}(A) u^0\}.
\]
Multiplying both sides on the left by the operator $A$, we obtain
\[
A^p u^0 \in \myspan\{u_i, A Q_{p-2}(A) u^0\} \subseteq \widetilde{U}.
\]
Thus, the subspace $\widetilde{U}$ is invariant under the operator $A$: $A\widetilde{U} \subseteq \widetilde{U}$. Then the Ritz vectors in $\widetilde{U}$ coincide with the eigenvectors of the restriction $A|_{\widetilde{U}}$ of the operator $A$ to the subspace $\widetilde{U}$. The proof of Lemma~2.2.3 is complete.

Lemma~2.2.3 guarantees that the subspace $\widetilde{U}$ can always be represented as an orthogonal sum $\widetilde{U} = U \oplus \widetilde{V}$, where $U$ is an invariant subspace and for the subspace $\widetilde{V}$ the strict inequalities (2.2.25) necessarily hold. This means that conditions (2.2.25) are natural.

Thus, Lemma~2.2.2 is completely proved, and we can return to verifying the statement of Theorem~2.2.2 proper. To bring out the essence of the technique of reduction to the ``low-dimensional'' case, we formulate the final stage of the proof as a lemma as well.

\textbf{Lemma 2.2.4.} For Theorem~2.2.2 to hold, it is necessary and sufficient that Lemma~2.2.1 hold, i.e., Theorem~2.2.2 and Lemma~2.2.1 are equivalent.

\textit{Proof of Lemma 2.2.4.} Necessity is obvious. To verify sufficiency, we embed the three-dimensional space $H^3$ into $H$ in a special way as a subspace $H^3 \subset H$ (assuming $\dim H > 3$). We introduce the operator $P^3$ of orthogonal projection onto $H^3$ in $H$ and set $A^3 = P^3 A$. The constructions of $H^3$ and of the two-dimensional subspace $\widetilde{U}^3 \subset H^3$ will be chosen in such a way that
\begin{gather*}
\lambda_1^3 \leqslant \lambda_1, \tag{2.2.28} \\
\lambda_2^3 = \lambda_i, \tag{2.2.29} \\
% === Orig. p. 67 [s35R] ===
\lambda_3^3 \leqslant \lambda_{i+1}, \tag{2.2.30} \\
\tilde{\lambda}_1^3 \geqslant \tilde{\lambda}_{i-1}, \tag{2.2.31} \\
\tilde{\lambda}_2^3 = \tilde{\lambda}_i, \tag{2.2.32} \\
\sin^2 \ang(u_2^3; \tilde{u}_2^3) = \sin^2 \ang(u_i; \tilde{u}_i). \tag{2.2.33}
\end{gather*}
Note that (2.2.28) already follows obviously from the preliminary description of $H^3$ and $A^3$ given above by means of the Courant--Fischer variational principle.

Instead of the hypotheses of Theorem~2.2.2, we shall use the stronger conditions
\[
\tilde{\lambda}_{i-1} > \lambda_i, \qquad \tilde{\lambda}_i > \lambda_{i+1}. \tag{2.2.34}
\]
Indeed, when $\tilde{\lambda}_{i-1} \leqslant \lambda_i$ or $\tilde{\lambda}_i \leqslant \lambda_{i+1}$, it is easy to construct an example of $\widetilde{U}$ in which $\sin^2 \ang(u_i, \tilde{u}_i) = 1$. The situation is nontrivial only under conditions (2.2.34).

From (2.2.29)--(2.2.31), (2.2.34), and the variational principle it follows that
\[
\tilde{\lambda}_1^3 > \lambda_2^3 \geqslant \tilde{\lambda}_2^3 > \lambda_3^3.
\]
Under such restrictions, the right-hand side in (2.2.23) is nondecreasing in $\lambda_1^3$, $\lambda_3^3$ and nonincreasing in $\tilde{\lambda}_1^3$. Then the desired statement (2.2.7) follows from (2.2.28)--(2.2.31) and (2.2.23).

It remains to describe precisely the construction of $H^3$ and $\widetilde{U}^3$.

In the case $i = 2$, it is quite simple:
\[
\widetilde{U}^3 = \myspan\{\tilde{u}_1, \tilde{u}_2\}, \qquad H^3 = \myspan\{\tilde{u}_1, \tilde{u}_2, u_2\}.
\]
Obviously, we have
\[
\tilde{u}_1^3 = \tilde{u}_1, \qquad \tilde{u}_2^3 = \tilde{u}_2, \qquad u_2^3 = u_2,
\]
% === Orig. p. 68 [s36L] ===
and, moreover, $\lambda_3^3 \leqslant \lambda_3$ by virtue of the variational principle. Consequently, all relations (2.2.28)--(2.2.33) are satisfied.

In the case of an arbitrary $i \in [2, p]$, the construction becomes somewhat more involved:
\[
\widetilde{U}^3 = \myspan\{\tilde{u}_1^*, \tilde{u}_i\}, \qquad H^3 = \myspan\{\tilde{u}_1^*, \tilde{u}_i, u_i\},
\]
where $\tilde{u}_1^* \in \myspan\{\tilde{u}_1, \ldots, \tilde{u}_{i-1}\} \setminus 0$ is chosen so that $P_1 \tilde{u}_1^* = P_1 \tilde{u}_i$ if $P_1 \tilde{u}_i \neq 0$, and arbitrarily otherwise (here $P_1$ is the orthogonal projector onto $\myspan\{u_1, \ldots, u_{i-1}\}$). Such a vector $\tilde{u}_1^*$ exists because
\[
\operatorname{codim}\{v \in H : P_1 v \in \myspan\{P_1 \tilde{u}_i\}\} = i - 2,
\]
and
\[
\tilde{u}_1^3 = \tilde{u}_1^*, \qquad \tilde{u}_2^3 = \tilde{u}_i, \qquad u_2^3 = u_i.
\]
The inequality $\tilde{\lambda}_1^3 \geqslant \tilde{\lambda}_{i-1}$ is obvious, whereas $\lambda_3^3 \leqslant \lambda_{i+1}$ follows from
\[
u_3^3 \perp \myspan\{u_1, \ldots, u_i\}.
\]

Thus, all relations (2.2.28)--(2.2.33) are satisfied. Lemma~2.2.4, and along with it Theorem~2.2.2, is proved.

\textit{Proof of the attainability of estimate (2.2.7).} Let us show that for a given $i \in [2, p]$ and prescribed numbers $\tilde{\lambda}_{i-1}, \tilde{\lambda}_i$ satisfying the conditions
\[
\lambda_{i-1} \geqslant \tilde{\lambda}_{i-1} > \lambda_i \geqslant \tilde{\lambda}_i > \lambda_{i+1},
\]
there exists a space $\widetilde{U}$ such that (2.2.7) turns into an equality. By virtue of (2.2.23), it follows from the scheme of the proof of Lemma~2.2.4 that for this it suffices that the nonstrict inequalities (2.2.28), (2.2.30), and (2.2.31) turn into equalities for a given $\widetilde{U}$. Let us give such an example of $\widetilde{U}$: let the vectors $\tilde{u}_i$ and $\tilde{u}_{i-1}$ belong to $H^3 = \myspan\{u_1, u_i, u_{i+1}\}$, while the remaining Ritz vectors from $\widetilde{U}$ lie in the orthogonal complement to $H^3$ and coincide with suitable eigenvectors of the operator $A$. The existence of the vectors
% === Orig. p. 69 [s36R] ===
$\tilde{u}_i, \tilde{u}_{i-1}$ described above for prescribed $\lambda_1 \geqslant \tilde{\lambda}_{j-1} > \lambda_j \geqslant \tilde{\lambda}_j > \lambda_{j+1}$ follows % [math note: index j instead of i in original]
from the already mentioned Exercise~10.1.2 in~\citep{Parlett1983}. Consequently, estimate (2.2.7) is sharp on the class of all subspaces $\widetilde{U}$ satisfying the hypotheses of Theorem~2.2.2 (admissible subspaces).

\textit{The proof} of the attainability of estimate (2.2.8) on the class of all admissible subspaces $\widetilde{U}$ and operators $A$ without restrictions on the magnitude of the eigenvalue $\lambda_1$ follows from the preceding arguments and the fact that (2.2.8) is obtained from (2.2.7) as $\lambda_1 \to \infty$. Thus, estimate (2.2.8) is sharp in a narrower sense. However, unlike (2.2.7), it uses the minimum possible information; i.e., any attempt to eliminate any of the parameters $\tilde{\lambda}_{i-1}, \lambda_i, \tilde{\lambda}_i$, or $\lambda_{i+1}$ from (2.2.7) (in the same way as $\lambda_1$ is eliminated in (2.2.8)) leads only to the trivial statement $\sin^2 \ang(u_i, \tilde{u}_i) \leqslant 1$.

\textit{Proof of Theorem 2.2.3.} As usual, in the proof we assume $L = E$ and $A = M$. Let us state Theorem~2.2.3 in the simplest special case of a two-dimensional space.

\textbf{Lemma 2.2.5.} Let $\dim H^2 = 2$ and $\dim \widetilde{U}^2 = 1$; then
\[
\frac{\lambda_1^2 - \tilde{\lambda}_1^2}{\tilde{\lambda}_1^2 - \lambda_{\min}^2} \leqslant \tan^2 \ang(\widetilde{U}^2, U^2), \tag{2.2.35}
\]
where $\lambda_1^2, \lambda_{\min}^2$ are the eigenvalues of the operator $A^2$, $U^2$ is the (one-dimensional) invariant subspace corresponding to $\lambda_1^2$, and $\lambda_{\min}^2$ is the % [math note: \lambda_{\min}^2 in original instead of \tilde{\lambda}_1^2]
Ritz value on the subspace $\widetilde{U}^2$ (also one-dimensional).

In fact, equality even holds in (2.2.35), whose verification is obvious.

In accordance with the technique of reduction to the ``low-dimensional'' case, by analogy with the proof of Theorem~2.2.2, we establish the equivalence of Theorem~2.2.3 and Lemma~2.2.5 in the following formulation.

\textbf{Lemma 2.2.6.} For Theorem~2.2.3 to hold, it is necessary and sufficient that Lemma~2.2.5 hold.

% === Orig. p. 70 [s37L] ===
\textit{Proof of Lemma 2.2.6.} Necessity is obvious. To verify sufficiency, we choose $H^2$ and $\widetilde{U}^2$ as subspaces in $H$ ($\dim H > 2$) so that the following relations hold (for a fixed $i \in [1, p]$):
\begin{gather*}
\lambda_1^2 \geqslant \lambda_i, \tag{2.2.36} \\
\lambda_{\min}^2 \geqslant \lambda_{\min}, \tag{2.2.37} \\
\tilde{\lambda}_1^2 \leqslant \tilde{\lambda}_i, \tag{2.2.38} \\
\tan^2 \ang(\widetilde{U}^2, U^2) \leqslant \tan^2 \ang(\widetilde{U}, U) \tag{2.2.39}
\end{gather*}
(for $i = 1$, equality holds in (2.2.36)). Then estimate (2.2.10) of Theorem~2.2.3 will obviously follow from (2.2.35).

It remains to specify the exact form of $H^2$ and $\widetilde{U}^2$.

We introduce the subspaces
\[
\overline{V}_i = \myspan\{\bar{u}_1, \ldots, \bar{u}_i\}, \qquad V_i = \myspan\{u_1, \ldots, u_i\}. \tag{2.2.40}
\]
We exclude from consideration the situation trivial for estimate (2.2.10),
\[
\tan \ang(\widetilde{U}, U) = \infty.
\]
Then all the vectors $\bar{u}_i$ in (2.2.11) are well defined, and $\dim \overline{V}_i = \dim V_i$.

The Courant--Fischer variational principle yields
\[
\tilde{\lambda}_i \geqslant \bar{\lambda}_i \equiv \min_{v \in \overline{V}_i} \lambda(v).
\]
Choose $\bar{v}_i \in \overline{V}_i \setminus 0$ so that $\lambda(\bar{v}_i) = \bar{\lambda}_j$, and % [math note: \bar{\lambda}_j in original instead of \bar{\lambda}_i]
define $v_i \equiv P\bar{v}_i$, where $P$ is the orthogonal projector onto $U$. We have $v_i \in V_i \setminus 0$; consequently, $\lambda(v_i) \geqslant \lambda_i$. In the case $\bar{v}_i = v_i$, obviously $\tilde{\lambda}_i = \lambda_i$. Henceforth we assume $\bar{v}_i \neq v_i$; then $\bar{v}_i - v_i \in U^\perp \setminus 0$ and $\lambda(\bar{v}_i - v_i) \in [\lambda_{\min}, \lambda_{p+1}]$.

% === Orig. p. 71 [s37R] ===
Set $H^2 \equiv \myspan\{\bar{v}_i, v_i\}$ and $\widetilde{U}^2 = \myspan\{\bar{v}_i\}$. Introduce the orthogonal projector $P^2$ onto $H^2$ and set $A^2 = P^2 A$. From the relations obtained above it follows that the eigenvectors of $A^2$ are $v_i$ and $\bar{v}_i - v_i$, with the corresponding eigenvalues
\[
\lambda_1^2 = \lambda(v_i) \geqslant \lambda_i, \qquad \lambda_{\min}^2 = \lambda(\bar{v}_i - v_i) \in [\lambda_{\min}, \lambda_{p+1}].
\]
To the subspace $\widetilde{U}^2 \subset H^2$ there corresponds the Ritz value
\[
\tilde{\lambda}_1^2 = \lambda(\bar{v}_i) = \bar{\lambda}_i \leqslant \tilde{\lambda}_i.
\]
Moreover, since $\widetilde{U}^2 = \myspan\{\bar{v}_i\}$, we have
\[
\tan^2 \ang(\widetilde{U}^2, U^2) \leqslant \tan^2 \ang(\overline{V}_i, V_i) \leqslant \tan^2 \ang(\widetilde{U}, U). \tag{2.2.41}
\]
Thereby, all conditions (2.2.36)--(2.2.39) are satisfied.

If (2.2.39) is replaced by the first inequality in (2.2.41), then these arguments lead to the estimate
\[
0 \leqslant (\lambda_i - \tilde{\lambda}_i) / (\tilde{\lambda}_i - \lambda_{\min}) \leqslant \tan^2 \ang(\overline{V}_i, V_i), \tag{2.2.42}
\]
which slightly improves (2.2.10) and will be needed subsequently in the proof of Theorem~2.3.1.

The proof of Lemma~2.2.6, and along with it of Theorem~2.2.3, is complete.

\textit{Proof of the attainability of estimate (2.2.10) for some $i \in [1, p]$.} Let us choose $\widetilde{U}$ so that $\tilde{u}_j = u_j$ for $j \neq i$, and $\tilde{u}_i \in \myspan\{u_i, u_{\min}\}$. Then all inequalities in (2.2.36)--(2.2.39) become equalities, and it remains to recall that equality actually holds in (2.2.35) as well. Consequently, estimate (2.2.10) is attainable.

\textit{Proof of Theorem 2.2.4 on the asymptotic closeness of the vectors $\tilde{u}_i$, $\bar{u}_i$, and $u_i^*$.} As usual, in the proof we assume without loss of generality that $L = E$ and $A = M$. Suppose that the vectors $u_i$ are unit vectors and that $(\bar{u}_i, u_j) = \delta_{ij}$. For
% === Orig. p. 72 [s38L] ===
convenience, let us redefine
\[
\varepsilon_i \equiv \sin \ang(u_i, \widetilde{U}) \sim \ang(u_i, \widetilde{U}).
\]

We begin the verification with the statement
\[
\ang(\bar{u}_i, u_i^*) = O(\varepsilon \varepsilon_i),
\]
or, more precisely, even the formally stronger relation
\[
\|\bar{u}_i - u_i^*\| = O(\varepsilon \varepsilon_i) \tag{2.2.43}
\]

The argument is based on the relations
\[
|(\bar{u}_i - u_i^*, u_j^*)| = \varepsilon_i \varepsilon_j |(v_i, v_j)| = O(\varepsilon_i \varepsilon_j), \tag{2.2.44}
\]
in which $v_i = (u_i - u_i^*) / \varepsilon_i$, $\|v_i\| = 1$, $i, j = 1, \ldots, p$. These relations follow from the chain of equalities
\begin{align*}
(\bar{u}_i - u_i^*, u_j^*) &= (\bar{u}_i - u_i^*, u_j) = \\
&= \delta_{ij} - (u_i^*, u_j) = \varepsilon_i (v_i, u_j) = \varepsilon_i \varepsilon_j (v_i, v_j).
\end{align*}
In a similar manner we obtain
\[
(u_i^*, u_j^*) = \delta_{ij} - \varepsilon_i \varepsilon_j (v_i, v_j), \tag{2.2.45}
\]
i.e., the vectors $u_1^*, \ldots, u_p^*$ are ``almost'' orthonormal; consequently, they form a basis in $\widetilde{U}$. Then for an arbitrary unit vector $w \in \widetilde{U}$, on the one hand, we have
\[
w = \sum_{j=1}^p x_j u_j^*, \qquad \sum_{j=1}^p x_j^2 = 1 + O(\varepsilon^2),
\]
% === Orig. p. 73 [s38R] ===
and on the other hand, using (2.2.44), we see that
\begin{align*}
|(\bar{u}_i - u_i^*, w)| &\leqslant \sum_{j=1}^p |x_j| |(\bar{u}_i - u_i^*, u_j)| = % [math note: u_j without asterisk in original; cf. (2.2.44)]
\\
&= \sum_{j=1}^p |x_j| O(\varepsilon \varepsilon_i).
\end{align*}
Taking into account the fact that each of the quantities $|x_j|$ is bounded, we arrive at the asymptotic relation
\[
|(\bar{u}_i - u_i, w)| = O(\varepsilon \varepsilon_i). % [math note: asterisk omitted on u_i^* in original]
\]
Choosing here
\[
w = (\bar{u}_i - u_i^*) / \|\bar{u}_i - u_i^*\|,
\]
we obtain (2.2.43).

The last nontrivial step in the proof of Theorem~2.2.4 consists in verifying the validity of the representation
\[
\tilde{u}_i = \sum_{j=1}^p x_j u_j^*, \qquad x_j = \delta_{ij} + O(\varepsilon \varepsilon_i), \tag{2.2.46}
\]
which, by virtue of (2.2.45), implies the statement
\[
\|\tilde{u}_i - u_i^*\| = O(\varepsilon \varepsilon_i). \tag{2.2.47}
\]

To prove (2.2.46), we define the vector $X = (x_1, \ldots, x_p)^T$ as an eigenvector of the matrix eigenvalue problem of the Rayleigh--Ritz method with respect to the basis $\{u_1^*, \ldots, u_p^*\}$:
\begin{gather*}
\widehat{A} X = \tilde{\lambda} \widehat{E} X, \qquad \widehat{A} \equiv ((Au_i^*, u_j^*)), \qquad \widehat{E} \equiv ((u_i^*, u_j^*)), \tag{2.2.48} \\
i, j = 1, \ldots, p.
\end{gather*}

% === Orig. p. 74 [s39L] ===
From statement (2.2.45) and its analogue in the inner product $(\cdot, *)_A$, we obtain the expressions
\[
\delta_{ij} + O(\varepsilon_i \varepsilon_j) \qquad \text{and} \qquad \lambda_i \delta_{ij} + O(\varepsilon_i \varepsilon_j)
\]
for the entries of the matrices $\widehat{E}$ and $\widehat{A}$, respectively.

Let us show that the eigenvector $X$ of problem (2.2.48) corresponding to the eigenvalue $\tilde{\lambda}_i$ has components $x_i = \delta_{ij} + O(\varepsilon \varepsilon_i)$. % [math note: x_i instead of x_j in original; cf. (2.2.46)]
We use a well-known result of perturbation theory~\citep[\S\,24]{Collatz1968ru} to analyze the perturbation of an eigenvector of a diagonal matrix, setting
\begin{gather*}
\Delta A \equiv \widehat{A} - \diag(\lambda_1, \ldots, \lambda_p), \\
\Delta E \equiv \widehat{E} - \diag(1, \ldots, 1).
\end{gather*}
We have the eigenvalue correction
\[
\tilde{\lambda}_i - \lambda_i = ([\Delta A - \lambda_i \Delta E) e_i, e_i) + O(\varepsilon^4) = O(\varepsilon \varepsilon_i) % [math note: unmatched square bracket [\Delta A - \lambda_i \Delta E) in original]
\]
and the equation for the eigenvector correction
\[
\begin{aligned}
&\diag(\lambda_1 - \lambda_i, \ldots, \lambda_p - \lambda_i)(X - e_i) = \\
&\quad = (\tilde{\lambda}_i - \lambda_i) e_i - (\Delta A - \lambda_i \Delta E) e_i + O(\varepsilon^4),
\end{aligned} \tag{2.2.49}
\]
where $e_i$---the $i$th coordinate vector---is an eigenvector of the matrix $\diag(\lambda_1, \ldots, \lambda_p)$. Since
\[
\|(\Delta A - \lambda_i \Delta E) e_i\| = O(\varepsilon \varepsilon_i),
\]
the norm of the right-hand side of equation (2.2.49) is bounded by $O(\varepsilon \varepsilon_i)$. Then this quantity bounds the solution of (2.2.49) as well; consequently, representation (2.2.46) holds.

% === Orig. p. 75 [s39R] ===
All the remaining assertions of Theorem~2.2.4 are trivial consequences of the already proved relations (2.2.43) and (2.2.47). The proof of Theorem~2.2.4 is complete.

\textit{The proof of Theorem 2.2.5} is based on the observation that the proof of Theorem~2.2.4 does not use the assumption that $\lambda_1, \ldots, \lambda_p$ are consecutive and largest eigenvalues of the operator $A$. Without loss of generality, we assume $L = E$.

Set $A = (\alpha F + \beta G)^{-1} F$ if $\beta \neq 0$, or $A = (\alpha F + \beta G)^{-1} G$ if $\alpha \neq 0$.

This is a self-adjoint operator in the space $H_{\alpha F + \beta G}$ with eigenvectors $u_1, \ldots, u_p, \ldots$, where the vectors $u_1, \ldots, u_p$ correspond to simple eigenvalues, as follows from the hypothesis of Theorem~2.2.5 on the simplicity of the eigenvalues of the problem $Fu = \eta Gu$. Applying the Rayleigh--Ritz method with the subspace $\widetilde{U}$ to the operator $A$ in $H_{\alpha F + \beta G}$, we obtain Ritz vectors $\tilde{u}_i^\eta$ coinciding with the vectors from (2.2.14). From Theorem~2.2.4 we obtain
\[
\ang(\tilde{u}_i^\eta; u_i^\eta) \asymp \varepsilon_i, \qquad \ang(\tilde{u}_i^\eta; \bar{u}_i^\eta) = O(\varepsilon \varepsilon_i), \qquad i = 1, \ldots, p, \tag{2.2.50}
\]
where the vectors $u_1^\eta, \ldots, u_p^\eta$ coincide up to numbering with the vectors $u_1, \ldots, u_p$, and the vectors $\bar{u}_i^\eta \in \widetilde{U}$ are defined % [orig: minor spelling slip in 1986 Russian text for "are defined"]
by the conditions
\[
(\bar{u}_i^\eta, u_j^\eta)_{\alpha F + \beta G} = 0, \qquad i \neq j, \qquad i, j = 1, \ldots, p,
\]
which are equivalent to the conditions
\[
(\bar{u}_i^\eta, u_j^\eta) = 0, \qquad i \neq j, \qquad i, j = 1, \ldots, p.
\]

Comparing the latter equalities with (2.2.11), we verify that the vectors $\bar{u}_i^\eta$ coincide up to numbering with the corresponding vectors $\bar{u}_i$. Then relations (2.2.13) and (2.2.50) lead to (2.2.15). Theorem~2.2.5 is proved.

% =============================================================================
% SOURCE MODULE: part_05_ch2_c.tex
% =============================================================================
% === Orig. p. 76 [s40L] ===
\monosection{2.3}{Abstract subspace iteration method}

Let a subspace $U^0$ ($\dim U^0 = p$) be given, which we regard as an initial approximation to the invariant subspace $U = \myspan\{u_1, \ldots, u_p\}$. Some operator $F_k$, which is a function $F_k(\cdot)$ of the operator $A = L^{-1}M$, will be regarded as the transition operator from $U^0$ to the $k$th approximation---the subspace
\[
U^k = F_k U^0. \tag{2.3.1}
\]

Let us define the quantities
\[
\begin{aligned}
\sigma_j^k &\equiv \max_{j > p} |F_k(\lambda_j)|^2 / |F_k(\lambda_i)|^2, \\ % [math note: in original, \sigma_j^k instead of \sigma_i^k]
\bar{\sigma}_i^k &\equiv \min_{j \in [1, i]} \sigma_j^k.
\end{aligned}
\tag{2.3.2}
\]

For $\bar{\sigma}_p^k < 1$, the convergence rate of the ``iterations'' (2.3.1) can be estimated from the inequality
\[
\tan^2 \ang_L(U^k; U) \leqslant \bar{\sigma}_p^k \tan^2 \ang_L(U^0; U), \tag{2.3.3}
\]
which holds for any finite $\bar{\sigma}_p^k$.

However, in $U^k$ there exist vectors---namely, the vectors from (2.2.11)---that converge faster than the subspace $U^k$ as a whole.

\textbf{Lemma 2.3.1.} Let the vectors $\bar{u}_i^k \in U^k$ be defined by conditions (2.2.11). Then
\[
\tan^2 \ang_L(\bar{u}_i^k, u_i) \leqslant \sigma_i^k \tan^2 \ang_L(U^0, U). \tag{2.3.4}
\]

Of greatest practical interest are error estimates for the Ritz values and Ritz vectors in the subspace $U^k$. From (2.3.4) and the results of Theorem~2.2.4, it is easy to obtain sharp asymptotic estimates for the desired quantities. However, when studying the efficiency of methods for
% === Orig. p. 77 [s40R] ===
mesh problems, such asymptotic estimates are insufficient. It turns out to be possible to generalize to the subspace iterations (2.3.1) those estimates of Theorem~1.1.1 on whose left-hand side only the Rayleigh quotient depends on $u^k$, namely (1.1.9) and (1.1.10).

\textbf{Theorem 2.3.1.} Let $\lambda_i^k$ ($i = 1, \ldots, p$) denote the Ritz values in the subspace $U^k$. They satisfy the estimates
\[
0 \leqslant (\lambda_i - \lambda_i^k) / (\lambda_i^k - \lambda_{\min}) \leqslant \bar{\sigma}_i^k \tan^2 \ang_L(U^0, U). \tag{2.3.5}
\]
If $\bar{\sigma}_p^k < 1$,
\[
\lambda_p^0 \equiv \min_{u \in U^0 \setminus 0} \lambda(u) > \lambda_{p+1} \tag{2.3.6}
\]
and $|F_k(\lambda_1)| \geqslant \ldots \geqslant |F_k(\lambda_p)|$, then the following estimate also holds:
\[
0 \leqslant (\lambda_p - \lambda_p^k) / (\lambda_p^k - \lambda_{p+1}) \leqslant \bar{\sigma}_p^k (\lambda_p - \lambda_p^0) / (\lambda_p^0 - \lambda_{p+1}). \tag{2.3.7}
\]
Each of the estimates (2.3.5) and (2.3.7) is attainable for a given $i \in [1, p]$ in the same sense as the estimates of Theorem~1.1.1.

The overbar on $\bar{\sigma}_i^k$ in (2.3.5) and (2.3.7) cannot, in general, be removed.

We have not succeeded in generalizing estimate (2.3.7) to the case $i < p$.

Error estimates for the Ritz vectors can be obtained by means of (2.3.5), (2.2.9), (2.2.8) or (2.3.5), (2.3.4), (2.2.2).

Finally, note that for $F_k = E$ we have $U^k = U^0$ and $\bar{\sigma}_i^k = 1$, in which case estimate (2.3.5) reduces to estimate (2.2.10) of Theorem~2.2.3.

\textit{Proof of Lemma 2.3.1 and estimate (2.3.3).} As usual, we assume $L = E$ and $A = M$. Let us choose an arbitrary vector $u^0 \in U^0 \setminus 0$ and write the decomposition
\[
u^0 = P u^0 + P^\perp u^0. \tag{2.3.8}
\]
% === Orig. p. 78 [s41L] ===
We have
\[
\begin{aligned}
u^k &= F_k u^0 = F_k P u^0 + F_k P^\perp u^0 \\
&= P F_k P u^0 + P^\perp F_k P^\perp u^0.
\end{aligned}
\tag{2.3.9}
\]
Consequently,
\[
\begin{aligned}
\tan^2 \ang(u^k, U) &= \frac{\|P^\perp F_k P^\perp u^0\|^2}{\|P F_k P u^0\|^2} \\
&\leqslant \frac{\max\limits_{j > p} |F_k(\lambda_j)|^2}{\min\limits_{i \leqslant p} |F_k(\lambda_i)|^2} \frac{\|P^\perp u^0\|^2}{\|P u^0\|^2} = \bar{\sigma}_p^k \tan^2 \ang(u^0, U),
\end{aligned}
\tag{2.3.10}
\]
whence, in particular, it is seen that $u^k \neq 0$. This shows that $\dim U^k = p$, and that the operator $F_k : U^0 \to U^k$ is an isomorphism. Then from (2.3.10) we conclude that
\[
\max_{u^k \in U^k} \tan^2 \ang(u^k, U) \leqslant \bar{\sigma}_p^k \max_{u^0 \in U^0} \tan^2 \ang(u^0, U).
\]

By virtue of the properties of the principal angle between subspaces from \S\,2.1, this inequality coincides with inequality (2.3.3).

To verify (2.3.4), let us choose in (2.3.8) $u^0 = \bar{u}_i^0$, the vector from (2.2.11) belonging to the subspace $U^0$. Then $u^k$ in (2.3.9) is the vector from (2.2.11) belonging to the subspace $U^k$: $u^k = \bar{u}_i^k$. In the denominator of relation (2.3.10) we have
\[
P \bar{u}_i^k = P F_k P \bar{u}_i^0 = F_k(\lambda_i) P \bar{u}_i^0.
\]
Consequently,
\[
\tan^2 \ang(\bar{u}_i^k, u_i) \leqslant \sigma_i^k \tan^2 \ang(\bar{u}_i^0, u_i),
\]
% === Orig. p. 79 [s41R] ===
whence, by virtue of the relations
\[
\tan^2 \ang(\bar{u}_i^0; u_i) = \tan^2 \ang(\bar{u}_i^0; U) \leqslant \tan^2 \ang(U^0; U),
\]
estimate (2.3.4) follows.

This completes the proof of Lemma~2.3.1 and estimate (2.3.3).

Note that similar arguments prove the validity of the estimate
\[
\tan^2 \ang(\overline{V}_i^k; V_i) \leqslant \bar{\sigma}_i^k \tan^2 \ang(\overline{V}_i^0; V_i), \tag{2.3.11}
\]
where $V_i = \myspan\{u_1, \ldots, u_i\}$ and $\overline{V}_i^n = \myspan\{\bar{u}_1^n, \ldots, \bar{u}_i^n\}$, $n = 0, k$. This estimate will be needed below in proving estimate (2.3.5) of Theorem~2.3.1.

\textit{Proof of Theorem 2.3.1.} Without loss of generality we set $L = E$.

Let us establish estimate (2.3.5). Recalling results (2.2.42) and (2.3.11) and combining them, we obtain the inequality
\[
\frac{\lambda_i - \lambda_i^k}{\lambda_i^k - \lambda_{\min}} \leqslant \bar{\sigma}_i^k \tan^2 \ang(\overline{V}_i^0; V_i). \tag{2.3.12}
\]
Estimate (2.3.5) follows immediately from this, since
\[
\tan^2 \ang(\overline{V}_i^0; V_i) = \tan^2 \ang(\overline{V}_i^0; U) \leqslant \tan^2 \ang(U^0; U).
\]

To derive (2.3.7), we will need the following auxiliary lemma.

\textbf{Lemma 2.3.2.} If for a vector $u \in H \setminus 0$ and a real-valued function $F(\lambda)$ the inequalities
\[
\begin{aligned}
\lambda(u) &\in (\lambda_{i+1}, \lambda_i], \\
|F(\lambda)| &\leqslant 1 && \text{for} \quad \lambda \in \{\lambda_{\min}, \ldots, \lambda_{i+1}\}, \\
|F(\lambda)| &\geqslant 1 && \text{for} \quad \lambda \in \{\lambda_i, \ldots, \lambda_1\},
\end{aligned}
\]
% === Orig. p. 80 [s42L] ===
hold, then we have the estimate
\[
\lambda(u) \leqslant \lambda(F(A)u). \tag{2.3.13}
\]

\textit{Proof of Lemma 2.3.2.} Let us represent $F(\lambda)$ in factored form:
\[
F(\lambda) = F_{\lambda_1}(\lambda) \cdot F_{\lambda_2}(\lambda) \cdot \cdots \cdot F_{\lambda_{\min}}(\lambda),
\]
where by definition
\[
\begin{aligned}
F_{\lambda_j}(\lambda) &\equiv 1 && \text{for} \quad \lambda \neq \lambda_j, \\
F_{\lambda_j}(\lambda_j) &= F(\lambda_j), && j \in \{1, 2, \ldots, \min\}.
\end{aligned}
\]

To prove (2.3.13), consider first the special case where
\[
F(\lambda) \equiv 1 \qquad \text{for} \quad \lambda \neq \lambda_j,
\]
i.e., where $F(\lambda) = F_{\lambda_j}(\lambda)$. Define the subspace
\[
H^2 = \myspan\{u, F_{\lambda_j}(A)u\},
\]
which is two-dimensional in the nontrivial case $F_{\lambda_j}(\lambda_j) \neq 1$, and let $P^2$ be the orthogonal projector onto this subspace. The operator $A^2 = P^2 A$ has eigenvectors $u_j$ and $u - u_j$ corresponding to the eigenvalues $\lambda_j$ and $\lambda(u - u_j)$. Note that the following expansions in these eigenvectors hold:
\[
\begin{aligned}
u &= u_j + (u - u_j), \\
F_{\lambda_j}(A)u &= F(\lambda_j)u_j + (u - u_j).
\end{aligned}
\]
% === Orig. p. 81 [s42R] ===
Since the Rayleigh quotients for the operators $A$ and $A^2$ coincide in $H^2$, verifying the inequality
\[
\lambda(u) \leqslant \lambda(F_{\lambda_j}(A)u) \tag{2.3.14}
\]
can in fact be carried out in the two-dimensional space $H^2$, where it becomes obvious. It is also clear that
\[
\lambda(F_{\lambda_j}(A)u) \leqslant \lambda_j \qquad \text{for} \quad j \leqslant i. \tag{2.3.15}
\]
Applying result (2.3.14) successively with account taken of (2.3.15) for $j = i, i - 1, \ldots, 1$, and then applying (2.3.14) for all $j > i$, we obtain the chain of inequalities
\[
\begin{aligned}
\lambda(u) &\leqslant \lambda(F_{\lambda_i}(A)u) \leqslant \ldots \\
&\leqslant \lambda(F_{\lambda_1}(A) \cdot \cdots \cdot F_{\lambda_i}(A)u) \leqslant \ldots \leqslant \lambda(F(A)u).
\end{aligned}
\]

Lemma~2.3.2 is proved.

Now let us proceed directly to the derivation of estimate (2.3.7).

Choose a vector $u^0 \in U^0$ such that
\[
\lambda(F_k(A)u^0) = \lambda_p^k.
\]
Since the operator $F_k(A) : U^0 \to U^k$ is an isomorphism, such a choice is always possible. We have
\[
\lambda(u^0) \geqslant \lambda_p^0 > \lambda_{p+1}.
\]
Suppose that $\lambda(u^0) \in (\lambda_{i+1}, \lambda_i]$ for some $i \in [1, p]$. By virtue of the conditions imposed on the function $F_k(\cdot)$ in Theorem~2.3.1, the hypotheses of Lemma~2.3.2 are satisfied. Applying it gives
\[
\lambda(u^0) \leqslant \lambda(F_k(A)u^0) = \lambda_p^k,
\]
and, invoking the variational principle, we conclude that
\[
\lambda_{p+1} < \lambda_p^0 \leqslant \lambda(u^0) \leqslant \lambda_p^k \leqslant \lambda_p. \tag{2.3.16}
\]

% === Orig. p. 82 [s43L] ===
Define the step function
\[
\widehat{F}(\lambda_i) =
\begin{cases}
\min\limits_{j \in [1, p]} |F_k(\lambda_j)| & \text{for} \quad \lambda_i \geqslant \lambda_p, \\[2ex]
\max\limits_{j \geqslant p+1} |F_k(\lambda_j)| & \text{for} \quad \lambda_i < \lambda_p
\end{cases}
\]
and the vector $\hat{u} \equiv \widehat{F}(A) u^0$.

Applying Lemma~3.3.2 twice, we arrive at the inequalities % [math note: in original, 3.3.2 instead of 2.3.2]
\[
\lambda(u^0) \leqslant \lambda(\hat{u}) \leqslant \lambda_p^k. \tag{2.3.17}
\]

Now everything is ready for constructing a two-dimensional problem. Let us introduce the subspace $H^2 \equiv \myspan\{u^0, \hat{u}\}$, the orthogonal projector $P^2$ onto $H^2$, and the operator $A^2 \equiv P^2 A$. The operator $A^2$ has eigenvectors $P u^0$ and $P^\perp u^0$ corresponding to the eigenvalues $\lambda(P u^0) > \lambda(P^\perp u^0)$, where $P$ is the orthogonal projector onto $U$ and $P^\perp = E - P$. We have the following expansions in these eigenvectors:
\[
\begin{aligned}
u^0 &= P u^0 + P^\perp u^0, \\[1ex]
\hat{u} &= c\bigl(P u^0 + \sqrt{\bar{\sigma}_p^k}\, P^\perp u^0\bigr),
\end{aligned}
\]
where $c \neq 0$ is some constant and $\bar{\sigma}_p^k$ is defined in (2.3.2). Using these expansions, one easily derives (a two-dimensional problem!) the equality
\[
\frac{\lambda(P u^0) - \lambda(\hat{u})}{\lambda(\hat{u}) - \lambda(P^\perp u^0)} = \bar{\sigma}_p^k \frac{\lambda(P u^0) - \lambda(u^0)}{\lambda(u^0) - \lambda(P^\perp u^0)}. \tag{2.3.18}
\]
% === Orig. p. 83 [s43R] ===
Relation (2.3.18) is the simplest special case (for a two-dimensional space) of the nonstrict inequality (2.3.7). On the other hand, taking into account inequalities (2.3.16) and (2.3.17) and replacing $\lambda(P u^0)$ in (2.3.18) by $\lambda_p \leqslant \lambda(P u^0)$ and $\lambda(P^\perp u^0)$ by $\lambda_{p+1} \geqslant \lambda(P^\perp u^0)$, we see that (2.3.7) follows from (2.3.18) by virtue of the monotonic dependence of the function
\[
\frac{\lambda(P u^0) - \lambda(\hat{u})}{\lambda(\hat{u}) - \lambda(P^\perp u^0)} \, \frac{\lambda(u^0) - \lambda(P^\perp u^0)}{\lambda(P u^0) - \lambda(u^0)}
\]
on all of its variables. Thus, statement (2.3.7) and its two-dimensional version (2.3.18) are equivalent. This completes the proof of Theorem~2.3.1.

Examples of attainability of the estimates in Theorem~2.3.1.

The nonstrict inequality (2.3.5) becomes an equality for a given $i \in [1, p]$ when in the subspace $U^0$ all Ritz vectors $\tilde{u}_j^0$, except $\tilde{u}_i^0$, coincide with the corresponding eigenvectors $u_j$, while the vector $\tilde{u}_i^0$ belongs to $\myspan\{u_i, u_{\min}\}$. At the same time, the operator $F_k$ must satisfy
\[
\begin{aligned}
\max_{j \geqslant p+1} |F_k(\lambda_j)| &= |F_k(\lambda_{\min})|, \\
\min_{j \leqslant i} |F_k(\lambda_j)| &= |F_k(\lambda_i)|.
\end{aligned}
\tag{2.3.19}
\]

The nonstrict inequality (2.3.7) becomes an equality when in the subspace $U^0$ all Ritz vectors $\tilde{u}_j^0$ coincide with the corresponding eigenvectors $u_j$, except for the vector $\tilde{u}_p^0$, which must lie in $\myspan\{u_p, u_{p+1}\}$. At the same time, the operator $F_k$ must satisfy the condition
\[
\max_{j > p+1} |F_k(\lambda_j)| = |F_k(\lambda_{p+1})|. \tag{2.3.20} % [math note: in original, j > p+1 instead of j \geqslant p+1 (or j > p)]
\]

Attainability is easily verified directly by following the proof of Theorem~2.3.1.

% === Orig. p. 84 [s44L] ===
\monosection{2.4}{Simple subspace iteration with a constant shift and iterations with Chebyshev parameters}

For $i = 1, \ldots, p$, let us introduce the quantities
\[
\bar{\sigma}_i(F_k) \equiv \max_{\lambda \in [\lambda_{\min}, \lambda_{p+1}]} |F_k(\lambda)|^2 / \min_{j \leqslant i} |F_k(\lambda_i)|^2. \tag{2.4.1} % [math note: in original, \lambda_i instead of \lambda_j under \min_{j \leqslant i}]
\]
From definitions (2.3.2) we obtain the obvious inequalities
\[
\bar{\sigma}_i^k \leqslant \bar{\sigma}_i(F_k), \quad i \in [1, p].
\]
Recalling now statements (2.3.5) and (2.3.7) of Theorem~2.3.1, we conclude that the dependence of the convergence rate of the iterations $U^k = F_k U^0$ on the choice of $F_k$ can be characterized by the quantities $\bar{\sigma}_i(F_k)$. This naturally leads to the following definition.

\textbf{Definition of optimality} of subspace iterations. We say that a method $U^k = F_k U^0$ with a given transition operator $F_k$ is optimal over a certain class of operators $\{F_k\}$ if, at the corresponding function $F_k(\cdot)$, all the quantities $\bar{\sigma}_i(F_k)$, $i \in [1, p]$, simultaneously attain their minimum values over the corresponding class of functions $\{F_k(\cdot)\}$.

Note that a similar definition is easily formulated for the vector iterations of Chapter~1.

Let us consider specific examples of optimal subspace iteration methods with an a priori choice of iteration parameters.

Suppose that in (2.3.1) the operator $F_k$ does not depend on $U^0$ and is a polynomial in $A = L^{-1}M$ of degree $k$: $F_k = F_k(A)$. For example,
\[
U^{n+1} = (A - \gamma^n E) U^n, \quad n = 0, 1, \ldots, k - 1. \tag{2.4.2}
\]

Consider the problem of a priori optimization of the choice of parameters $\gamma^n$---the roots of this polynomial---from the condition of minimizing the quantities in (2.4.1).

% === Orig. p. 85 [s44R] ===
In the case of a constant shift, the optimal choice of $\gamma^n$ is
\[
\gamma^n = \gamma \equiv (\lambda_{p+1} + \lambda_{\min}) / 2, \tag{2.4.3}
\]
for which all $\bar{\sigma}_i(F_k)$ simultaneously attain their minimum values
\begin{gather*}
\bar{\sigma}_i(F_k) = \bar{\sigma}_i^k = [(1 - \xi_i) / (1 + \xi_i)]^{2k}, \tag{2.4.4} \\
\xi_i \equiv (\lambda_i - \lambda_{p+1}) / (\lambda_i - \lambda_{\min}). \tag{2.4.5}
\end{gather*}

In the case of variable shifts in (2.4.2), the optimal choice of $\gamma^n$ is
\[
\gamma^n = \{\lambda_{p+1} + \lambda_{\min} - \tau^n (\lambda_{p+1} - \lambda_{\min})\} / 2, \tag{2.4.6}
\]
where $\tau^n \in [-1, 1]$ are all $k$ roots of the polynomial $T_k(\tau)$. In this case, $F_k(\lambda)$ deviates least from zero on the interval $[\lambda_{\min}, \lambda_{p+1}]$, and all $\bar{\sigma}_i(F_k)$ simultaneously attain their minimum values
\[
\bar{\sigma}_i(F_k) = \bar{\sigma}_i^k = \bigl|T_k\{(1 + \xi_i) / (1 - \xi_i)\}\bigr|^{-2}, \tag{2.4.7}
\]

with $\xi_i$ defined in (2.4.5). The proof of this fact was suggested by V.\,I.~Lebedev \citep{KnyazevLebedev1983} and is given below.

A three-term analogue of the Chebyshev iterations (2.4.2), (2.4.6), which is a block version of the vector iterations (1.2.1), (1.2.7), is widely known; see, for example, \citep{Rutishauser1969,WilkinsonReinsch1976ru,Parlett1983}. The preceding arguments show that this method is also optimal over the class of methods whose transition operators are polynomials in $A$.

% === Orig. p. 86 [s45L] ===
Thus, in the sense of the formal definition given above, these subspace iterations are optimal. This resolves the doubt expressed in \citep{Parlett1983} concerning the possibility of a rigorous proof of optimality for such methods. It is important to note that the proposed way of defining optimality is analogous to the well-known one for Chebyshev iterative methods for solving systems of equations, which confirms the appropriateness of this new definition.

Convergence rate estimates for algorithms (2.4.2) are obtained as a corollary of Theorem~2.3.1. In estimates (2.3.5) and (2.3.7) of that theorem, one should take the quantities $\bar{\sigma}_i^k$ from (2.4.4) when applying algorithm $\{(2.4.2), (2.4.3)\}$ and from (2.4.7) for the algorithm with Chebyshev parameters $\{(2.4.2), (2.4.6)\}$. Moreover, the estimates also turn out to be attainable, since conditions (2.3.19) and (2.3.20) are satisfied for the functions $F_k(\cdot)$ generated by methods $\{(2.4.2), (2.4.3)\}$ and $\{(2.4.2), (2.4.6)\}$. % [orig: (2.4.2.)]

In implementing the two-term scheme with Chebyshev parameters $\{(2.4.2), (2.5.6)\}$, no essentially new difficulties arise compared with the corresponding vector iterations $\{(1.2.1), (1.2.7)\}$. % [math note: in original, (2.5.6) instead of (2.4.6)]
All the remarks of \S\,1.2 remain valid regarding the possibility of ordering the parameters $\tau^n$ in (2.4.6) to ensure stability, computing improved bounds for $\lambda_{p+1}$ and $\lambda_{\min}$ in the course of the iteration, etc.\ (see also \citep{WilkinsonReinsch1976ru}).

Let us mention here the paper \citep{AklDilgerIrons1982} on the practical application of method (2.4.2) with variable (though nonoptimal) shifts. % [bib note: third author Irons omitted in original in-text citation]

The \textit{proof} of optimality for the methods with Chebyshev parameters obviously reduces to verifying the following statement on the Chebyshev polynomial.

\textbf{Lemma 2.4.1.} Over all polynomials $F_k(\cdot)$ of degree at most $k$, the quantities $\bar{\sigma}_i(F_k)$, $i = 1, \ldots, p$, simultaneously attain their minimum values at the polynomial
\[
F_k^*(\lambda) = T_k\left(\frac{\lambda_{p+1} + \lambda_{\min} - 2\lambda}{\lambda_{p+1} - \lambda_{\min}}\right), \tag{2.4.8}
\]
i.e., at the Chebyshev polynomial of the first kind of degree $k$ mapped to the interval $[\lambda_{\min}, \lambda_{p+1}]$.

% === Orig. p. 87 [s45R] ===
\textit{Proof of Lemma 2.4.1.} We use the classical fact that the minimum of each of the quantities
\[
\sigma_i(F_k) \equiv \max_{\lambda \in [\lambda_{\min}, \lambda_{p+1}]} |F_k(\lambda)|^2 / |F_k(\lambda_i)|^2,
\]
$i \in [1, p]$, is attained at the polynomial $F_k^*(\cdot)$ from (2.4.8) and equals
\[
\min_{\{F_k\}} \sigma_i(F_k) = |F_k^*(\lambda_i)|^{-2}.
\]
Consequently,
\[
\max_{j \leqslant i} \min_{\{F_k\}} \sigma_j(F_k) = |F_k^*(\lambda_i)|^{-2}.
\]
Substituting this expression for the maximin into the well-known inequality relating the maximin and minimax,
\[
\max_{j \leqslant i} \min_{\{F_k\}} \sigma_j(F_k) \leqslant \min_{\{F_k\}} \max_{j \leqslant i} \sigma_j(F_k),
\]
we obtain
\[
|F_k^*(\lambda_i)|^{-2} \leqslant \min_{\{F_k\}} \max_{j \leqslant i} \sigma_j(F_k).
\]
Using definition (2.4.1), we have
\[
|F_k^*(\lambda_i)|^{-2} \leqslant \min_{\{F_k\}} \bar{\sigma}_i(F_k).
\]
Statement (2.4.8) of the lemma now follows from the fact that
\[
|F_k^*(\lambda_i)|^{-2} = \bar{\sigma}_i(F_k^*).
\]
This completes the proof of Lemma~2.4.1.

% === Orig. p. 88 [s46L] ===
\monosection{2.5}{Block method of steepest ascent}

Let an initial approximation---a subspace $U^0$, $\dim U^0 = p$---be given. The iterative scheme of the block method of steepest ascent is defined as follows: in the subspace
\[
U^n + A U^n, \quad A = L^{-1}M, \quad n = 0, 1, \ldots, \tag{2.5.1}
\]
whose dimension, in general, ranges from $p$ to $2p$, the Rayleigh--Ritz method is carried out to extract $p$ Ritz vectors
\[
u_1^{n+1}, \ldots, u_p^{n+1},
\]
corresponding to the $p$ largest Ritz values
\[
\lambda_1^{n+1} \geqslant \ldots \geqslant \lambda_p^{n+1},
\]
and, finally, the subspace $U^{n+1}$ is defined by the equality
\[
U^{n+1} = \myspan\{u_1^{n+1}, \ldots, u_p^{n+1}\}.
\]
Details can be found, for example, in \citep{CullumWilloughby1985}.

Theorem~2.3.1 is not directly applicable. However, if starting from the initial approximation $U^n$ one performs a single iteration of method (2.4.2), (2.4.3) and denotes the Ritz values in the resulting subspace $(A - \gamma^n E) U^n$ by
\[
\tilde{\lambda}_1^{n+1} \geqslant \ldots \geqslant \tilde{\lambda}_p^{n+1},
\]
then Theorem~2.3.1 yields the estimate
\[
\frac{\lambda_i - \tilde{\lambda}_i^{n+1}}{\lambda_i^{n+1} - \lambda_{\min}} \leqslant \bar{\sigma}_i \tan^2 \ang_L(U^n; U), \tag{2.5.2} % [math note: in original, the tilde over \lambda_i^{n+1} in the denominator is omitted]
\]
$i \in [1, p]$, and, for $\lambda_p^n > \lambda_{p+1}$, also the estimate
\[
\frac{\lambda_p - \tilde{\lambda}_p^{n+1}}{\tilde{\lambda}_p^{n+1} - \lambda_{p+1}} \leqslant \bar{\sigma}_p \frac{\lambda_p - \lambda_p^n}{\lambda_p^n - \lambda_{p+1}}, \tag{2.5.3}
\]
% === Orig. p. 89 [s46R] ===
where $\lambda_p^n$ is the smallest Ritz value for the subspace $U^n$,
\begin{gather*}
\bar{\sigma}_i = (1 - \xi_i)^2 / (1 + \xi_i)^2, \tag{2.5.4} \\
\xi_i = (\lambda_i - \lambda_{p+1}) / (\lambda_i - \lambda_{\min}) \tag{2.5.5}
\end{gather*}
(cf.\ (2.4.4), (2.4.5)). Invoking the inequalities $\tilde{\lambda}_i^{n+1} \leqslant \lambda_i^{n+1}$, which follow from the variational principle, we note that $\tilde{\lambda}_i^{n+1}$ can be replaced by $\lambda_i^{n+1}$ in estimates (2.5.2) and (2.5.3). It is now natural to try to derive estimates in terms of the characteristics of the subspace $U^0$. An inequality of the form (2.5.2) is unsuitable for this purpose, since its left- and right-hand sides are expressed in different terms. From inequality (2.5.3) we obtain, for $\lambda_p^0 > \lambda_{p+1}$, that
\begin{gather*}
\frac{\lambda_p - \lambda_p^k}{\lambda_p^k - \lambda_{p+1}} \leqslant \bar{\sigma}_p^k \frac{\lambda_p - \lambda_p^0}{\lambda_p^0 - \lambda_{p+1}}, \tag{2.5.6} \\
\bar{\sigma}_p^k = [(1 - \xi_p) / (1 + \xi_p)]^{2k}, \tag{2.5.7}
\end{gather*}
since $\xi_p$ in (2.5.5) does not depend on $U^n$ (cf.\ (2.4.4), (2.4.5), and (1.3.3)).

Estimates of the convergence rate of $\lambda_i^k$ to $\lambda_i$ for $i < p$ of the type (2.5.6) with $\bar{\sigma}_i^k$ from (2.4.4) are unknown to the author. Instead, it is not difficult to obtain the weaker estimates
\[
\frac{\lambda_i - \lambda_i^k}{\lambda_i^k - \lambda_{i+1}} \leqslant \hat{\sigma}_i^k \frac{\lambda_i - \lambda_i^0}{\lambda_i^0 - \lambda_{i+1}}, \quad i = 1, \ldots, p, \tag{2.5.8}
\]
% === Orig. p. 90 [s47L] ===
\[
\begin{aligned}
\hat{\sigma}_i^k &\equiv [(1 - \hat{\xi}_i) / (1 + \hat{\xi}_i)]^{2k}, \\
\hat{\xi}_i &\equiv (\lambda_i - \lambda_{i+1}) / (\lambda_i - \lambda_{\min})
\end{aligned}
\]
for $\lambda_i^0 > \lambda_{i+1}$ (cf.\ the definition of $\xi_i$ in (2.4.5)). Note, however, that for a given $i \in [1, p)$ this estimate does not involve $p$, so it remains unclear how the error quantity
\[
(\lambda_i - \lambda_i^k) / (\lambda_i^k - \lambda_{i+1})
\]
behaves as $p$ increases. For instance, when $i = 1$, estimate (2.5.8) coincides with the convergence rate estimate (1.3.3) for the vector method of steepest ascent in \S\,1.3. Thus, the question of whether the block version of the method of steepest ascent has a superior convergence rate compared with the vector version remains open. Note that in the asymptotic regime (as $U^0 \to U$) this question can be answered in the affirmative on the basis of the asymptotic properties of Ritz vectors stated in Theorem~2.2.4.

Alongside (2.5.1), another, less expensive version of the block method of steepest ascent is also known: in the subspace $U^n$, $\dim U^n = p$, one determines the Ritz vectors $u_1^n, \ldots, u_p^n$, performs one iteration of the method of steepest ascent on each of them,
\[
\hat{u}_i^{n+1} = A u_i^n - \alpha_i^n u_i^n,
\]
where $\alpha_i^n$ is chosen as in (1.3.1), and forms the subspace
\[
U^{n+1} = \myspan\{\hat{u}_1^{n+1}, \ldots, \hat{u}_p^{n+1}\};
\]
see, for example, \citep{LongsineMcCormick1980}. A similar algorithm was first proposed in \citep{Samokish1958}, where the vectors $u_i^n$ in $U^n$ were chosen by Gram--Schmidt orthogonalization rather than by the Rayleigh--Ritz method. The possibility of estimating the asymptotic convergence rate as $U^0 \to U$ was also noted there. In the general case, the validity of this algorithm has apparently not been justified to date.

% === Orig. p. 91 [s47R] ===
\monosection{2.6}{Block Lanczos-type and conjugate gradient methods}

Let us define the ``block'' Krylov subspace
\[
\mathcal{K}^k(U^0) \equiv U^0 + A U^0 + \cdots + A^k U^0. \tag{2.6.1}
\]
Theoretically, the block Lanczos method consists in finding the subspace $U^k \subset \mathcal{K}^k(U^0)$ ($\dim U^k = p$) spanned by the $p$ Ritz vectors of problem (0.1) in the subspace $\mathcal{K}^k(U^0)$ corresponding to the $p$ largest Ritz values $\lambda_1^k \geqslant \ldots \geqslant \lambda_p^k$; see, for example, \citep{Parlett1983,Ikramov1982,GolubVanLoan1983,CullumWilloughby1985}. It is easy to estimate the convergence rate of this method using the following observation. By the Courant--Fischer variational principle, the values $\lambda_i^k$ are no smaller than those obtained in the Chebyshev subspace iterations $\{(2.4.2), (2.4.6)\}$. Therefore, for the block Lanczos method, the exact same estimates of Theorem~2.2.5 with $\bar{\sigma}^k$ from (2.4.7) hold. % [math note: in original, Theorem 2.2.5 with \bar{\sigma}^k instead of Theorem 2.3.1 with \bar{\sigma}_i^k]
Estimate (2.3.5) slightly refines the well-known result of Underwood (see, e.g., \citep{Ikramov1982,GolubVanLoan1983}).

An iterative version of the block Lanczos method is known (see, e.g., \citep{CullumWilloughby1985}), in which the process is terminated every $k$ steps and the approximation $U^k$ is used as the initial one at the next stage. For such a version, estimate (2.3.7), which can be applied recursively, is particularly suitable. Convergence rate estimates toward $\lambda_i$, $i < p$, of the type (2.3.7) with $\bar{\sigma}_i^k$ from (2.4.7) are unknown to the author. Note that the block method of steepest ascent of \S\,2.5 can be interpreted as the iterative version of the block Lanczos method with a single ``inner'' iteration; cf.\ (2.5.1) and (2.6.1).

Block versions of conjugate gradient methods have not been sufficiently justified theoretically. By analogy with the case of vector iterations, it is natural to conjecture that none of these versions (including the one in \citep{CullumWilloughby1985}) is equivalent to the block Lanczos method, even though the vectors involved in $k$ iterations of the block conjugate gradient method form a basis of the subspace $\mathcal{K}^k(U^0)$. % [orig: period missing at right margin after \mathcal{K}^k(U^0)]
The lack of global optimality prevents us here
% === Orig. p. 92 [s48L] ===
from using the convergence rate estimates of the block Lanczos method. Probably, only as $U^0 \to U$ do all the versions become asymptotically equivalent.

A block analogue of the locally optimal scheme (1.5.6) is easily justified:
\[
V^{n+1} = A U^n + U^n + U^{n-1}, \quad U^{n+1} \subset V^{n+1}, \quad n = 0, 1, \ldots, k - 1, \tag{2.6.2}
\]
where $U^{n+1}$ is chosen as the linear span of the Ritz vectors for problem (0.1) in the subspace $V^{n+1}$ corresponding to the $p$ largest Ritz values $\lambda_1^{n+1} \geqslant \ldots \geqslant \lambda_p^{n+1}$. Obviously, these values are no smaller than those obtained in a single iteration of the block method of steepest ascent with the initial approximation $U^n$ (see \S\,2.5). Therefore, estimates (2.5.6)--(2.5.8) also hold for algorithm (2.6.2).

Returning to the block Lanczos method, let us consider the possibility of using as the initial approximation $U^0$ the Krylov subspace $U^0 = \myspan\{u^0, \allowbreak A u^0, \allowbreak \ldots, \allowbreak A^{p-1} u^0\}$. In this case, the subspace $\mathcal{K}^k(U^0)$ coincides with the subspace $\mathcal{K}^{k+p-1}(u^0)$ and has, in general, dimension $k + p$ rather than $(k + 1)p \geqslant k + p$. However, the picture changes dramatically when actual computations are carried out in the presence of roundoff errors. Suppose that the following implementation of the block Lanczos method is used. One performs simple subspace iteration:
\[
u_i^{n+1} = (A - \lambda_{\min} E) u_i^n, \quad n = 0, \ldots, k - 1, \quad i = 1, \ldots, p, \tag{2.6.3}
\]
where $\{u_1^0, \ldots, u_p^0\}$ is a basis of the subspace $U^0$, computes the inner products $(u_i^n, u_j^n)_L$ and $(u_i^{n+1}, u_j^n)_L$, $i, j = 1, \ldots, p$, and assembles from them the two matrices of the generalized eigenvalue problem of the Rayleigh--Ritz method for the subspace
\[
\mathcal{K}^k(U^0) = \myspan\{u_1^0, \ldots, u_p^0, \ldots, u_1^k, \ldots, u_p^k\}.
\]
In the case of (2.6.3), in actual computations these two matrices will be nearly (but not exactly!) singular, and in the iterative scheme (2.6.4) it actually suffices to carry out computations only for $i = p$. % [math note: formula (2.6.4) does not exist in Sect. 2.6; by meaning, (2.6.3) is intended, while the beginning of the sentence refers to choosing U^0 as a Krylov subspace]
Thus, $k + p - 1$ multiplications by the operator $A - \lambda_{\min} E$ are performed, and $p + 1$ vectors need to be stored in computer memory. One may conjecture that such a method will yield more accurate approximations to closely spaced eigenvalues $\lambda_1, \ldots, \lambda_p$ than $k + p - 1$ steps of the Lanczos method of \S\,1.4.

% === Orig. p. 93 [s48R] ===
\monosection{2.7}{Orthogonalization against previously computed eigenvectors}

Let $\widetilde{V}$ be the linear span of the already computed approximations to eigenvectors of problem (0.1), i.e., the subspace $\widetilde{V}$ is close to some invariant subspace $V$ of the operator $A = L^{-1}M$. The eigenvectors not yet known lie in $V^{\perp_L}$, the orthogonal complement to $V$ in $H_L$. One can take advantage of this when computing them by iterative methods, supplementing the scheme with orthogonalization. Take, as an example, the simple iteration (1.2.1) with orthogonalization:
\[
u^{n+1} = \widetilde{P}^\perp A u^n - \gamma^n u^n, \quad n = 0, 1, \ldots, \quad u^0 \in \widetilde{V}^\perp, \tag{2.7.1}
\]
where $\widetilde{P}^\perp$ is the orthogonal projector onto the orthogonal complement to $\widetilde{V}$. It is important in which inner product orthogonality is defined.

The sought solutions lie in $V^{\perp_L}$, and method (2.7.1) ensures that $u^n \in \widetilde{V}^\perp$. It is natural to require that these subspaces be close whenever $\widetilde{V}$ is close to $V$ in $H_L$, i.e., say, when the quantity $\varepsilon = \tan \ang_L(\widetilde{V}; V)$ is small. Let some operator $K = K^* > 0$ define the metric of $H_K$ in which $\widetilde{V}^{\perp_K}$ is defined. Note that
\[
\widetilde{V}^{\perp_K} = (L^{-1} K \widetilde{V})^{\perp_L}, \tag{2.7.2}
\]
and use the property of gaps between subspaces (see \S\,2.1) stating that the gap between two subspaces is equal to the gap between their orthogonal complements. Then
\[
\tan \ang_L(\widetilde{V}^{\perp_K}; V^{\perp_L}) = \tan \ang_L(L^{-1} K \widetilde{V}; V). \tag{2.7.3}
\]
From (2.7.3) it is seen that the most natural choice $K = L$ ensures
\[
\tan \ang_L(\widetilde{V}^{\perp_L}; V^{\perp_L}) = \varepsilon. \tag{2.7.4}
\]
% === Orig. p. 94 [s49L] ===
If $V$ corresponds to the largest eigenvalues, a slight gain is achieved for $K = M - \lambda_{\min} L$, when
\[
\tan \ang_L\bigl((L^{-1}M - \lambda_{\min} E)\widetilde{V}; V\bigr) < \tan \ang_L(\widetilde{V}; V);
\]
see (2.3.3).

The choice of the inner product of $H_K$ fixes the subspace $\widetilde{V}^{\perp_K}$, whereas a specific projector $\widetilde{P}^{\perp_K}$ can be prescribed independently.

For definiteness, below we set $K = L$.

How does the convergence rate of the iterations change under additional projection?

The simplest situation to analyze is when $\widetilde{V} = V$. In this case, $A u \in \widetilde{V}^{\perp_L}$ for $u \in \widetilde{V}^{\perp_L}$, and for any choice of $\widetilde{P}^{\perp_L}$ we have $\widetilde{P}^{\perp_L} A u = A u$; that is, in iterations (2.7.1) all $u^n \in \widetilde{V}^{\perp_L}$, even if the projection is omitted altogether (assuming, of course, the absence of roundoff errors). Method (2.7.1) for problem (0.1) is equivalent to method (1.2.1) for the problem
\[
L^{-1}M\big|_{V^{\perp_L}} u = \lambda^\perp u, \quad u \in V^{\perp_L}, \tag{2.7.5}
\]
and its convergence rate is easily estimated by analogy with the results of \S\,1.2.

If $\widetilde{V} \neq V$, then for most operators $\widetilde{P}^{\perp_L}$ the transition operator from $u^n$ to $u^{n+1}$ in (2.7.1) fails to be self-adjoint on the set $\widetilde{V}^{\perp_L}$ in the inner product of $H_L$. To ensure self-adjointness, it suffices to take $\widetilde{P}^{\perp_L}$ orthogonal in $H_L$: $\widetilde{P}^{\perp_L} = \widetilde{P}_L^{\perp_L}$.

Let us show that in this case the approach of backward error analysis is applicable. One may regard method (2.7.1) as being equivalent to method (1.2.1) applied, instead of (0.1), to the problem
\[
\widetilde{P}_L^{\perp_L} L^{-1}M\big|_{\widetilde{V}^{\perp_L}} u = \tilde{\lambda}^\perp u, \quad u \in \widetilde{V}^{\perp_L}, \tag{2.7.6}
\]
% === Orig. p. 95 [s49R] ===
in which the operator $\widetilde{P}_L^{\perp_L} L^{-1}M$ is self-adjoint in the space $H_L$. Applying the theory of Chapter~1 yields convergence rate estimates for (2.7.1) in terms of problem (2.7.6). To return to the original notation of problem (0.1), it suffices to establish, on the basis of (2.7.4), how much the eigenvalues $\tilde{\lambda}_i^\perp$ and eigenvectors of problem (2.7.6) differ from the corresponding ones in problem (2.7.5), i.e., to estimate the error of the Rayleigh--Ritz method for problem (0.1) in the subspace $\widetilde{V}^{\perp_L}$. If $V$ contains all eigenvectors of problem (0.1) corresponding to the largest eigenvalues, then $V^{\perp_L}$ corresponds only to the smallest ones, and one can apply analogues of Theorems~2.2.1--2.2.3 with $U = V^{\perp_L}$. In the case of an arbitrary subspace, more general results are needed, for example,
\[
-(\lambda_i^\perp - \lambda_{\min})\varepsilon^2 \leqslant \tilde{\lambda}_i^\perp - \lambda_i^\perp \leqslant (\lambda_1 - \lambda_i^\perp)\varepsilon^2,
\]
and Theorem~2.2.1 reformulated for an arbitrary subspace \citep{Saad1983}: if the eigenvalue $\lambda_i^\perp$ corresponds to the eigenvector $u_i^\perp \in V^{\perp_L}$ and the Ritz value $\tilde{\lambda}_i^\perp$ corresponds to the Ritz vector $\tilde{u}_i^\perp \in \widetilde{V}^{\perp_L}$, then
\[
\begin{aligned}
&\sin^2 \ang_L(u_i^\perp; \tilde{u}_i^\perp) \leqslant (1 + \gamma^2 / d^2) \sin^2 \ang_L(u_i^\perp; \widetilde{V}^{\perp_L}), \\[1ex]
&\gamma \equiv \|\widetilde{P}_L L^{-1}M (E - \widetilde{P}_L)\|_L \\[1ex]
&d = \min_{\tilde{\lambda}_j^\perp \neq \tilde{\lambda}_i^\perp} |\tilde{\lambda}_j^\perp - \lambda_i| % [math note: in original, \lambda_i without superscript \perp]
\end{aligned}
\]
where the minimum is taken over all Ritz values $\tilde{\lambda}_j^\perp$ for the subspace $\widetilde{V}^{\perp_L}$ and $\widetilde{P}_L$ is the orthogonal projector in $H_L$ onto $\widetilde{V}$.

In \citep{Saad1983}, a weaker statement is given: there exists a Ritz vector $\tilde{u}_i^\perp$ such that\ldots\ However, from the proof in \citep{Saad1980} it is easy to see that the statement actually holds for any Ritz vector $\tilde{u}_i^\perp$, as stated above.

% === Orig. p. 96 [s50L] ===
\monosection{2.8}{Another way of reducing the generalized eigenvalue problem to a standard one}

Using, for $M > 0$, the operator $A = -M^{-1}L$ in $H_M$ instead of the operator $A = L^{-1}M$ in $H_L$ was discussed in detail above in \S\,1.6 for the case of vector iterations. The case of subspace iterations involves no essential differences, and carrying out this substitution in \S\S\,2.2--2.6 presents no difficulty.

We formulate the subspace analogue of relation (1.6.8) linking the angles in $H_M$ and $H_L$ as a lemma.

\textbf{Lemma 2.8.1.} For $M > 0$, the following statement holds:
\[
\tan^2 \ang_M(\widetilde{U}; U) / \tan^2 \ang_L(\widetilde{U}; U) \in [\lambda_{\min}/\lambda_1; \lambda_{p+1}/\lambda_p]. \tag{2.8.1}
\]

\textit{Proof of Lemma 2.8.1.} Let us establish the equality
\[
\tan^2 \ang_M(\widetilde{U}; U) = \tan^2 \ang_L\bigl(\sqrt{L^{-1}M}\,\widetilde{U}; U\bigr). \tag{2.8.2}
\]
We have the representation (see \S\,2.1):
\[
\tan^2 \ang_M(\widetilde{U}; U) = \max_{\tilde{u} \in \widetilde{U}} \|(E - P)\tilde{u}\|_M^2 / \|P\tilde{u}\|_M^2, \tag{2.8.3}
\]
where $P$ is the projector onto the subspace $U$, which is simultaneously orthogonal in $H_M$ and $H_L$. To pass to the angle in $H_L$, it suffices to replace the norms of the space $H_M$ in (2.8.3) by the norms of the space $H_L$. To do this, we write
\[
\|(E - P)\tilde{u}\|_M^2 = \bigl\|\sqrt{L^{-1}M}\,(E - P)\tilde{u}\bigr\|_L^2 = \bigl\|(E - P)\sqrt{L^{-1}M}\,\tilde{u}\bigr\|_L^2.
\]
The operators $\sqrt{L^{-1}M}$ and $(E - P)$ commute because $P$ is the projector onto an invariant subspace of the operator $L^{-1}M$.

Transforming the denominator on the right-hand side of (2.8.3) similarly, we arrive at (2.8.2).

% === Orig. p. 97 [s50R] ===
Now, taking (2.8.2) into account, we obtain statement (2.8.1) by applying estimate (2.3.3) twice. The first time, we set
\[
U^0 = \widetilde{U}; \quad U^k = \sqrt{L^{-1}M}\,\widetilde{U}, \quad F_k = \sqrt{L^{-1}M}, \quad \bar{\sigma}_p^k = \lambda_{p+1}/\lambda_p.
\]
The second time, we set
\[
U^0 = \sqrt{L^{-1}M}\,\widetilde{U}; \quad U^k = \widetilde{U}, \quad F_k = \sqrt{M^{-1}L}, \quad \bar{\sigma}_p^k = \lambda_1/\lambda_{\min}.
\]
These two values of $\bar{\sigma}_p^k$ form the right and left endpoints of the interval in statement (2.8.1). Lemma~2.8.1 is proved.

The fact that $U = \myspan\{u_1, \ldots, u_p\}$ plays an essential role here. If instead it is known only that the subspace $U$ is invariant under the operator $L^{-1}M$, then estimate (2.8.1) is replaced by the estimate
\[
\tan^2 \ang_M(\widetilde{U}; U) / \tan^2 \ang_L(\widetilde{U}; U) \in [\lambda_{\min}/\lambda_1, \lambda_1/\lambda_{\min}], \tag{2.8.4}
\]
and the interval in (2.8.4) cannot, in general, be narrowed. To prove (2.8.4), it suffices to use representation (2.8.3) and replace the norms in $H_M$ by the norms in $H_L$ via the obvious inequalities
\[
\lambda_{\min} \|\cdot\|_L \leqslant \|\cdot\|_M \leqslant \lambda_1 \|\cdot\|_L. % [math note: in original, squares on the norms are omitted: should be \|\cdot\|_L^2 \leqslant \|\cdot\|_M^2 \leqslant \lambda_1 \|\cdot\|_L^2]
\]

For arbitrary subspaces $V$ and $W$, statements of the type (2.8.4) fail to hold in general. For example, the condition
\[
\tan^2 \ang_L(V; W) < \infty
\]
does not necessarily imply that
\[
\tan^2 \ang_M(V; W) < \infty.
\]
From this we conclude that in such a situation the tangent of the angle is an unfortunate choice of quantity, and should be abandoned in favor of, say, the sine:
\[
\sin^2 \ang_M(V; W) / \sin^2 \ang_L(V; W) \in [\lambda_{\min}/\lambda_1, \lambda_1/\lambda_{\min}]. % [orig: period omitted at end of chapter]
\]

% =============================================================================
% SOURCE MODULE: part_06_ch3.tex
% =============================================================================
% === Orig. p. 98 [s51L] ===
\monochapter{3}{The Temple--Lehmann Method for Two-Sided Eigenvalue Bounds}

In this chapter we study the Temple--Lehmann method---a generalized version of the Rayleigh--Ritz method that makes it possible to compute upper and lower bounds for eigenvalues.

In \S\,3.1, the problem statement is given: bases of the $p$-dimensional subspaces $\widetilde{U}$ and $A\widetilde{U}$ are assumed known, where the operator $A = L^{-1}M$ is self-adjoint in $H_L$, and it is required to compute two-sided bounds for the eigenvalues $\lambda_1, \ldots, \lambda_p$. For $p = 1$, a solution to this problem is given by the Temple method.

In \S\,3.2, we formulate the Temple--Lehmann method---a block analogue of the Temple method. Properties of the Temple--Lehmann method analogous to those of the Rayleigh--Ritz method are described: the equivalence of the operator and matrix formulations and the existence of Courant--Fischer-type variational principles. Conditions under which lower and upper bounds are guaranteed are specified. The Kato--Temple and Kato methods are formulated, and it is proved that they yield worse eigenvalue bounds than the Temple--Lehmann method (the comparison with the Kato--Temple method is carried out only in a special case).

Error estimates for the Temple--Lehmann method are presented and proved in \S\,3.3. They show that the eigenvalues are approximated to the same order as in the Rayleigh--Ritz method, while the lower approximations are more accurate than the Ritz values.

In \S\,3.4, we study a modification of the Temple--Lehmann method generated by the new choice of the operator $A = -M^{-1}L$ in the inner product of the space $H_M$.

% === Orig. p. 99 [s51R] ===
\monosection{3.1}{Problem statement and bounds for a single eigenvalue}

Suppose that the operator $A$ is self-adjoint in $H_L$ and shares the spectrum and eigenvectors of problem (0.1), for instance, $A = L^{-1}M$. In the iterative methods with operator $A$ from Chapter~1 ($p = 1$) and Chapter~2 ($p > 1$), at each iteration, for a given $p$-dimensional subspace $\widetilde{U}$, the $p$-dimensional subspace $A\widetilde{U}$ is computed. How can one make optimal use of this information to compute two-sided eigenvalue bounds?

A similar problem for $p = 1$ was posed and solved in \citep{Kato1949} by means of the following Temple-type functional:
\[
\lambda^\alpha(u) = (Au, [A - \alpha E]u)_L / (u, [A - \alpha E]u)_L.
\tag{3.1.1}
\]

\textbf{Lemma 3.1.1.} If $\alpha \leqslant \lambda_{\min}$, then $\lambda_1 \geqslant \lambda^\alpha(u)$; if $\lambda(u) > \alpha \geqslant \lambda_2$, then $\lambda^\alpha(u) \geqslant \lambda_1$.

It is easy to verify that the following representation holds:
\[
\lambda^\alpha(u) = \lambda(u) + \frac{\|Au - \lambda(u)u\|_L^2 / \|u\|_L^2}{\lambda(u) - \alpha}.
\tag{3.1.2}
\]
Then the assertions of the lemma are equivalent to the inequalities
\[
\begin{aligned}
[\lambda_1 - \lambda(u)][\lambda(u) - \lambda_2]
&\leqslant \|Au - \lambda(u)u\|_L^2 / \|u\|_L^2 \\
&\leqslant [\lambda_1 - \lambda(u)][\lambda(u) - \lambda_{\min}],
\end{aligned}
\tag{3.1.3}
\]
of which the first holds for $\lambda(u) > \lambda_2$.

The inequalities of Lemma~3.1.1 and formula (3.1.3) are known as Temple's inequalities \citep{Temple1928} and are widely used to compute two-sided eigenvalue bounds; see, e.g., \citep{Parlett1983,Nikolskii1976,Rektorys1985ru,ReedSimon1982ru,Chatelin1983}. The corresponding methods for computing eigenvalue bounds are called Temple methods.

% === Orig. p. 100 [s52L] ===
The bounds for $\lambda_1$ in Lemma~3.1.1 are optimal because the lower bound is exact for all vectors $u \in \myspan\{u_1, u_{\min}\}$, whereas the upper bound is exact for $u \in \myspan\{u_1, u_2\}$ \citep{Kato1949}.

Inequalities (3.1.3) make it easy to estimate the approximation error for $\lambda(u) > \lambda_2$ and $\lambda(u) > \alpha$:
\[
|\lambda_1 - \lambda^\alpha(u)| \leqslant \frac{\max\{|\lambda_2 - \alpha|, |\alpha - \lambda_{\min}|\}}{\lambda(u) - \alpha} [\lambda_1 - \lambda(u)].
\tag{3.1.4}
\]

Minimizing the right-hand side of (3.1.4) with respect to $\alpha$ leads to $\alpha = (\lambda_2 + \lambda_{\min}) / 2$. For such $\alpha$, therefore, the functional $\lambda^\alpha(u)$ approximates $\lambda_1$ with the best error estimate.

The best upper bound for $\lambda_1$ is obtained for $\alpha = \lambda_2$, and the best lower bound for $\alpha = \lambda_{\min}$.

Let us consider the question of evaluating the functional $\lambda^\alpha(u)$ when the vector $L^{-1}Mu$ is unknown. For $A = L^{-1}M$, we write the functional $\lambda^\alpha(\cdot)$ from (3.1.1) in the form
\[
\lambda^\alpha(u) = \frac{(L^{-1}Mu, Mu) - \alpha(Mu, u)}{(Mu, u) - \alpha(Lu, u)},
\tag{3.1.5}
\]
and representation (3.1.2) in the form
\[
\lambda^\alpha(u) = \frac{(Mu, u)}{(Lu, u)} + \frac{(L^{-1}[M - \lambda(u)L]u, [M - \lambda(u)L]u)}{(Mu, u) - \alpha(Lu, u)}.
\tag{3.1.6}
\]

It can be seen from formulas (3.1.5) and (3.1.6) that one cannot avoid solving systems of equations of the form $Lv = f$, whatever form is used to write the functional $\lambda^\alpha(\cdot)$ with $A = L^{-1}M$ in the inner product of the space $H_L$.

% === Orig. p. 101 [s52R] ===
\monosection{3.2}{Formulation of the Temple--Lehmann method and its properties}

For a subspace $\widetilde{U}$ ($\dim \widetilde{U} = p$), we determine the smallest Ritz value $\tilde{\lambda}_p$, fix a parameter $\alpha < \tilde{\lambda}_p$, and introduce on $\widetilde{U}$ the functional $\lambda^\alpha(u)$ from (3.1.1). The functional $\lambda^\alpha(\cdot)$ has exactly $p$ stationary values
\[
\lambda_1^\alpha \geqslant \ldots \geqslant \lambda_p^\alpha,
\]
which are the eigenvalues of the problem
\[
\widetilde{P}_L A[A - \alpha E]\big|_{\widetilde{U}} u^\alpha = \lambda^\alpha \widetilde{P}_L [A - \alpha E]\big|_{\widetilde{U}} u^\alpha
\tag{3.2.1}
\]
or its matrix analogue:
\[
(\widehat{AA} - \alpha\widehat{A})X = \lambda^\alpha(\widehat{A} - \alpha\widehat{E})X, \quad X = (x_1, \ldots, x_p)^T.
\tag{3.2.2}
\]

Here, as usual, $\widetilde{P}_L$ is the orthogonal projector in $H_L$ onto $\widetilde{U}$, and the symmetric matrices $\widehat{AA}$, $\widehat{A}$, and $\widehat{E}$ are defined entrywise by
\[
(A\hat{u}_k, A\hat{u}_n)_L, \quad (A\hat{u}_k, \hat{u}_n)_L \quad \text{and} \quad (\hat{u}_k, \hat{u}_n)_L, \quad k, n = 1, \ldots, p,
\]
with respect to some basis $\hat{u}_1, \ldots, \hat{u}_p \in \widetilde{U}$.

There is also a variational characterization of the values $\lambda_i^\alpha$:
\[
\begin{aligned}
\lambda_i^\alpha &= \max_{\widetilde{V} \subseteq \widetilde{U},\, \dim \widetilde{V} = i} \;\, \min_{u \in \widetilde{V} \setminus 0} \lambda^\alpha(u) \\
&= \min_{\widetilde{W} \subseteq \widetilde{U},\, \dim \widetilde{W} = p - i + 1} \;\, \max_{u \in \widetilde{W} \setminus 0} \lambda^\alpha(u).
\end{aligned}
\tag{3.2.3}
\]

The proofs of all these assertions are analogous to the corresponding ones for the Rayleigh--Ritz method, since for $u \in \widetilde{U}$ the functional $\lambda^\alpha(u)$ is actually the Rayleigh quotient for the operator $A$ in the inner product $(L[A - \alpha E]\cdot, *)$ on the subspace $\widetilde{U}$. Note that for $\alpha > \lambda_{\min}$ this inner product on the whole space $H$ is
% === Orig. p. 102 [s53L] ===
indefinite. In this case it is possible that the values $\lambda_i^\alpha$ bound the eigenvalues $\lambda_i$ from above, in contrast to the ordinary Ritz values $\tilde{\lambda}_i \leqslant \lambda_i$. The justification of such a two-sided bounding property is due to Lehmann; see, e.g., \citep{Lehmann1949a,GoerischAlbrecht1984a}.

\textbf{Theorem 3.2.1.} If $\alpha \leqslant \lambda_{\min}$, then $\lambda_i^\alpha \leqslant \lambda_i$; if $\alpha \in [\lambda_{p+1}, \tilde{\lambda}_p)$, then $\lambda_i \leqslant \lambda_i^\alpha$ ($i = 1, \ldots, p$).

Theorem~3.2.1 is the block analogue of Lemma~3.1.1. Thus, Lehmann's method (3.2.2) is the block analogue of the Temple method % [orig: misspelled "Lemmana" in Russian]
from \S\,3.1. The name ``Temple--Lehmann method'' is used to distinguish this method from many others proposed by Lehmann.

A proof of this theorem can also be found in \citep[pp.~226--227]{Parlett1983}, where a matrix formulation of problem (3.2.1) different from (3.2.2) is given. It should be pointed out that in practical implementation the use of that scheme (which is also repeated in \citep[pp.~145--146]{VoevodinKuznetsov1984}) requires greater computational cost than using (3.2.2).

An important property of the values $\lambda_i^\alpha$ is their monotonicity with respect to $\alpha$:

\textbf{Lemma 3.2.1.} The quantities $\lambda_i^\alpha$ are nondecreasing with respect to $\alpha < \tilde{\lambda}_p$. If the subspace $\widetilde{U}$ contains no eigenvectors, then the quantities $\lambda_i^\alpha$ are strictly increasing with respect to $\alpha < \tilde{\lambda}_p$.

This assertion follows directly from the monotonicity of the functional $\lambda^\alpha(\cdot)$ with respect to $\alpha < \tilde{\lambda}_p$ and equalities (3.2.3).

Note that as $\alpha \to -\infty$ the values $\lambda_i^\alpha$ reduce to the Ritz values $\tilde{\lambda}_i$. It then follows from Lemma~3.2.1 that for $\alpha = \lambda_{\min}$ the quantities $\lambda_i^\alpha$ provide more accurate lower bounds for the eigenvalues $\lambda_i$ than $\tilde{\lambda}_i$.

The question naturally arises: are the eigenvectors $u_i^\alpha$ of problem (3.2.1) substantially better approximations to $u_i$ than the Ritz vectors $\tilde{u}_i$? Asymptotically (as $\widetilde{U} \to U$), the answer is negative---these vectors asymptotically coincide.

\textbf{Lemma 3.2.2.} Let
\[
\varepsilon_i \equiv \ang(u_i, \widetilde{U}), \quad i = 1, \ldots, p, \quad \varepsilon = \max \varepsilon_i \to 0, \quad \varepsilon^3 = O(\varepsilon_i).
\]
Then
\[
\ang(u_i^\alpha, \tilde{u}_i) = O(\varepsilon\varepsilon_i), \quad \ang(u_i^\alpha; u_i) \sim \ang(\tilde{u}_i; u_i) \sim \varepsilon_i.
\]

% === Orig. p. 103 [s53R] ===
Lemma~3.2.2 is a special case of Theorem~2.2.5: in (2.2.14) one should set
\[
G = A - \alpha E, \quad F = AG,
\]
cf.\ (2.2.14) and (3.2.1). The justification was given above, at the end of \S\,2.2. Thus, Lemma~3.2.2 establishes the asymptotic equivalence of the vectors $u_i^\alpha$ and $\tilde{u}_i$:
\[
\ang(u_i^\alpha; u_i) \sim \ang(\tilde{u}_i, u_i).
\tag{3.2.4}
\]

At the same time, for the eigenvalue error we have equivalence only in order:
\[
|\lambda_i - \lambda_i^\alpha| \asymp \lambda_i - \tilde{\lambda}_i,
\]
or, in a more precise formulation,
\[
|\lambda_i - \lambda_i^\alpha| \sim c^\alpha (\lambda_i - \tilde{\lambda}_i),
\tag{3.2.5}
\]
where $c^\alpha < 1$, for example, for $\alpha \in (-\infty; \lambda_{\min})$. As an illustration of relation (3.2.5) for $p = 1$, we point out inequality (3.1.4).

An alternative to the Temple--Lehmann method is the Kato--Temple method, which makes it possible, given good approximations to the eigenvectors $u_i$, to obtain two-sided bounds for each $\lambda_i$ independently. It is based on the following fact generalizing the first inequality in (3.1.2) \citep{Kato1949,Chatelin1983,Parlett1983}: if $\lambda(u) \in [\lambda_{i+1}, \lambda_i]$, then % [math note: in the original, (3.1.2) is cited instead of (3.1.3)]
\[
[\lambda_i - \lambda(u)][\lambda(u) - \lambda_{i+1}] \leqslant \|Au - \lambda(u)u\|_L^2 / \|u\|_L^2.
\tag{3.2.6}
\]
For example, if a number $\alpha \in [\lambda_{i+1}, \lambda(u))$ is known, then from (3.2.6) we deduce the bound
% === Orig. p. 104 [s54L] ===
\[
\lambda_i \leqslant \lambda^\alpha(u).
\tag{3.2.7}
\]

This method of deriving eigenvalue bounds, proposed in \citep{Kato1949} and generalizing the Temple method of \S\,3.1, is called the Kato--Temple method.

When the subspace $\widetilde{U}$ is close to $U$, the Ritz vectors $\tilde{u}_i$ can be taken as suitable approximations to the eigenvectors $u_i$.

Let us compare the accuracy of the Kato--Temple and Temple--Lehmann methods in the simplest special case $i = 1$, $p = 2$.

\textbf{Lemma 3.2.3.} Let $p = 2$ and $\widetilde{U} \to U$. To bound the eigenvalue $\lambda_1$ from above, we use the quantities $\lambda^{\lambda_2}(\tilde{u}_1)$ (the Kato--Temple method) and $\lambda_1^{\lambda_3}$ (the Temple--Lehmann method), and to bound it from below, the Ritz value $\tilde{\lambda}_1$. Then for the error ratio we have
\[
\frac{\lambda^{\lambda_2}(\tilde{u}_1) - \tilde{\lambda}_1}{\lambda_1^{\lambda_3} - \tilde{\lambda}_1} \longrightarrow \frac{\lambda_1 - \lambda_3}{\lambda_1 - \lambda_2} > 1.
\tag{3.2.8}
\]

The asymptotic relation (3.2.8) shows that for a sufficiently ``accurate'' subspace $\widetilde{U}$ the Temple--Lehmann method yields better eigenvalue bounds than the Kato--Temple method, especially when $\lambda_1$ is close to $\lambda_p$.

In \citep{Kato1949}, another algorithm for computing two-sided eigenvalue bounds was also proposed, close to the Temple--Lehmann method in its range of applicability and computational cost. Let us formulate this method as applied to upper bounds $\rho_i \geqslant \lambda_i$, using the notation of the present work and the quantities
\[
\varepsilon_i \equiv \|A\tilde{u}_i - \tilde{\lambda}_i \tilde{u}_i\|_L,
\]
where $\tilde{\lambda}_i, \tilde{u}_i$ are the Ritz values and Ritz vectors. The number $\rho_i$ is defined as the largest root of the equation \citep{Kato1949}:
\[
\sum_{j=i}^p \varepsilon_j^2 (\rho - \tilde{\lambda}_j)^{-1} (\tilde{\lambda}_j - \alpha)^{-1} = 1
\tag{3.2.9}
\]
% === Orig. p. 105 [s54R] ===
under the condition $\alpha \in [\lambda_{p+1}, \tilde{\lambda}_p)$. Kato also gives a simpler bound for $\rho_i$:
\[
\rho_i \leqslant \tilde{\lambda}_i + \sum_{j=p}^i \varepsilon_j^2 (\tilde{\lambda}_j - \alpha)^{-1}, % [math note: in the original, the summation limits are written with j=p at the bottom and i at the top (cf. j=i..p in (3.2.9))]
\tag{3.2.10}
\]
which was also included in the book \citep{Chatelin1983}.

\textbf{Lemma 3.2.4.} If $\alpha \in [\lambda_{p+1}, \tilde{\lambda}_p)$, then $\lambda_i^\alpha \leqslant \rho_i$, $i = 1, \ldots, p$.

Thus, the conditions of applicability of the Temple--Lehmann and Kato methods coincide, their computational costs are roughly the same, and in accuracy Kato's method is inferior. The fact that $\lambda_i^\alpha$ provide more accurate bounds can also be proved in the most general case (for unbounded operators and when $\lambda_1, \ldots, \lambda_p$ are not necessarily extreme eigenvalues) considered in \citep{Kato1949}. A justification of the Temple--Lehmann method for this case can be found in \citep{GoerischAlbrecht1984a}.

\textit{Proof of Lemma 3.2.3.} Denote the left-hand side of (3.2.8) by $\varkappa$ and write
\[
\varkappa = [\lambda^{\lambda_2}(\tilde{u}_1) - \lambda(\tilde{u}_1)] / [\lambda^{\lambda_3}(u_1^\alpha) - \lambda(\tilde{u}_1)].
\tag{3.2.11}
\]
Assume that $\|\tilde{u}_1\|_L = \|u_1^\alpha\|_L = 1$. Next, from Lemma~3.2.2 and % [orig: period after 1 omitted]
the asymptotic equality (3.2.4) we have
\[
\tilde{u}_1 = u_1 + O(\varepsilon), \quad u_1^\alpha = \tilde{u}_1 + O(\varepsilon^2).
\]
We use representation (3.1.2):
\[
\lambda^\alpha(u) = \lambda(u) + \frac{\|Au - \lambda(u)u\|_L^2 / \|u\|_L^2}{\lambda(u) - \alpha}.
\]
Then in expression (3.2.11) for $\varkappa$ the numerator is a quantity of order $O(\varepsilon^2)$. The denominator in (3.2.11) takes the form
\[
\lambda(u_1^\alpha) - \lambda(\tilde{u}_1) + \|Au_1^\alpha - \lambda(u_1^\alpha)u_1^\alpha\|_L^2 \|u_1^\alpha\|_L^{-2} [\lambda(u_1^\alpha) - \lambda_3]^{-1}
\]
% === Orig. p. 106 [s55L] ===
and is likewise of second order in $\varepsilon$. Here the difference of the Rayleigh quotients is of order $O(\varepsilon^3)$, since
\begin{gather*}
|\lambda(u_1^\alpha) - \lambda(\tilde{u}_1)| = O(\|u_1^\alpha - \tilde{u}_1\|_L \times \|\operatorname{grad} \lambda(\tilde{u}_1)\|_L), \\
\|\operatorname{grad} \lambda(\tilde{u}_1)\|_L \asymp \|A\tilde{u}_1 - \lambda(\tilde{u}_1)\tilde{u}_1\|_L = O(\varepsilon).
\end{gather*}

Then, up to $O(\varepsilon)$, $\varkappa$ is determined by the quantity
\[
\begin{aligned}
&\frac{\lambda(u_1^\alpha) - \lambda_3}{\lambda(\tilde{u}_1) - \lambda_2} \times \frac{\|A\tilde{u}_1 - \lambda(\tilde{u}_1)\tilde{u}_1\|_L^2}{\|Au_1^\alpha - \lambda(u_1^\alpha)u_1^\alpha\|_L^2} \times \frac{\|u_1^\alpha\|_L^2}{\|\tilde{u}_1\|_L^2} \\
&\qquad \to (\lambda_1 - \lambda_3) / (\lambda_1 - \lambda_2),
\end{aligned}
\]
which proves Lemma~3.2.3. The last asymptotic relation holds by virtue of
\[
\lambda(\tilde{u}_1) = \lambda_1 + O(\varepsilon^2), \quad \lambda(u_1^\alpha) = \lambda_1 + O(\varepsilon^2),
\]
and the simple fact that in its upper part the last factors tend to unity since $u_1^\alpha = \tilde{u}_1 + O(\varepsilon^2)$. Lemma~3.2.3 is proved.

\textit{Proof of Lemma 3.2.4.} Let us reconstruct the line of reasoning from \citep{Kato1949}. By dimension arguments, there exists a vector
\[
w \in \myspan^{\perp}\{u_{i+1}, \ldots, u_p\} \cap \myspan\{\tilde{u}_i, \ldots, \tilde{u}_p\}.
\]
By definition, $\lambda(w) \in [\tilde{\lambda}_p, \tilde{\lambda}_i]$, and one can apply the Kato--Temple method (3.2.7):
\[
\lambda_i \leqslant \lambda^\alpha(w), \quad \alpha \in [\lambda_{p+1}, \tilde{\lambda}_p).
\]
However, the vector $w$ is actually unknown; therefore, Kato represents it in the form
% === Orig. p. 107 [s55R] ===
\[
w = \sum_{j=i}^p \gamma_j \tilde{u}_j,
\]
and, in order to eliminate the unknowns $\gamma_j$, establishes the validity of the inequality $\lambda^\alpha(w) \leqslant \rho_i$, with $\rho_i$ from (3.2.9), for arbitrary coefficients $\gamma_i$. Consequently, % [math note: \gamma_i instead of \gamma_j in the original]
\[
\rho_i \geqslant \max_{u \in \myspan\{\tilde{u}_i, \ldots, \tilde{u}_p\}} \lambda^\alpha(u).
\]
Now observe, on the basis of the minimax variational principle in (3.2.3), that on the right-hand side of this inequality the expression is no smaller than $\lambda_i^\alpha$. Thus, $\rho_i \geqslant \lambda_i^\alpha$. This completes the proof of Lemma~3.2.4.

At the end of this section, let us consider the special case of computing lower bounds when $A \geqslant 0$, in which the preceding results become almost obvious.

Choose $\alpha = 0$; then the functional $\lambda^\alpha(\cdot)$ reduces to
\[
\lambda^0(u) = (Au, Au)_L / (u, Au)_L
\tag{3.2.12}
\]
---a Schwarz-type quotient \citep[p.~179]{Collatz1968ru}. From the Cauchy--Bunyakovsky--Schwarz inequality we deduce the inequalities
\[
\lambda^0(u) \geqslant \sqrt{(Au, Au)_L / (u, u)_L} \geqslant \lambda(u).
\]

Thus, among the three functionals written out here, it is $\lambda^0(\cdot)$ that provides the best lower approximation to the eigenvalues, and it should be recommended for use in subspace iterations when $A > 0$ (cf.\ the choice in the book \citep[pp.~311--321]{Parlett1983}).

Of course, the functional $\lambda^\alpha(\cdot)$ can be used in place of the Rayleigh quotient not only in numerical methods, but also in the discretization of continuous problems, for instance, for spectral problems with integral operators \citep{Fichera1978,Chatelin1983}. Such an application can be interpreted as an example of Krylov's method of moments \citep[p.~87]{LuchkaLuchka1985}.

% === Orig. p. 108 [s56L] ===
\monosection{3.3}{Error estimates for the Temple--Lehmann method}

We formulate the error estimates for the Temple--Lehmann method of \S\,3.2 in the same terms as the error estimates for the Rayleigh--Ritz method in Theorem~2.3.2. % [math note: in the original, Theorem 2.3.2 is cited instead of Theorem 2.2.3]

\textbf{Theorem 3.3.1.} Let the numbers $\lambda_1^\alpha \geqslant \ldots \geqslant \lambda_p^\alpha$ be the eigenvalues of problem (3.2.1) or (3.2.2) of the Temple--Lehmann method. Then, if $\alpha \leqslant \lambda_{\min}$,
\[
0 \leqslant \frac{\lambda_i - \lambda_i^\alpha}{\lambda_i^\alpha - \lambda_{\min}} \leqslant \frac{\lambda_{p+1} - \alpha}{\lambda_i - \alpha} \tan^2 \ang_L(\widetilde{U}; U);
\tag{3.3.1}
\]
if $\alpha \in [\lambda_{p+1}, \tilde{\lambda}_p)$, then
\[
0 \leqslant \frac{\lambda_i^\alpha - \lambda_i}{\lambda_i^\alpha - \lambda_{\min}} \leqslant \frac{\alpha - \lambda_{\min}}{\lambda_p - \alpha} \tan^2 \ang_L(\widetilde{U}; U).
\tag{3.3.2}
\]

As $\alpha \to -\infty$, inequalities (3.3.1) reduce to (2.2.10). The quantity $(\lambda_{p+1} - \alpha) / (\lambda_i - \alpha)$ indicates the approximate gain in accuracy when approximating the eigenvalues $\lambda_i$ by $\lambda_i^\alpha$ as compared with the Ritz values $\tilde{\lambda}_i$ (see Theorem~2.2.3). Recall that the presence of such a gain is guaranteed by the monotonicity Lemma~3.2.1.

\textit{Proof of Theorem 3.3.1.} To verify inequality (3.3.1), it suffices to observe that $\lambda_i^\alpha$ are the Ritz values for the operator $A$ in $H_L$ and the subspace $\sqrt{A - \alpha E}\,\widetilde{U}$. The condition $\alpha \leqslant \lambda_{\min}$ ensures the existence and self-adjointness in $H_L$ of the operator $\sqrt{A - \alpha E}$. Setting in Theorem~2.3.1
\[
U^0 = \widetilde{U}, \quad U^k = \sqrt{A - \alpha E}\,\widetilde{U}, \quad F_k = \sqrt{A - \alpha E}, \quad \bar{\sigma}_i^k = \frac{\lambda_{p+1} - \alpha}{\lambda_i - \alpha},
\]
we see that inequality (2.3.5) leads to (3.3.1).

We were unable to justify inequality (3.3.2) in a similar manner, so we carry out a direct verification, assuming, as usual, $L = E$.

% === Orig. p. 109 [s56R] ===
Fix $i \in [1, p]$ and introduce the subspaces
\[
\overline{W}_i \equiv \myspan\{\bar{u}_i, \ldots, \bar{u}_p\}, \quad W_i \equiv \myspan\{u_i, \ldots, u_p\},
\]
where $\bar{u}_j$, $j = i, \ldots, p$, are the vectors from (2.2.11) in the subspace $\widetilde{U}$ (see \S\,2.2). In the nontrivial case where $\tan^2 \ang(\widetilde{U}, U) < \infty$, the following assertions hold:
\[
\dim \overline{W}_i = \dim W_i, \quad \tan^2 \ang(\overline{W}_i, W_i) \leqslant \tan^2 \ang(\widetilde{U}; U).
\]

From the minimax principle (3.2.3) it follows that
\[
\lambda_i^\alpha \leqslant \bar{\lambda}_i^\alpha \equiv \max_{w \in \overline{W}_i \setminus 0} \lambda^\alpha(\bar{w}). % [math note: in the original, w under max has no bar, whereas the argument of \lambda^\alpha(\bar{w}) has a bar]
\]

Choose a vector $\bar{w}_i \in \overline{W}_i \setminus 0$ such that $\lambda^\alpha(\bar{w}_i) = \bar{\lambda}_i^\alpha$, and define the vector $w_i \equiv P\bar{w}_i$, where $P$ is the orthogonal projector onto $U$. We have $w_i \in W_i \setminus 0$. If $\bar{w}_i = w_i$, then clearly $\lambda_i^\alpha = \lambda_i$; % [orig: period after U omitted]
henceforth we therefore assume $\bar{w}_i \neq w_i$. Then $v_i \equiv \bar{w}_i - w_i \in U^\perp \setminus 0$ and
\[
\bar{\lambda}_i^\alpha = \frac{(A[A - \alpha E]w_i, w_i) + (A[A - \alpha E]v_i, v_i)}{([A - \alpha E]w_i, w_i) + ([A - \alpha E]v_i, v_i)}.
\]

Taking into account that $\alpha \in [\lambda_{p+1}, \lambda_p)$, we make the substitution
\[
\hat{w}_i \equiv \sqrt{[A - \alpha E]|_U}\, w_i, \quad \hat{v}_i \equiv \sqrt{[\alpha E - A]|_{U^\perp}}\, v_i.
\tag{3.3.3}
\]

We obtain the representation
\[
\bar{\lambda}_i^\alpha = \frac{(A\hat{w}_i, \hat{w}_i) - (A\hat{v}_i, \hat{v}_i)}{(\hat{w}_i, \hat{w}_i) - (\hat{v}_i, \hat{v}_i)},
\tag{3.3.4}
\]
% === Orig. p. 110 [s57L] ===
from which the basic identity follows:
\[
\frac{\bar{\lambda}_i^\alpha - \lambda(\hat{w}_i)}{\bar{\lambda}_i^\alpha - \lambda(\hat{v}_i)} = \frac{\|\hat{v}_i\|^2}{\|\hat{w}_i\|^2}.
\tag{3.3.5}
\]
In the left-hand side of (3.3.5) we replace the quantities as follows: $\bar{\lambda}_i^\alpha$ by $\lambda_i^\alpha \leqslant \bar{\lambda}_i^\alpha$, $\lambda(\hat{w}_i)$ by $\lambda_i \geqslant \lambda(\hat{w}_i)$, and $\lambda(\hat{v}_i)$ by $\lambda_{\min} \leqslant \lambda(\hat{v}_i)$. The first inequality was already stated above, the second follows from $\hat{w}_i \in W_i \setminus 0$, and the third is obvious. Under this replacement, the left-hand side of (3.3.5) does not increase, since
\[
\lambda(\hat{w}_i) \geqslant \lambda_p > \lambda_{p+1} \geqslant \lambda(\hat{v}_i).
\]
Thus, we arrive at the inequality
\[
(\lambda_i^\alpha - \lambda_i) / (\lambda_i^\alpha - \lambda_{\min}) \leqslant \|\hat{v}_i\|^2 / \|\hat{w}_i\|^2.
\]

Now let us show that the right-hand side of (3.3.5) does not exceed the right-hand side of inequality (3.3.2), which will justify the latter.

Recalling definition (3.3.3), we write the inequalities
\[
\frac{\|\hat{v}_i\|^2}{\|\hat{w}_i\|^2} \leqslant \frac{\|(\alpha E - A)|_{U^\perp}\|}{\|(A - \alpha E)|_U^{-1}\|^{-1}} \cdot \frac{\|v_i\|^2}{\|w_i\|^2} \leqslant \frac{\alpha - \lambda_{\min}}{\lambda_p - \alpha} \cdot \frac{\|v_i\|^2}{\|w_i\|^2}. % [math note: in the original, one vertical bar is missing on the left of the norm of \hat{w}_i: |\hat{w}_i\|^2]
\]
Finally, we have
\[
\frac{\|v_i\|^2}{\|w_i\|^2} = \tan^2 \ang(\bar{w}_i; W_i) \leqslant \tan^2 \ang(\overline{W}_i; W_i) \leqslant \tan^2 \ang(\widetilde{U}; U).
\]
Theorem~3.3.1 is proved.

% === Orig. p. 111 [s57R] ===
\monosection{3.4}{Another way of reducing the generalized eigenvalue problem to a standard one}

Above, in \S\S\,3.1--3.3, we considered methods that use the operator $A = L^{-1}M$ and the inner product of the space $H_L$. As was already noted in \S\,1.6 and \S\,2.8, another choice is also possible, provided only that the operator $A$ is self-adjoint in some inner product. In the present section we examine how the Temple--Lehmann method of \S\,3.2 changes, and what the error estimates of \S\,3.3 become, for the following choice of the operator $A$ when $M > 0$:
\[
A = -M^{-1}L, \quad A = A^* \quad \text{in} \quad H_M.
\tag{3.4.1}
\]

In all cases, the bounds $\lambda^\alpha$ are sought as the eigenvalues of problems of the form
\[
\widehat{F}X = \lambda^\alpha \widehat{G}X
\tag{3.4.2}
\]
with $p \times p$ matrices $\widehat{F}$ and $\widehat{G}$ defined as Gram matrices, in various inner products, of the vectors $\hat{u}_1, \ldots, \hat{u}_p$ forming a basis of the subspace $\widetilde{U}$. For instance, in the case $A = L^{-1}M$ in $H_L$ we have (cf.\ (3.2.2)):
\begin{gather*}
\widehat{F} = \{(L^{-1}M\hat{u}_i, [L^{-1}M - \alpha E]\hat{u}_j)_L\}, \\
\widehat{G} = \{(\hat{u}_i, [L^{-1}M - \alpha E]\hat{u}_j)_L\}.
\end{gather*}

For the choice (3.4.1), instead of (3.1.1) we obtain a new formula for the functional $\lambda^\alpha(u)$:
\[
\begin{aligned}
\lambda^\alpha(u) &= (u, [E - \alpha M^{-1}L]u)_M \\
&\quad \big/ (M^{-1}Lu, [E - \alpha M^{-1}L]u)_M,
\end{aligned}
\tag{3.4.3}
\]
% === Orig. p. 112 [s58L] ===
which leads to the new matrices $\widehat{F}$ and $\widehat{G}$:
\begin{gather*}
\widehat{F} = \{(\hat{u}_i, [E - \alpha M^{-1}L]\hat{u}_j)_M\}, \\
\widehat{G} = \{(M^{-1}L\hat{u}_i, [E - \alpha M^{-1}L]\hat{u}_j)_M\}.
\end{gather*}

An analogue of Theorem~3.3.1 is as follows.

\textbf{Theorem 3.4.1.} If $\alpha \in [0, \lambda_{\min}]$, then
\[
0 \leqslant \frac{\lambda_i - \lambda_i^\alpha}{\lambda_i^\alpha - \lambda_{\min}} \leqslant \frac{\lambda_i}{\lambda_{\min}} \cdot \frac{\lambda_{p+1} - \alpha}{\lambda_i - \alpha} \tan^2 \ang_L(\widetilde{U}; U);
\tag{3.4.4}
\]
if $\alpha \in [\lambda_{p+1}, \tilde{\lambda}_p)$, then
\[
0 \leqslant \frac{\lambda_i^\alpha - \lambda_i}{\lambda_i^\alpha - \lambda_{\min}} \leqslant \frac{\lambda_i}{\lambda_{\min}} \cdot \frac{\alpha - \lambda_{\min}}{\lambda_p - \alpha} \tan^2 \ang_L(\widetilde{U}; U).
\tag{3.4.5}
\]

Comparing the error estimates of Theorems~3.3.1 and 3.4.1 shows that the original version of the Temple--Lehmann method from \S\,3.2 appears to provide sharper eigenvalue bounds. In practice, however, the choice of modification is, as a rule, predetermined by the nature of the information actually available. For instance, when using the Temple--Lehmann method in subspace iterations, it is natural to make the choice of the operator $A$ in the method consistent with that in the iterative scheme. In particular, the version from \S\,3.2 is efficiently implemented in the subspace iterations of \S\S\,2.4--2.6.

Other modifications of the Temple--Lehmann method are also possible, likewise based on replacing the Rayleigh quotient by a functional of the form
\[
\lambda^Q(u) \equiv (Au, u)_{LQ} / (u, u)_{LQ},
\]
where the operator $Q$ is a polynomial or rational function of $A = L^{-1}M$.

Note that in the proof of Theorem~3.3.1 we used the functional $\lambda^Q(\cdot)$ with $Q(\lambda)$ being the Heaviside step function with a jump between $\lambda_{p+1}$ and $\lambda_p$; see (3.3.4).

% === Orig. p. 113 [s58R] ===
It should be recalled that as $\widetilde{U} \to U$, the stationary vectors of the functional $\lambda^Q(\cdot)$ on the subspace $\widetilde{U}$, defined by a formula of the type $F = AG$ in (2.2.14), asymptotically coincide with the Ritz vectors, as stated in Theorem~2.2.5. However, the stationary values $\lambda_i^Q$ may differ substantially for different $Q$ even asymptotically. Thus, using the functional $\lambda_i^Q$ asymptotically leads % [math note: in the original, "of the functional \lambda_i^Q" (with subscript i)]
to approximations to the eigenvectors that are not substantially new, while providing new bounds for the eigenvalues. In particular, if the metric of $H_{LQ}$ is indefinite, then in a number of cases upper bounds for $\lambda_i$ are computed.

\textit{Proof of Theorem 3.4.1.} To verify inequality (3.4.4), we apply the result (3.3.1) to the choice (3.4.1). After trivial transformations we obtain the inequalities
\[
0 \leqslant \frac{\lambda_i - \lambda_i^\alpha}{\lambda_i^\alpha - \lambda_{\min}} \leqslant \frac{\lambda_i^2}{\lambda_{\min}\lambda_{p+1}} \cdot \frac{\lambda_{p+1} - \alpha}{\lambda_i - \alpha} \tan^2 \ang_M(\widetilde{U}; U).
\tag{3.4.6}
\]
Recalling the proofs of Theorems~3.3.1 and 2.3.1, we observe that inequality (3.4.6) can be strengthened by replacing $\widetilde{U}$ with $\overline{V}_i$ and $U$ with $V_i$. From Lemma~2.8.1 it follows that
\[
\tan^2 \ang_M(\overline{V}_i; V_i) \leqslant \frac{\lambda_{p+1}}{\lambda_i} \tan^2 \ang_L(\overline{V}_i; V_i).
\]
Taking into account the inequality
\[
\tan^2 \ang_L(\overline{V}_i; V_i) \leqslant \tan^2 \ang_L(\widetilde{U}; U),
\]
we arrive at (3.4.4).

To derive (3.4.5), we employ a different approach. Observe that the functionals $\lambda^\alpha(\cdot)$ from (3.1.1) and $\lambda^\alpha(\cdot)$ from (3.4.3) are related by the identity
\[
\lambda_{(3.1.1)}^\alpha(v) = \lambda_{(3.4.3)}^\alpha(u), \quad v = \sqrt{M^{-1}L}\,u.
\]
Thus, the Temple--Lehmann method with functional (3.4.3) on the subspace $\widetilde{U}$ can be interpreted as the Temple--Lehmann method with functional (3.1.1) on the subspace $\sqrt{M^{-1}L}\,\widetilde{U}$. Now we use the strengthened version of bound (3.3.2):
% === Orig. p. 114 [s59L] ===
\[
0 \leqslant \frac{\lambda_i^\alpha - \lambda_i}{\lambda_i^\alpha - \lambda_{\min}} \leqslant \frac{\alpha - \lambda_{\min}}{\lambda_p - \alpha} \tan^2 \ang_L\bigl(\sqrt{M^{-1}L}\,\overline{W}_i; W_i\bigr)
\]
(see the proof of Theorem~3.3.1). Here we have used the fact that the vectors from (2.2.11) in the subspace $\sqrt{M^{-1}L}\,\widetilde{U}$ are obtained from the vectors of type (2.2.11) in the subspace $\widetilde{U}$ under the action of the operator $\sqrt{M^{-1}L}$. Recall that
\[
W_i = \myspan\{u_i, \ldots, u_p\}, \quad \overline{W}_i \subseteq \myspan^{\perp_L}\{u_1, \ldots, u_{i-1}\},
\]
and apply inequality (2.3.3) with $U^0 = \overline{W}_i$ and $F_k = \sqrt{M^{-1}L}$. Then
\[
\bar{\sigma}_p^k = \max_{j > p} |F_k(\lambda_j)|^2 \Big/ \min_{j \in [i, p]} |F_k(\lambda_j)|^2 = \frac{\lambda_i}{\lambda_{\min}}.
\]
Consequently,
\[
\tan^2 \ang_L\bigl(\sqrt{M^{-1}L}\,\overline{W}_i; W_i\bigr) \leqslant \frac{\lambda_i}{\lambda_{\min}} \tan^2 \ang_L(\overline{W}_i; W_i) \leqslant \frac{\lambda_i}{\lambda_{\min}} \tan^2 \ang_L(\widetilde{U}; U).
\]
Theorem~3.4.1 is proved.

At the end of this section, we note that when applying the methods considered above, one cannot avoid solving systems of equations of the form $(\alpha L + \beta M)v = f$. At the same time, it is of practical interest to construct a modification of the Temple--Lehmann method that would require solving systems of equations only with some auxiliary operator $B = B^* > 0 : 0 < \delta_0 B \leqslant L \leqslant \delta_1 B$ (see Chapter~4). For $p = 1$, from (3.1.5), (3.1.6), and Lemma~3.1.1 we obtain for $\lambda(u) > \lambda_2$ that
\begin{gather*}
\lambda_1 \leqslant \frac{\delta_0^{-1}(B^{-1}Mu, Mu) - \alpha(Mu, u)}{(Mu, u) - \alpha(Lu, u)},
\tag{3.4.7} \\
\lambda_1 \leqslant \frac{(Mu, u)}{(Lu, u)} + \delta_0^{-1} \frac{(B^{-1}[M - \lambda(u)L]u, [M - \lambda(u)L]u)}{(Mu, u) - \alpha(Lu, u)}.
\tag{3.4.8}
\end{gather*}
Relation (3.4.7), unlike (3.4.8), can be generalized to the case $p > 1$ \citep{GoerischAlbrecht1984a}; however, bound (3.4.7) is, generally speaking, not exact even for $u = u_1$.

% =============================================================================
% SOURCE MODULE: part_07_ch4_a.tex
% =============================================================================
% === Orig. p. 115 [s59R] ===
\monochapter{4}{Iterative Methods with an Auxiliary Operator}

We consider solving problem (0.1) by vector iterations with an auxiliary operator for $p = 1$ and by their block analogues for $p > 1$. The chapter consists of six sections.

In \S\,4.1, we define the concept of iterations with an auxiliary operator, introduce the two-stage form of the methods via an auxiliary spectral problem, describe the relationship between the eigenvalues of this problem and those of problem (0.1), and derive convergence rate estimates for the outer iterations.

In \S\,4.2, we study the convergence rate of an abstract two-stage vector iteration method with a finite number of inner iterations.

In \S\,4.3, we consider applications of the theory of \S\,4.2 to methods with Chebyshev parameters and Lanczos-type methods. For these methods, we derive convergence rate estimates that turn out to be asymptotically analogous to the corresponding estimates for optimal iterative methods for solving systems of linear equations.

In \S\,4.4, we investigate the simplest two-stage algorithms with a single inner iteration---the modified gradient methods. The general theory of \S\,4.2 is applied to them, and then specialized results are proved that yield sharper estimates in several important cases.

In \S\,4.5, we estimate the convergence rate of block analogues of the modified gradient methods from \S\,4.4. % [orig: period after section number in 1986 print]

In \S\,4.6, we propose a new way of choosing the metric when orthogonalizing against previously computed eigenvectors in iterative methods with an auxiliary operator. The presence of an auxiliary operator introduces certain specific features compared with the analysis in \S\,2.7.

% === Orig. p. 116 [s60L] ===
\monosection{4.1}{Two-stage form of the methods}

Suppose that we are given an auxiliary operator $B = B^* > 0$ for which systems of equations of the form $Bu = f$ can be solved efficiently, and the quantity $\delta$,
\[
\delta \equiv \delta_0 / \delta_1, \quad 0 < \delta_0 B \leqslant L \leqslant \delta_1 B, \tag{4.1.1}
\]
is as close to unity as possible. We use this operator to perform an equivalent transformation of problem (0.1) to the form
\[
B^{-1} M u = \lambda B^{-1} L u, \quad B^{-1} M = (B^{-1} M)^*, \quad B^{-1} L = (B^{-1} L)^* > 0, \tag{4.1.2}
\]
where taking the adjoint and the operator inequality are now understood in the space $H_B$.

Thus, the equivalent transformation of problem (0.1) into problem (4.1.2) does not take us outside the class of symmetric eigenvalue problems. (Here and throughout, it is assumed that passing from (0.1) to (4.1.2) is always accompanied by simultaneously replacing $H$ with $H_B$.) Note also that the values of the Rayleigh quotient $\lambda(\cdot)$ and of the functionals $\tan^2 \varphi(\cdot)$, $r^2(\cdot)$, and $\Delta\lambda(\cdot)$ from \S\,1.1 remain unchanged; for instance, % [orig: extra space before comma in 1986 print]
\[
\lambda(\cdot) = \frac{(M\cdot, \cdot)}{(L\cdot, \cdot)} = \frac{(B^{-1} M\cdot, \cdot)_B}{(B^{-1} L\cdot, \cdot)_B}.
\]

What do we gain by passing from (0.1) to (4.1.2) when applying iterative methods?

Consider two examples: a gradient-ascent-type method with ``inversion'' of the operator $L$,
\[
u^{n+1} = u^n + \alpha^n \{ L^{-1} M u^n - \lambda(u^n) u^n \} \tag{4.1.3}
\]
(cf.\ (1.2.1)), and without ``inversion,''
\[
u^{n+1} = u^n + \alpha^n \{ M u^n - \lambda(u^n) L u^n \}. \tag{4.1.4}
\]
As is well known, the name ``gradient ascent'' stems from the fact that the vectors in braces in (4.1.3) and (4.1.4) are codirectional with the gradients of the Rayleigh quotient at the point $u^n$ in the spaces $H_L$ and $H$, respectively.

% === Orig. p. 117 [s60R] ===
Let us replace (0.1) with (4.1.2) (and $H$ with $H_B$!) and assume that the quantities $\alpha^n$ in (4.1.3) and (4.1.4) remain unchanged. Evidently, algorithm (4.1.3) retains its original form, whereas (4.1.4) transforms into
\[
u^{n+1} = u^n + \alpha^n B^{-1} \{ M u^n - \lambda(u^n) L u^n \}. \tag{4.1.5}
\]
Note that method (4.1.5), as applied to problem (0.1), is known as modified gradient ascent, since the vector $u^{n+1} - u^n$ in (4.1.5) is codirectional with the gradient of the Rayleigh quotient at the point $u^n$ in the space $H_B$. The modification consists precisely in using the auxiliary inner product of the space $H_B$. Method (4.1.5) is studied in detail below in \S\,4.4.

Thus, we have established that method (4.1.3) is invariant under the equivalent transformation of problem (0.1), in contrast to method (4.1.4), which generates method (4.1.5). Related to these properties of methods (4.1.3) and (4.1.4) is the important fact that the eigenvectors of the transition operator from $u^n$ to $u^{n+1}$ in method (4.1.3) coincide with the corresponding eigenvectors of problem (0.1), whereas in methods (4.1.4) and (4.1.5), generally speaking, they do not.

From a computational standpoint, the difference between methods (4.1.3) and (4.1.5) is that in (4.1.5) one can, on the one hand, avoid the costly solution of systems of equations of the form $Lu = f$ and, on the other hand, by choosing the operator $B$ ``close'' to $L$, achieve convergence just as rapid as in the iterations (4.1.4) (see \S\,1.4). The quantity $\delta$ from (4.1.1) serves as a measure of closeness between the operators $B$ and $L$, and the operator $B$ may be interpreted as a ``model'' of the operator $L$.

The use of an auxiliary operator $B$ \citep{Kantorovich1948} for spectral problems (at the differential level) was proposed in \citep{Samokish1958}, where a method of type (4.1.5) was written down and its asymptotic convergence rate was estimated. Part of these results was generalized in \citep{Petryshyn1968}. The first nonasymptotic results were obtained in \citep{GodunovOgnevaProkopov1970}. The investigations in \citep{Ruhe1975} continue those initiated by Samokish. They clarify the asymptotic relationship between iterations of the form (4.1.5) and iterative methods for solving the homogeneous equation
\[
(M - \lambda_1 L) u = 0 \tag{4.1.6}
\]
% === Orig. p. 118 [s61L] ===
with preliminary improvement of the condition number through implicit reduction to the equation
\[
B^{-1} (M - \lambda_1 L) u = 0. \tag{4.1.7}
\]

Since 1977, series of investigations have been carried out on the detailed study of modified gradient methods for mesh spectral problems. Some of these results are summarized in \citep{DyakonovOrekhov1980,Dyakonov1983,Dyakonov1986en}. For example, under a natural assumption on the initial approximation, a (nonasymptotic) convergence rate bound no worse than a geometric progression has been proved for these methods, and used to justify the efficiency of choosing an operator $B$ spectrally equivalent to $L$. In \citep{Prikazchikov1980,PrikazchikovKhimich1984}, similar studies deal with the alternating triangular method.

Various aspects of the theory and practical implementation of methods with an auxiliary operator are examined, for instance, in \citep{Prokopov1968,Strakhovskaya1977,Orekhov1980,ZhukPrikazchikov1982}.

Of current interest are questions concerning the possibility of employing Chebyshev parameters and Lanczos-type algorithms in iterative methods with an auxiliary operator. The first question was answered in \citep{GavrilinIlin1981,Ilin1985}, and the second, for $B = E$, in \citep{Scott1981}. A block variant of the algorithm from \citep{Scott1981} is described in \citep{Zhihao1983}. The foundation for constructing such algorithms is the two-stage form of the methods, whose study we now undertake following the results of \citep{Knyazev1983b}.

Let us fix a parameter
\[
\lambda^n = \lambda \in (\lambda_{p+1}, \lambda_p] \tag{4.1.8}
\]
and consider the auxiliary spectral problem
\[
(M - \lambda L) v = \mu B v \tag{4.1.9}
\]
with eigenvalues $\mu_1 \geqslant \mu_2 \geqslant \ldots \geqslant \mu_{\min}$.

% === Orig. p. 119 [s61R] ===
\textbf{Theorem 4.1.1.} The following bounds hold:
\begin{gather*}
0 \leqslant \delta_0 (\lambda_i - \lambda) \leqslant \mu_i \leqslant \delta_1 (\lambda_i - \lambda), \quad i = 1, \ldots, p; \tag{4.1.10} \\
- \delta_1 (\lambda - \lambda_j) \leqslant \mu_j \leqslant - \delta_0 (\lambda - \lambda_j) < 0, \quad j > p. \tag{4.1.11}
\end{gather*}

For $p = 1$, let $v_1$ denote an eigenvector of problem (4.1.9) corresponding to $\mu_1$; for $p > 1$, let us introduce the invariant subspace $V_p$ corresponding to $\mu_1, \ldots, \mu_p$.

\textbf{Theorem 4.1.2.} For $p = 1$, the following inequality holds:
\[
0 \leqslant \lambda_1 - \lambda(v_1) \leqslant (\lambda_1 - \lambda) \left[ 1 + \frac{4\delta}{(1 - \delta)^2} \frac{\lambda_1 - \lambda_2}{\lambda_1 - \lambda} \right]^{-1}. \tag{4.1.12}
\]

One may conjecture that for $p > 1$ inequality (4.1.12) carries over to the following:
\[
0 \leqslant \lambda_p - \lambda(V_p) \leqslant (\lambda_p - \lambda) \left[ 1 + \frac{4\delta}{(1 - \delta)^2} \frac{\lambda_p - \lambda_{p+1}}{\lambda_p - \lambda} \right]^{-1}, \tag{4.1.13}
\]
where by definition $\lambda(V) \equiv \min \lambda(v)$, $v \in V \setminus 0$, is the minimum Ritz value of problem (0.1) for the subspace $V$.

Consider the following iterative process: for $n = 0, 1, \ldots$, given a vector $u^n \neq 0$, one computes $\lambda^n \equiv \lambda(u^n)$, finds the vector $v_1$ for problem (4.1.9) with $\lambda = \lambda^n$, and takes it as $u^{n+1}$. The block analogue of this process is obtained by replacing the vectors $u^n$ with $p$-dimensional subspaces $U^n$. It follows from (4.1.12) that for $\lambda_0 > \lambda_{p+1}$ the sequence $\lambda^n$ in such a process converges to $\lambda_p$ no slower than a geometric progression with ratio $(1 - \delta)^2 \times (1 + \delta)^{-2}$ and asymptotically quadratically. Asymptotically quadratic convergence of the block variant was proved in \citep{Zhihao1983}. % [math note: \lambda_0 has a subscript instead of superscript \lambda^0 in 1986 print]

As shown by the results of Chapters 1--2, the complexity of computing $v_1$ and $V_p$ is characterized mainly by the quantity $\xi_p = (\mu_p - \mu_{p+1}) / (\mu_p - \mu_{\min})$. Inequalities (4.1.10) and (4.1.11) make it possible, given known $\delta_0$ and $\delta_1$ in (4.1.1), to estimate this quantity a priori, or at least the nature of its dependence on the mesh size.

\textbf{Remark 4.1.1.} Let us point out that (4.1.12) implies the convergence of $u^n$ to $u_1$ by virtue of (2.2.9); however, (4.1.13) does not imply the convergence of $V^n$ to $U = \myspan\{u_1, \ldots, u_p\}$, and in fact such convergence fails to hold for % [math note: written V^n in 1986 print; on p. 120 the same subspace is denoted U^n]
% === Orig. p. 120 [s62L] ===
most operators $B$. By virtue of the inequality
\[
\sin^2 \ang_L (u_p; \widetilde{U}) \leqslant \{ \lambda_p - \lambda(\widetilde{U}) \} / \{ \lambda_p - \lambda_{p+1} \} \tag{4.1.14}
\]
(cf.\ the results of \S\,2.1), bound (4.1.13) ensures that the subspace $U^n$ contains a good approximation to $u_p$. Let $u_p^n$ denote the Ritz vector in the subspace $U^n$ corresponding to the Ritz value $\lambda_p^n = \lambda(U^n)$, and let $\lambda_{p-1}^n$ denote the next Ritz value after $\lambda_p^n$ in nondecreasing order. It is natural to regard $u_p^n$ as an approximation to $u_p$, but in order to obtain a bound on their closeness via (2.2.7) or (2.2.8), one must prove that $\lambda_{p-1}^n > \lambda_p$. Such a result has been proved only for the block modified ascent method of \S\,4.5.

\textbf{Remark 4.1.2.} Suppose that $\|v_1\|_B = \|u_1\|_B = 1$, and consider the vector $v_1$ as a function of $\lambda$: $v_1 = v_1(\lambda)$. Evidently, $v_1(\lambda_1) = u_1$. One can prove that the vector-valued function $v_1(\lambda)$ is differentiable at the point $\lambda = \lambda_1$. This means that as $\lambda \to \lambda_1$ the direction of the error vector $(v_1(\lambda) - u_1) / (\lambda_1 - \lambda)$ stabilizes, and hence known acceleration techniques can be applied. For $p > 1$, an analogous property likely holds.

In the iterative method described above, it is assumed that the vector $v_1$ can always be computed exactly. In practice, the vector $u^{n+1} = v_1$ is computed by means of inner iterations, which, first, converge only under the condition
\[
\tan \ang_B (v_1; u^n) < \infty \tag{4.1.15}
\]
(see \S\,1.1) and, second, yield only an approximation to $v_1$.

The validity of condition (4.1.15) in the general case $\lambda^0 < \lambda_2$ has not been proved, and it is therefore unknown whether the iterative process converges for every initial approximation $u^0$ (disregarding (4.1.15), such convergence was justified in \citep{Scott1981,GavrilinIlin1981,Ilin1985}).

The second and very substantial feature of an actual process is that the vector $u^{n+1}$ is merely an approximation to $v_1$. The case where $u^{n+1}$ is computed by means of vector iterations for solving the spectral problem (4.1.9) is studied in the next section. Note, in particular, that method (4.1.5) can likewise be viewed as a two-stage method in which a single inner iteration is performed.

% === Orig. p. 121 [s62R] ===
Before proving Theorems 4.1.1 and 4.1.2, let us introduce the Rayleigh quotient for problem (4.1.9):
\[
\mu(\cdot) \equiv [ m(\cdot) - \lambda l(\cdot) ] / b(\cdot).
\]
Here and below we use the following notation for quadratic forms:
\[
m(\cdot) \equiv (M\cdot, \cdot), \quad l(\cdot) \equiv (L\cdot, \cdot), \quad b(\cdot) \equiv (B\cdot, \cdot).
\]

The following identities are obvious:
\[
\begin{aligned}
\mu(\cdot) &= [ \lambda(\cdot) - \lambda ] l(\cdot) / b(\cdot), \\
\lambda(\cdot) &= \lambda + \mu(\cdot) b(\cdot) / l(\cdot).
\end{aligned}
\tag{4.1.16}
\]

\textit{Proof of Theorem 4.1.1.} From (4.1.16) we have
\[
\min_{j=0,1} \varphi_j(v) \leqslant \mu(v) \leqslant \max_{j=0,1} \varphi_j(v), \tag{4.1.17}
\]
where $\varphi_j(v) \equiv \delta_j \{ \lambda(v) - \lambda \}$, $v \in H \setminus 0$. Fix an arbitrary index $i \leqslant \dim H$. Since the operations $\min\{\cdot\}$ and $\max\{\cdot\}$ are monotone increasing, and hence so is the functional
\[
\max_{V_i \subseteq H,\, \dim V_i = i} \; \min_{v \in V_i \setminus 0} \{\cdot\},
\]
the latter can be applied to all parts of inequalities (4.1.17). In the new inequalities thus obtained, the middle expression equals $\mu_i$ by the Courant--Fischer minimax principle for problem (4.1.9):
\[
\mu_i = \max_{V_i \subseteq H,\, \dim V_i = i} \; \min_{v \in V_i \setminus 0} \mu(v).
\]
To transform the expressions on the left and on the right, we use the directly verifiable fact that for functions of the form $\delta_j f(x)$ in the variables $j \in \{0, 1\}$ and $x \in X$ (where $X$ is some set), when $\delta_1 \geqslant \delta_0 > 0$, the operations $\min$ or $\max$ over $j$ commute with $\min$ or $\max$ over $x$. We arrive at the inequalities
\[
\min_{j=0,1} \varphi_j^i \leqslant \mu_i \leqslant \max_{j=0,1} \varphi_j^i, \tag{4.1.18}
\]
where
% === Orig. p. 122 [s63L] ===
\[
\varphi_j^i \equiv \max_{V_i \subseteq H,\, \dim V_i = i} \; \min_{v \in V_i \setminus 0} \varphi_j(v).
\]

From the Courant--Fischer minimax principle for problem (0.1) and the definition of $\varphi_j(v)$, it follows that $\varphi_j^i = \delta_j (\lambda_i - \lambda)$. Substituting this expression for $\varphi_j^i$ into (4.1.18) and taking into account that $\lambda \in (\lambda_{p+1}, \lambda_p]$, we obtain all the inequalities in (4.1.10) and (4.1.11). Theorem 4.1.1 is proved.

Note that each of the inequalities in (4.1.10) and (4.1.11) is attainable. To verify this, it suffices to choose an operator $B$ such that the eigenvectors of problems (0.1) and (4.1.9) coincide and the only eigenvalues of the operator $B^{-1} L$ are the numbers $\delta_0$ and $\delta_1$.

\textit{Proof of Theorem 4.1.2.} Set $\hat{\lambda} \equiv \lambda(v_1)$. By virtue of (4.1.10) with $p = 1$ and (4.1.16), the inequality $\hat{\lambda} \geqslant \lambda$ holds; consequently, $\hat{\lambda} \in (\lambda_2, \lambda_1]$, and we can apply Temple's inequality (3.1.3) written in the form
\[
(\lambda_1 - \hat{\lambda})(\hat{\lambda} - \lambda_2) l(v_1) \leqslant \| (M - \hat{\lambda} L) v_1 \|_{L^{-1}}^2.
\]
We transform the right-hand side of this inequality into the expression
\[
\begin{aligned}
\| (M - \lambda L) v_1 \|_{L^{-1}}^2 &- (\hat{\lambda} - \lambda)^2 l(v_1) = \\
&= (\hat{\lambda} - \lambda)^2 l(v_1) \{ \| B v_1 \|_{L^{-1}}^2 \| v_1 \|_L^2 \| v_1 \|_B^{-4} - 1 \}
\end{aligned}
\]
In the first term of the difference inside the braces, we can make the substitution
\[
x \equiv B^{1/2} v_1, \quad S \equiv B^{1/2} L^{-1} B^{1/2}
\]
and apply the Kantorovich inequality \citep[p.~366]{GlazmanLyubich1969}
\[
\|x\|_S^2 \|x\|_{S^{-1}}^2 \|x\|^{-4} \leqslant \frac{1}{4} \{ 2 + \operatorname{cond}(S) + 1 / \operatorname{cond}(S) \},
\]
where the condition number satisfies $\operatorname{cond}(S) \leqslant \delta^{-1}$. Combining the above arguments, we obtain the inequality
\[
(\lambda_1 - \hat{\lambda})(\hat{\lambda} - \lambda_2) \leqslant \delta^{-1} / 4 \, (1 - \delta)^2 (\hat{\lambda} - \lambda)^2. \tag{4.1.19}
\]
From the inequality $\hat{\lambda} \geqslant \lambda$ and (4.1.19), it follows that $\hat{\lambda} \geqslant \lambda^*$, where $\lambda^*$ is the rightmost root of the quadratic equation
\[
\eta (\lambda^* - \lambda)^2 - (\lambda_1 - \lambda^*)(\lambda^* - \lambda^2) = 0, \quad \eta \equiv \delta^{-1} / 4 \, (1 - \delta)^2. % [math note: \lambda^2 instead of \lambda_2 in 1986 print]
\]
% === Orig. p. 123 [s63R] ===
Let us write this quadratic equation in standard form:
\[
\begin{aligned}
&a (\lambda^*)^2 + b \lambda^* + c = 0, \quad a \equiv \eta + 1, \\
&b \equiv - (2\eta\lambda + \lambda_1 + \lambda_2), \quad c \equiv \eta \lambda^2 + \lambda_1 \lambda_2.
\end{aligned}
\]
Performing straightforward manipulations, we find
\[
\lambda_1 - \lambda^* = \frac{2a\lambda_1 + b - \sqrt{b^2 - 4ac}}{2a} = \frac{a\lambda_1^2 + b\lambda_1 + c}{a\lambda_1 + b/2 + \sqrt{b^2/4 - ac}}.
\]
Observing that
\[
b^2 / 4 - ac = \left( \frac{\lambda_1 - \lambda_2}{2} \right)^2 + \eta (\lambda_1 - \lambda)(\lambda - \lambda_2) \geqslant \left( \frac{\lambda_1 - \lambda_2}{2} \right)^2,
\]
and substituting the expressions for $a$, $b$, and $c$, we obtain
\[
\lambda_1 - \lambda^* \leqslant \eta (\lambda_1 - \lambda)^2 [ \eta (\lambda_1 - \lambda) + \lambda_1 - \lambda_2 ]^{-1},
\]
whence (4.1.12) follows immediately. Theorem 4.1.2 is proved.

\textit{Proof of inequality (4.1.14).} Recall that the vector $\bar{u}_p$ from (2.2.11) satisfies the relations
\[
\bar{u}_p \in \widetilde{U}, \quad \bar{u}_p \perp_L \myspan\{ u_1, \ldots, u_{p-1} \}.
\]

From this we conclude that the following inequalities hold:
\[
\begin{aligned}
\sin^2 \ang_L (u_p, \widetilde{U}) &\leqslant \sin^2 \ang_L (u_p, \bar{u}_p) \leqslant \\
&\leqslant \{ \lambda_p - \lambda(\bar{u}_p) \} / \{ \lambda_p - \lambda_{p+1} \} \leqslant \{ \lambda_p - \lambda(\widetilde{U}) \} / \{ \lambda_p - \lambda_{p+1} \},
\end{aligned}
\]
the middle one of which follows from (1.1.3). Bound (4.1.14) is proved.

In conclusion, let us note that the term ``two-stage form of the methods'' (though perhaps not entirely felicitous) is well justified when applied to modified gradient methods. Indeed, a method of type (4.1.5) cannot be regarded as a ``full-fledged'' two-stage method, since only a single inner iteration is performed. At the same time, it is expedient to study this method as a two-stage one (see \S\,4.4).

% === Orig. p. 124 [s64L] ===
\monosection{4.2}{Abstract method with a finite number of inner iterations}

Assume $p = 1$ and consider the ``iterative'' process
\[
\lambda^n = \lambda(u^n), \quad u^{n+1} = F_n u^n, \quad n = 0, 1, \ldots, \tag{4.2.1}
\]
in which the operator $F_n$ is some polynomial $F_n(\cdot)$ in the operator $B^{-1} (M - \lambda^n L)$. Let us define the quantities
\begin{gather*}
\underline{\mu}^n = - \delta_1 (\lambda^n - \lambda_{\min}), \quad \bar{\mu}^n = - \delta_0 (\lambda^n - \lambda_2), \quad \eta^n = \delta_0 (\lambda_1 - \lambda^n), \tag{4.2.2} \\
\sigma^n = \max_{\mu \in [\underline{\mu}^n, \bar{\mu}^n]} |F_n(\mu)|^2 / |F_n(\eta^n)|^2 \leqslant \bar{\sigma}^n. \tag{4.2.3}
\end{gather*}

The convergence rate of the ``iterations'' (4.2.1) is estimated in the following theorem (an asymptotically analogous result was first obtained in \citep{Knyazev1983b}):

\textbf{Theorem 4.2.1.} If $\lambda^n > \lambda_2$, $\bar{\sigma}^n < 1$, and $\sigma^n$ is nonincreasing as a function of $\eta^n \geqslant 0$, then
\begin{gather*}
\Delta\lambda_{n+1} \leqslant \rho^n \Delta\lambda_n, \tag{4.2.4} \\
\begin{aligned}
\rho^n &= 1 - (1 - \bar{\sigma}^n) \max(\delta; g^n), \\
g^n &= \{ 1 + \delta^{-1} (1 - \delta) \bar{\sigma}^n \varepsilon^n \}^{-1} \times \\
&\times \{ 1 + \delta^{-1} / 4 \, (1 - \delta)^2 \varepsilon^n \}^{-1} \times \\
&\times \{ 1 + \delta^{-1} (1 - \delta) \sqrt{ \sigma^n \varepsilon^n / (\delta + (\bar{\sigma}^n - \delta) \varepsilon^n) } \}^{-1}, \quad \varepsilon^n \equiv \frac{\lambda_1 - \lambda^n}{\lambda_1 - \lambda_2}. % [math note: the first \sigma^n under the radical has no overbar in 1986 print, cf. (4.2.16)]
\end{aligned}
\tag{4.2.5}
\end{gather*}

\textbf{Remark 4.2.1.} The quantity $\sigma^n$ is nonincreasing in $\eta^n \geqslant 0$ if all roots of the polynomial $F_n(\mu)$ are nonpositive.

\textbf{Remark 4.2.2.} As $\lambda^n \to \lambda_1$ or as $\delta \to 1$, we have $\rho^n \to \sigma^n$, and estimate (4.2.4) is asymptotically sharp. Consequently, choosing the polynomial $F_n(\mu)$ so as to minimize $\sigma^n$ may be called asymptotic parameter optimization. % [math note: \rho^n \to \sigma^n in 1986 print assuming \bar{\sigma}^n = \sigma^n; for general \bar{\sigma}^n, (4.2.5) yields \rho^n \to \bar{\sigma}^n]

% === Orig. p. 125 [s64R] ===
\textbf{Remark 4.2.3.} If $\delta \to 0$, then it follows from (4.2.5) that $\rho^n = 1 - o(\delta)$, and estimate (4.2.4) becomes crude. Sharper estimates for small $\delta$ are proved independently for modified gradient methods in \S\,4.4. % [math note: lowercase o(\delta) in 1986 print; by meaning, (4.2.5) gives \rho^n = 1 - O(\delta)]

\textit{Proof of Theorem 4.2.1.} Let us set
\[
\begin{aligned}
u &\equiv u^n, \quad \lambda \equiv \lambda^n, \quad \hat{u} \equiv u^{n+1}, \quad \hat{\lambda} \equiv \lambda^{n+1}, \\
\underline{\mu} &\equiv \underline{\mu}^n, \quad \bar{\mu} \equiv \bar{\mu}^n, \quad \eta \equiv \eta^n
\end{aligned}
\]
and consider problem (4.1.9). We have $\mu(u) = 0$; hence, by virtue of (4.1.10), $\mu(u) \in (\mu_2, \mu_1]$. To estimate the convergence rate of method (4.2.1) in terms of (4.1.9), we use formula (1.1.9) from Theorem 1.1.1:
\[
\begin{aligned}
(\mu_1 - \hat{\mu}) / (\hat{\mu} - \mu_2) &\leqslant - \sigma \mu_1 / \mu_2, \quad \hat{\mu} \equiv \mu(\hat{u}), \\
\sigma \equiv \max_{\mu \in [\underline{\mu}, \bar{\mu}]} |P_{k^n}^n(\mu)|^2 / |P_{k^n}^n(\mu_1)|^2 &\leqslant \sigma^n \leqslant \bar{\sigma}^n < 1 % [math note: here and in (4.2.19), P_{k^n}^n denotes the polynomial F_n]
\end{aligned}
\]
(here we have used the fact that $\mu(u) = 0$). Let us rewrite this inequality in the form
\[
(1 - \hat{\mu} / \mu_1) / (1 - \hat{\mu} / \mu_2) \leqslant \sigma \tag{4.2.6}
\]
and observe that the left-hand side of (4.2.6) is increasing in $\mu_1 \geqslant \eta$ and decreasing in $\mu_2 \leqslant \bar{\mu}$ as well as in
\[
\hat{\mu} = (\hat{\lambda} - \lambda) l(\hat{u}) / b(\hat{u}) \leqslant \delta_1 (\hat{\lambda} - \lambda),
\]
moreover, $\sigma \leqslant \bar{\sigma}^n$. Using these four inequalities and making the corresponding substitutions, we transform (4.2.6) into the bound
\[
\frac{1 - \delta^{-1} (\hat{\lambda} - \lambda) / (\lambda_1 - \lambda)}{1 + \delta^{-1} (\hat{\lambda} - \lambda) / (\lambda - \lambda_2)} \leqslant \bar{\sigma}_n. % [math note: subscript n in \bar{\sigma}_n instead of superscript \bar{\sigma}^n in 1986 print]
\]
We transform it into the form
\[
\frac{\lambda_1 - \hat{\lambda}}{\hat{\lambda} - \lambda_2} \leqslant \rho_1 \frac{\lambda_1 - \lambda}{\lambda - \lambda_2}, \tag{4.2.7}
\]
where
\[
\rho_1 = 1 - \frac{\delta (1 - \bar{\sigma}^n)(\lambda_1 - \lambda_2)}{\lambda_1 - \lambda_2 - (\lambda_1 - \lambda)(1 - \bar{\sigma}^n)(1 - \delta)} \leqslant 1 - \delta (1 - \bar{\sigma}^n).
\]
% === Orig. p. 126 [s65L] ===
Here the quantity $\rho_1$ is deliberately coarsened so that the convergence rate estimate takes a simpler form. Thus, part of assertion (4.2.4), (4.2.5) is proved.

An asymptotically better estimate (as $\lambda \to \lambda_1$ or $\delta \to 1$) is obtained if, when bounding from above in (4.2.6) the quantity
\[
\frac{\hat{\mu}}{\mu_1} = \frac{\hat{\lambda} - \lambda}{\lambda(v) - \lambda} \cdot \frac{l(\hat{u}) b(v)}{b(\hat{u}) l(v)}, \quad v \equiv v_1
\]
one uses the facts that for small $\lambda_1 - \lambda$ or $1 - \delta$, in the first factor the value $\lambda(v)$ is close to $\lambda_1$:
\[
\frac{\lambda_1 - \lambda}{\lambda(v) - \lambda} \leqslant 1 + \frac{\delta^{-1}}{4} (1 - \delta)^2 \varepsilon, \quad \varepsilon \equiv \frac{\lambda_1 - \lambda}{\lambda_1 - \lambda_2} \tag{4.2.8}
\]
(which follows from (4.1.12)), while the second factor is close to unity.

Let us give a rigorous formulation of the latter statement and prove it.

\textbf{Proposition.} For $\hat{u} \neq 0$ and $v \neq 0$, the following inequality holds:
\[
\frac{l(\hat{u}) b(v)}{b(\hat{u}) l(v)} \leqslant 1 + (\delta^{-1} - 1) \sin \ang_B (\hat{u}; v). \tag{4.2.9}
\]

\textit{Proof of the Proposition.} If $(\hat{u}, v)_B = 0$, then (4.2.9) follows immediately from (4.1.1). Henceforth we assume $(\hat{u}, v)_B \neq 0$. Observing that scaling $\hat{u}$ or $v$ does not alter (4.2.9), we choose $\hat{u}$ and $v$ so that $b(\hat{u}) = 1$ and $(\hat{u} - v, v)_B = 0$, i.e., $\hat{u}$ is normalized in $H_B$ and $v$ is the $H_B$-orthogonal projection of $\hat{u}$ onto the one-dimensional subspace $\myspan\{v\}$. Then $\sin^2 \ang_B (\hat{u}; v) = b(\hat{u} - v) = 1 - b(v)$. Now, to verify (4.2.9), we can apply the technique of reduction to a ``low-dimensional'' problem. Indeed, set $H^2 \equiv \myspan\{\hat{u}, v\}$ (in the nontrivial case $\hat{u} \neq 0$) and consider the values of the quadratic forms $b(\cdot)$ and $l(\cdot)$ on the two-dimensional subspace $H^2$. The inequalities in (4.1.1) carry over to the inequalities % [math note: \hat{u} \neq 0 in 1986 print; by meaning, \hat{u} \neq v]
\[
0 < \delta_0^2 b(u) \leqslant l(u) \leqslant \delta_1^2 b(u), \quad u \in H^2 \setminus 0, \quad \delta^2 \equiv \delta_0^2 / \delta_1^2
\]
% === Orig. p. 127 [s65R] ===
(the superscript $2$ serves as a reminder that the two-dimensional case is being considered). Inequality (4.2.9) in $H^2$ takes the form
\[
\frac{l(\hat{u}) b(v)}{b(\hat{u}) l(v)} \leqslant 1 + ((\delta^2)^{-1} - 1) \sqrt{1 - b(v)}. \tag{4.2.10}
\]
From (4.2.10) and the obvious fact $\delta^2 \leqslant \delta < 1$, inequality (4.2.9) follows. It remains to prove (4.2.10) in the two-dimensional space $H^2$. % [math note: \delta^2 \leqslant \delta < 1 in 1986 print; since the superscript 2 does not denote squaring, by meaning \delta \leqslant \delta^2 \leqslant 1]

Let $z_0, z_1$ be an $H_B$-orthonormal basis of $H^2$ consisting of stationary points in $H^2 \setminus 0$ of the Rayleigh quotient $\delta(z) \equiv l(z) / b(z)$, with $\delta(z_j) = \delta_j^2$, $j = 0, 1$. Without loss of generality we may assume $\delta_0^2 = \delta^2$ and $\delta_1^2 = 1$ (henceforth the superscript $2$ in $\delta^2$ will be omitted). Let $a_0, a_1$ and $b_0, b_1$ denote the expansion coefficients of the vectors $\hat{u}$ and $v$, respectively, with respect to the basis $z_0, z_1$. In this new notation, let us write inequality (4.2.10),
\[
\frac{\delta a_0^2 + a_1^2}{a_0^2 + a_1^2} \cdot \frac{b_0^2 + b_1^2}{\delta b_0^2 + b_1^2} \leqslant 1 + (\delta^{-1} - 1) \sqrt{1 - b_0^2 - b_1^2} \tag{4.2.11}
\]
and the conditions $b(\hat{u}) = 1$, $(\hat{u} - v, v)_B = 0$:
\[
a_0^2 + a_1^2 = 1, \quad a_0 b_0 + a_1 b_1 = b_0^2 + b_1^2 \equiv b. \tag{4.2.12}
\]

Using (4.2.12), we transform (4.2.11) into the equivalent inequality
\[
(b_0^2 - a_0^2 b) / (\delta b_0^2 + b_1^2) \leqslant \delta^{-1} \sqrt{1 - b} \tag{4.2.13}
\]
and instead of (4.2.13) verify the stronger assertion
\[
b_0^2 - a_0^2 b \leqslant b \sqrt{1 - b}. \tag{4.2.14}
\]

From conditions (4.2.12) it follows that $a_0 = b_0 \pm b_1 \sqrt{b^{-1} - 1}$. Substituting this expression into the left-hand side of (4.2.14) transforms the latter into
\[
(b - 2d)(1 - b) \pm 2 \sqrt{d (b - d) b} \, \sqrt{1 - b}, \quad d \equiv b_1^2.
\]
% === Orig. p. 128 [s66L] ===
Dividing both sides of the inequality by $d b \sqrt{1 - b}$ and setting $c \equiv d^{-1}$, we arrive at the inequality
\[
\pm 2 \sqrt{c - b^{-1}} - 2 b^{-1} \sqrt{1 - b} \leqslant c (1 - \sqrt{1 - b}), \tag{4.2.15}
\]
which is equivalent to (4.2.14). Consider the nontrivial case in which the first term on the left-hand side of (4.2.15) is taken with the plus sign. For $c \geqslant b^{-1}$, the function $c (1 - \sqrt{1 - b}) - 2 \sqrt{c - b^{-1}}$ has a unique point of minimum $c^*$, determined by setting the derivative with respect to $c$ equal to zero: $1 - \sqrt{1 - b} - (c^* - b^{-1})^{-1/2} = 0$. We find $c^* = b^{-1} + (1 - \sqrt{1 - b})^{-2}$. Substituting $c = c^*$ into the nontrivial case of (4.2.15), we find after elementary manipulations that the inequality reduces to a valid identity.

The Proposition is proved.

To apply (4.2.9) to the proof of the theorem, recall that the vector $v$ is an eigenvector of problem (4.1.9) corresponding to the eigenvalue $\mu_1$. Then, as is known from (1.1.3),
\[
\sin \ang_B (\hat{u}; v) \leqslant \sqrt{ (\mu_1 - \hat{\mu}) / (\mu_1 - \mu_2) },
\]
and by means of (4.2.6) the right-hand side of this inequality can be bounded from above by $\sqrt{ - \bar{\sigma}^n \mu_1 / \mu_2 (1 - \bar{\sigma}^n \mu_1 / \mu_2)^{-1} }$. Substituting bounds for $\mu_1$ and $\mu_2$ using the inequalities $\mu_1 \leqslant \delta_1 (\lambda_1 - \lambda)$ and $\mu_2 \leqslant \bar{\mu}$, we finally obtain
\begin{gather*}
\frac{l(\hat{u}) b(v)}{b(\hat{u}) l(v)} \leqslant 1 + (\delta^{-1} - 1) \sqrt{ \bar{\sigma}^n \varepsilon (\delta + (\bar{\sigma}^n - \delta) \varepsilon^{-1} }, \tag{4.2.16} \\ % [math note: missing closing parenthesis before the exponent -1 in 1986 print; by meaning, (\delta + (\bar{\sigma}^n - \delta)\varepsilon)^{-1}]
\varepsilon \equiv (\lambda_1 - \lambda) / (\lambda_1 - \lambda_2).
\end{gather*}

Let us denote the product of the right-hand sides of inequalities (4.2.8) and (4.2.16) by $f^{-1}$ and proceed to transform the expressions in (4.2.6). From (4.2.6) follows the inequality
\[
(1 - \bar{\sigma}^n) / (1 - \bar{\sigma}^n \mu_1 / \mu_2) \leqslant \hat{\mu} / \mu_1.
\]
On the left-hand side of this inequality, using (4.1.10) and (4.1.11), we make the substitutions $\mu_1 \leqslant \delta_1 (\lambda_1 - \lambda)$ and $\mu_2 \leqslant \bar{\mu} < 0$, while we bound the right-hand side from above
% === Orig. p. 129 [s66R] ===
using the inequality
\[
\hat{\mu} / \mu_1 \leqslant f^{-1} (\hat{\lambda} - \lambda) / (\lambda_1 - \lambda).
\]
We obtain the relation
\[
\hat{\lambda} - \lambda \geqslant f \frac{(1 - \bar{\sigma}^n)(\lambda_1 - \lambda)(\lambda - \lambda_2)}{\lambda - \lambda_2 + \delta^{-1} \bar{\sigma}^n (\lambda_1 - \lambda)},
\]
which elementary manipulations bring to the form
\[
\frac{\lambda_1 - \hat{\lambda}}{\hat{\lambda} - \lambda_2} \leqslant \left\{ 1 - f \frac{1 - \bar{\sigma}^n}{1 + h} \right\} \frac{\lambda_1 - \lambda}{\lambda - \lambda_2}, \tag{4.2.17}
\]
where $h \equiv - (1 - f)(1 - \bar{\sigma}^n)\varepsilon + (\delta^{-1} - 1)\bar{\sigma}^n \varepsilon$. Let us set $g \equiv f / (1 + h)$ and drop the first (negative) term in $h$ so that the expression for $g$ takes a simpler form. Then bound (4.2.17) complements (4.2.7) and thereby leads to inequality (4.2.4) with $\rho^n$ given by (4.2.5). Theorem 4.2.1 is proved. % [orig: extra period after (4.2.7) in 1986 print]

\textit{Proof of the asymptotic attainability of estimate (4.2.4) with $\rho^n$ from (4.2.5) as $\delta \to 1$ or as $\lambda^n \to \lambda_1$.} As seen from (4.2.5), $\rho^n \approx \sigma^n$. Let us exhibit a vector $u^n$ and an operator $B$ such that
\[
\frac{\lambda_1 - \lambda^{n+1}}{\lambda^{n+1} - \lambda_2} = \sigma^n \frac{\lambda_1 - \lambda^n}{\lambda^n - \lambda_2}. \tag{4.2.18}
\]

Consider problems (0.1) and (4.1.9) in an $H_L$-orthonormal basis consisting of eigenvectors of (0.1); then the operator $L$ becomes the identity matrix and $M$ becomes diagonal. Let us choose
\[
B = \diag(\delta_0^{-1}, \delta_0^{-1}, \ldots, \delta_0^{-1}, \delta_1^{-1}).
\]
Then the eigenvectors of problems (0.1) and (4.1.9) coincide, $\lambda(v_1) = \lambda_1$, and the eigenvalues of (4.1.9) are given by relations of the form
\[
\mu_1 = \delta_0 (\lambda_1 - \lambda), \quad \mu_2 = - \delta_0 (\lambda - \lambda_2), \quad \mu_{\min} = - \delta_1 (\lambda - \lambda_{\min}),
\]
from which it follows that $\sigma = \sigma^n$. Take a vector $u^n = c_1 u_1 + c_2 u_2$. In this case, inequality (4.2.4) in the preceding proof turns into an equality (see \S\S\,1.1, 2.3). Regarding $P_{k^n}^n(\mu)$, one must assume the validity of the condition % [math note: (4.2.4) in 1986 print; by meaning, (4.2.6)]
\[
\max_{\mu \in [\underline{\mu}^n, \bar{\mu}^n]} |P_{k^n}^n(\mu)| = |P_{k^n}^n(\bar{\mu}^n)|, \tag{4.2.19}
\]
% === Orig. p. 130 [s67L] ===
which will indeed be satisfied for all nonvariational methods in \S\S\,4.3, 4.4. Observe that the vector $\hat{u}^{n+1}$ has the form $\hat{u}^{n+1} = \hat{c}_1 u_1 + \hat{c}_2 u_2$, and consequently $l(\hat{u}) / b(\hat{u}) \times b(v) / l(v) = 1$. Thus, we conclude that relation (4.2.8) holds. The asymptotic attainability of estimate (4.2.4) in Theorem 4.2.1 is proved. % [math note: \hat{u}^{n+1} and reference to (4.2.8) instead of (4.2.18) in 1986 print]

Theorem 4.2.1 admits a partial generalization to the case $\lambda^n \in (\lambda_{s+1}, \lambda_s]$. Let us redefine the quantities:
\[
\begin{aligned}
\bar{\mu}^n &= - \delta_0 (\lambda^n - \lambda_{s+1}), \quad \eta^n = \delta_0 (\lambda_s - \lambda^n), \\
\sigma^n &= \max_{\mu \in [\underline{\mu}^n, \bar{\mu}^n]} |F_n(\mu)|^2 \Big/ \min_{\mu \geqslant \eta^n} |F_n(\mu)|^2 \leqslant \bar{\sigma}^n.
\end{aligned}
\tag{4.2.20}
\]
If $\bar{\sigma}^n < 1$ and the quantity $\sigma^n$ is nonincreasing as a function of $\eta^n \geqslant 0$, then instead of (4.2.4) we have the bound
\[
\begin{aligned}
&(\lambda_s - \lambda^{n+1}) / (\lambda^{n+1} - \lambda_{s+1}) \leqslant \rho^n (\lambda_s - \lambda^n) / (\lambda^n - \lambda_{s+1}), \\
&\rho^n = 1 - \delta (1 - \bar{\sigma}^n).
\end{aligned}
\tag{4.2.21}
\]
Note that in (4.2.21) the possibility $\lambda^{n+1} > \lambda_s$ is not ruled out.

The proof of this fact is analogous to the proof of Theorem 4.2.1.

If we set $s = p$ in (4.2.20) and (4.2.21), then inequality (4.2.21) becomes a convergence rate estimate for the block analogue of method (4.2.1): % [orig: missing comma between (4.2.20) and (4.2.21) in 1986 print]
\[
\lambda^n \equiv \min_{u \in U^n} \lambda(u), \quad U^{n+1} = F_n U^n, \quad n = 0, 1, \ldots,
\]
when
\[
\lambda^n \in (\lambda_{p+1}, \lambda_p], \quad \dim U^n = p.
\]

The proof is likewise analogous to the proof of Theorem 4.2.1. The analogue of formula (4.2.6) is derived on the basis of bound (2.3.7) for problem (4.1.9).

% =============================================================================
% SOURCE MODULE: part_08_ch4_b.tex
% =============================================================================
% === Orig. p. 131 [s67R] ===
\monosection{4.3}{Methods with Chebyshev parameters and Lanczos-type methods}

When $F_n(\mu)$ is chosen in an asymptotically optimal way in the sense of Remark~4.2.2 as the Chebyshev polynomial of the first kind of degree $k^n$ mapped to the interval $[\underline{\mu}^n, \bar{\mu}^n]$, Theorem~4.2.1 holds with
\[
\sigma^n = \bigl|T_{k^n}\bigl\{(2\eta^n - \underline{\mu}^n - \bar{\mu}^n)/(\bar{\mu}^n - \underline{\mu}^n)\bigr\}\bigr|^{-2}. \tag{4.3.1}
\]
One should note the availability of the simple and asymptotically sharp estimate:
\[
\sigma^n \leqslant \bar{\sigma}^n = \bigl|T_{k^n}\bigl\{(1 + \xi)/(1 - \xi)\bigr\}\bigr|^{-2}, \tag{4.3.2}
\]
where
\[
\xi \equiv \delta(\lambda_1 - \lambda_2)/(\lambda_1 - \lambda_{\min}).
\]

The estimate (4.2.4), (4.2.5) is asymptotically sharp as $\lambda^n \to \lambda_1$ or $\delta \to 1$, since condition (4.2.19) is satisfied.

Theorem~4.2.1 with $\sigma^n$ from (4.3.1) also holds when $k^n$ steps of the Lanczos method with the operator $B^{-1}(M - \lambda^n L)$ in $H_B$ are applied to solve the generalized problem (4.1.9), or when $u^{n+1}$ is chosen from the condition
\[
\lambda(u^{n+1}) = \max_{u \in \myspan\{u^n,\, B^{-1}(M - \lambda^n L)u^n,\, \ldots,\, [B^{-1}(M - \lambda^n L)]^{k^n} u^n\}} \lambda(u). \tag{4.3.3}
\]

We leave issues of algorithm implementation aside here, referring the reader to \citep{GavrilinIlin1981,Knyazev1983b,Ilin1985,Scott1981}.

Let us pose the problem of optimizing the choice of the number of inner iterations. We call a sequence $\{k^n\}$ asymptotically optimal if $\inf \varlimsup_{n \to \infty} (\rho^n)^{1/k^n}$ is attained on it. Under condition (4.3.1), within the class of sequences $\{k^n\}$ having only one limit point (possibly infinite), sequences of the form $k^n \to \infty$, $k^n = o(n)$ are asymptotically optimal; for them,
\[
\varlimsup_{n \to \infty} (\rho)^{1/k^n} = (1 - \sqrt{\xi})^2 (1 + \sqrt{\xi})^{-2}. % [math note: in the original, (\rho)^{1/k^n} instead of (\rho^n)^{1/k^n}; this formula is left unnumbered, but below on pp. 132--133 it is referred to as (4.3.3)]
\]
This asymptotic equality shows that the estimate (4.2.4, 4.2.5) with $\sigma^n$ from (4.3.1) asymptotically coincides with the well-known convergence rate estimate for analogous optimal iterative methods for solving
% === Orig. p. 132 [s68L] ===
the homogeneous equation $(M - \lambda_1 L)u = 0$ with preconditioning by means of the operator $B$ (see, e.g., \citep{MarchukKuznetsov1975}).

\textit{Proof} of the validity of the estimate of Theorem~4.2.1 with $\sigma^n$ from (4.3.1) for Lanczos-type methods. Note that in the proof of Theorem~4.2.1, the explicit representation (4.2.1) of the vector $u^{n+1}$ is used only to derive the estimate (4.2.6). When $u^{n+1}$ is computed by $k^n$ steps of the Lanczos method with the operator $B^{-1}(M - \lambda^n L)$ in $H_B$ and initial guess $u^n$, the estimate (4.2.6) with $\sigma \leqslant \sigma^n$ from (4.3.1) also holds (see \S\,1.4). All remaining arguments in the proof of Theorem~4.2.1 stay unchanged.

In the case where $u^{n+1}$ is determined by condition (4.3.3), it suffices to use the facts that the quantity $\lambda^{n+1}$ is no smaller than in any method of the form (4.2.1) with a polynomial $F_n(\cdot)$ of degree at most $k^n$, and that for the Chebyshev method, as noted at the beginning of this section, the estimate (4.2.4, 4.2.5) holds with $\sigma^n$ from (4.3.1).

Thus, the applicability of the result of Theorem~4.2.1 with $\sigma^n$ from (4.3.1) to the two Lanczos-type methods under consideration is proved.

\textit{Proof} of the optimality of the choice $k^n \to \infty$, $k^n = o(n)$. Suppose that for some method the estimate (4.2.4, 4.2.5) holds with $\sigma^n$ from (4.3.1) (examples of such methods were given above). Let us pose the problem of asymptotic optimization of the choice of $k^n$ as $n \to \infty$: if $\inf \varlimsup (\rho^n)^{1/k^n}$ is attained on a sequence $\{k^n\}$, we call it asymptotically optimal (here and below we assume $n \to \infty$ and omit the corresponding notation under the limit sign). This definition has practical meaning when the computational cost of determining $u^{n+1}$ from $u^n$ is directly proportional to $k^n$. Let us show that under condition (4.3.1), within the class of sequences each of which converges to some value $k$ (the possibility $k = \infty$ is not excluded) and satisfies the constraint $k^n = o(n)$, the asymptotically optimal sequences are those with $k = \infty$, and for them assertion (4.3.3) holds. % [math note: the reference (4.3.3) in the original refers to the unnumbered asymptotic equality at the bottom of p. 131]
First let us prove that $\varepsilon^n / \delta^n \to 0$. % [math note: in the original, \varepsilon^n/\delta^n \to 0 (a typo for \varepsilon^n/\sigma^n \to 0, cf. p. 133)]
Since by (4.3.2) we have the inequality
\[
\sigma^n \leqslant (1 - \xi)^2 (1 + \xi)^2, % [math note: in the original, the minus sign in the exponent is omitted: it should be (1 + \xi)^{-2}, cf. (4.3.2)]
\]
the estimate (4.2.4, 4.2.5) guarantees the existence of an $\alpha \in [0, 1)$ such that $\varepsilon^n = o(\alpha^n)$. On the other hand, it follows from condition (4.3.1) that
% === Orig. p. 133 [s68R] ===
there exists a $\beta \in (0, 1)$ for which $\varlimsup \beta^{k^n} / \sigma^n < \infty$. Then
\[
\varepsilon^n / \sigma^n = \bigl\{o(\alpha^n) / \beta^{k^n}\bigr\} \times \bigl\{\beta^{k^n} / \sigma^n\bigr\} \to 0,
\]
which follows from the condition $k^n = o(n)$.

Now observe that directly from (4.2.5), using the fact that $\varepsilon^n / \sigma^n \to 0$, one can deduce that
\[
\rho^n / \sigma^n = 1 + O\bigl(\sqrt{\varepsilon^n / \sigma^n}\bigr),
\]
and then
\[
\varlimsup (\rho^n)^{1/k^n} = \varlimsup (\sigma^n)^{1/k^n}.
\]
The latter quantity can be written out explicitly using (4.3.1) and the condition $k^n \to k$. For finite $k$ it equals
\[
\bigl|T_k\bigl\{(1 + \xi)/(1 - \xi)\bigr\}\bigr|^{-2/k},
\]
and for $k = \infty$ it equals $(1 - \sqrt{\xi})^2 (1 + \sqrt{\xi})^{-2}$.

Taking into account that the function $|T_k(x)|^{1/k}$ is monotonically increasing in $k \geqslant 1$ for $x > 1$, we arrive at (4.3.3). % [math note: the reference (4.3.3) in the original refers to the unnumbered asymptotic equality at the bottom of p. 131]

In conclusion, we note that on the whole the problem of making effective use of Chebyshev parameters and Lanczos-type methods in iterative methods with an auxiliary operator remains unsolved. For example, the following important questions have been poorly studied:

(a) What happens as $\delta$ decreases? If the estimate of Theorem~4.2.1 is sharp for $\delta \approx 0$ (which is unknown), it implies that the number $k^n$ should be chosen small, i.e., one must in effect abandon two-stage methods and use modified gradient methods, in which $k^n = 1$. Can one propose a way of applying Chebyshev parameters and Lanczos-type methods that would lead to methods more effective than modified gradient methods even for small $\delta$?

(b) How should one optimally choose $k^n$---the number of inner iterations---in practice? The asymptotically optimal choice of the parameter $k^n$: $k^n \to \infty$, $k^n = o(n)$ admits various specific implementations. Several ways of determining $k^n$ are described in \citep{Scott1981}, although they lack theoretical justification.

(c) What is the efficiency of conjugate gradient methods with an auxiliary operator? We note the paper \citep{Savinov1981}, which contains some asymptotic results.

% === Orig. p. 134 [s69L] ===
\monosection{4.4}{Modified gradient methods}

By analogy with the iterative process (1.3.1), let us write the modified gradient method (4.1.5) in the form
\[
\begin{aligned}
u^{n+1} &= B^{-1}\{M u^n - \lambda^n L u^n\} + \gamma^n u^n, \\
\lambda^n &\equiv \lambda(u^n), \quad n = 0, 1, \ldots\,.
\end{aligned} \tag{4.4.1}
\]

As a corollary of Theorem~4.2.1, we find that for the method (4.4.1) with parameters
\[
\gamma^n = \delta_1(\lambda^n - \lambda_{\min}) \tag{4.4.2}
\]
the estimate (4.2.4, 4.2.5) holds, where
\[
\sigma^n = (\bar{\mu}^n - \underline{\mu}^n)^2 / (\eta^n - \underline{\mu}^n)^2 \leqslant \bar{\sigma}^n = (1 - \xi)^2. \tag{4.4.3}
\]

Recall that
\[
\xi \equiv \delta(\lambda_1 - \lambda_2)/(\lambda_1 - \lambda_{\min}).
\]

Asymptotically (as $\lambda^n \to \lambda_1$), in (4.2.5) we have
\[
\rho^n = (1 - \xi)^2 + o(1), \tag{4.4.4}
\]
and the estimate (4.2.4) is asymptotically sharp, since condition (4.2.19) is satisfied.

An asymptotically optimal choice of the parameters $\gamma_n$ in (4.4.1): % [math note: in the original, subscript \gamma_n instead of superscript \gamma^n]
\[
\gamma^n = \bigl\{\delta_1(\lambda^n - \lambda_{\min}) + \delta_0(\lambda^n - \lambda_2)\bigr\}/2 \tag{4.4.5}
\]
leads to a special case of the function $F_n(\cdot)$---the Chebyshev polynomial of degree $k^n = 1$ from \S\,4.3. In this case, the corresponding results of \S\,4.3 take the following form: for the method (4.4.1) with $\gamma^n$ from (4.4.5), the estimate (4.2.4, 4.2.5) holds with $\sigma^n$:
\[
\begin{aligned}
\sigma^n &= (\bar{\mu}^n - \underline{\mu}^n)^2 / (2\eta^n - \underline{\mu}^n - \bar{\mu}^n)^{-2} \\ % [math note: in the original, exponent -2 in the denominator (a typo for 2, cf. (4.3.1))]
&\leqslant \bar{\sigma}^n = (1 - \xi)^2 / (1 + \xi)^2,
\end{aligned} \tag{4.4.6}
\]
% === Orig. p. 135 [s69R] ===
asymptotically (as $\lambda^n \to \lambda_1$), the quantity $\rho^n$ in (4.2.5) is given by the relation
\[
\rho^n = (1 - \xi)^2 / (1 + \xi)^2 + o(1), \tag{4.4.7}
\]
and the estimate (4.2.4) is asymptotically sharp, since condition (4.2.19) is satisfied.

It also follows from the results of \S\,4.3 that the estimate (4.2.4, 4.2.5) with $\sigma^n$ from (4.4.6) holds for two variants of the method of steepest ascent, when in the iterative scheme (4.4.1) the parameter $\gamma^n$ is chosen from the condition of maximizing either $\mu(u^{n+1})$ or $\lambda(u^{n+1})$. In other words, in these methods of steepest ascent, at each iteration the Rayleigh--Ritz method is applied to problem (4.1.9) or (0.1), respectively, in the two-dimensional subspace
\[
\myspan\{u^n, w^n\}, \quad w^n \equiv B^{-1}(M - \lambda^n L)u^n.
\]
One can find $\gamma^n$ as one of the roots of the characteristic quadratic equation. For the first variant of the method, from (1.3.1) we have
\[
\gamma^n = \frac{1}{2}\left\{\sqrt{\mu^2(w^n) + 4\|w^n\|_B^2 \|u^n\|_B^{-2}} - \mu(w^n)\right\}, \tag{4.4.8}
\]
and the arguments of \S\,S.3 show the validity of the relation
\[
\gamma^n \in [-\bar{\mu}^n, -\underline{\mu}^n] = [\delta_0(\lambda^n - \lambda_2), \delta_1(\lambda^n - \lambda_{\min})].
\]
For the second variant of the method, the desired formula
\[
\begin{aligned}
\gamma^n &= \frac{1}{2}\left\{\sqrt{\mu^2(w^n) + 4\alpha^n} - \mu(w^n)\right\}, \\
\alpha^n &\equiv \bigl\{\|w^n\|_L^2 - \mu(w^n)(u^n, w^n)_L\bigr\} \|u^n\|_L^{-2}
\end{aligned} \tag{4.4.9}
\]
follows directly from the results of \citep{Samokish1958,Petryshyn1968}.

Note that the parameters $\gamma^n$ in (4.4.8) and (4.4.9) may differ substantially; in particular, in (4.4.9) the parameter $\gamma^n$ may be negative or zero, although asymptotically (as $\lambda^n \to \lambda_1$) $\gamma^n \approx -\mu(w^n) > 0$ in both formulas.

% === Orig. p. 136 [s70L] ===
In Remark~4.2.3 it was pointed out that in (4.2.5), $\rho^n = 1 - o(\delta)$ as $\delta \to 0$. For certain modified gradient methods, a sharper estimate is available for small $\delta$.

\textbf{Theorem 4.4.1.} If $\lambda^0 > \lambda_2$, then for the method $\{$(\textup{4.4.1}), (\textup{4.4.2})$\}$ the following estimate holds:
\begin{gather*}
\Delta\lambda_k \leqslant (1 - \xi)^k \Delta\lambda_0, \tag{4.4.10} \\
\xi = \delta(\lambda_1 - \lambda_2)/(\lambda_1 - \lambda_{\min}).
\end{gather*}
This estimate also holds for the method of steepest ascent $\{$(\textup{4.4.1}), (\textup{4.4.9})$\}$.

Note that inequality (4.4.10) is equivalent to the following:
\[
\Delta\lambda_{n+1} \leqslant (1 - \xi)\Delta\lambda_n, \tag{4.4.11}
\]
which can now be directly compared with the estimate (4.2.4, 4.2.5) for $\sigma^n$ from (4.4.3). Recalling relation (4.4.4), we see that asymptotically (as $\lambda^n \to \lambda_1$) the estimate (4.4.10) is crude. Thus, neither of these estimates is attainable. Apparently, a sharp convergence rate estimate for modified gradient methods is unknown. For the algorithm $\{$(\textup{4.4.1}), (\textup{4.4.2})$\}$, one may conjecture that the desired estimate has the form
\[
\Delta\lambda_{n+1} \leqslant (1 - \xi)^2 \Delta\lambda_n. \tag{4.4.12}
\]
Inequality (4.4.12) has been justified only in the practically uninteresting case where the eigenvectors of problems (0.1) and (4.1.9) coincide. The proof of this fact is analogous to the proof of Theorem~4.2.1 with natural simplifications.

A new approach to deriving sharp estimates for modified gradient methods is connected with the technique of reduction to a ``low-dimensional'' case.

\textbf{Theorem 4.4.2.} If for the method (4.4.1) in the special case $\dim H = 3$ the estimate
\[
\Delta\lambda_{n+1} \leqslant \bar{\sigma} \Delta\lambda_n, \quad \bar{\sigma} < 1, \tag{4.4.12} % [math note: in the original, the number (4.4.12) is assigned to two formulas on p. 136]
\]
holds, in which the quantity $\bar{\sigma}$ is a nonincreasing function of $\xi$, then the estimate (4.4.12) is also valid in the general case $\dim H \geqslant 3$.

% === Orig. p. 137 [s70R] ===
For $\dim H = 3$, inequality (4.4.12) takes the form of a very cumbersome algebraic inequality in eight independent variables. Verifying its validity analytically does not appear feasible, and carrying out numerical experiments on a computer is proposed for future work. In \citep{Babenko1985}, a similar approach is called ``demonstrative computation.''

\textit{Proof of Theorem 4.4.1.} Let us transform problem (0.1) to the form
\[
Lu = \lambda_\alpha^{-1} M_\alpha u, \quad M_\alpha \equiv M - \alpha L, \tag{4.4.13}
\]
where the parameter $\alpha < \lambda_{\min}$, so that $M_\alpha > 0$. The eigenvectors of problems (0.1) and (4.4.13) coincide, and the eigenvalues are related by the formula $\lambda_\alpha^{-1} = (\lambda - \alpha)^{-1}$. To compute the smallest eigenvalue of problem (4.4.13) and a corresponding eigenvector, let us apply the modified gradient method from \citep{DyakonovOrekhov1980}:
\[
\begin{aligned}
u^{n+1} &= u^n - \delta_1^{-1} B^{-1}\bigl(L - (\lambda_\alpha^n)^{-1} M_\alpha\bigr) u^n, \\
(\lambda_\alpha^n)^{-1} &= (L u^n, u^n) / (M_\alpha u^n, u^n).
\end{aligned} \tag{4.4.14}
\]
Comparing (4.4.1) and (4.4.14), we verify that the methods described by these iterative schemes are equivalent if in (4.4.1) we choose $\gamma^n = \delta_1(\lambda^n - \alpha)$. Under the condition $(\lambda_\alpha^n)^{-1} < (\lambda_2 - \alpha)^{-1}$, which is equivalent to the condition $\lambda^n > \alpha$, % [math note: in the original, \lambda^n > \alpha (for \lambda^n > \alpha, the inequality (\lambda_\alpha^n)^{-1} < (\lambda_2 - \alpha)^{-1} is equivalent to \lambda^n > \lambda_2)]
a convergence rate estimate is known for the method (4.4.14) \citep{DyakonovOrekhov1980}, which in the notation of problem (0.1) takes the following form:
\[
\lambda_1 - \lambda^{n+1} \leqslant \bigl\{1 - \delta(\lambda^n - \lambda_2)/(\lambda^n - \alpha)\bigr\}(\lambda_1 - \lambda^n). \tag{4.4.15}
\]
Inequality (4.4.15) is equivalent to the inequality
\[
\frac{\lambda_1 - \lambda^{n+1}}{\lambda^{n+1} - \lambda_2} \leqslant \left\{1 - \frac{\delta(\lambda_1 - \lambda_2)}{\lambda_1 - \alpha - (1 - \delta)(\lambda_1 - \lambda^n)}\right\} \frac{\lambda_1 - \lambda^n}{\lambda^n - \lambda_2}.
\]
The expression in braces can be increased by replacing the quantity $(1 - \delta)(\lambda_1 - \lambda^n)$ with zero in the denominator of the subtrahend. From this,
% === Orig. p. 138 [s71L] ===
the estimate
\[
\Delta\lambda_{n+1} \leqslant \bigl\{1 - \delta(\lambda_1 - \lambda_2)/(\lambda_1 - \alpha)\bigr\} \Delta\lambda_n \tag{4.4.16}
\]
follows immediately.

Recall that $\alpha < \lambda_{\min}$, and let us show that the estimate (4.4.16) also holds for $\alpha = \lambda_{\min}$. To this end, we verify the continuous dependence of both sides of inequality (4.4.16) on the parameter $\alpha$ at the point $\lambda_{\min}$ from the left and take the left-hand limit in (4.4.16) as $\alpha \to \lambda_{\min}$. On the right-hand side of inequality (4.4.16), only the expression in braces depends on $\alpha$, and from the relations $\lambda_1 > \lambda_2 \geqslant \lambda_{\min} > \alpha$ it is obvious that this dependence is continuous. On the left-hand side of (4.4.16), the quantity $\lambda^{n+1} \equiv \lambda(u^{n+1})$ depends on $\alpha$ as a composite function. The Rayleigh quotient $\lambda(\cdot)$ is a continuous function on $H \setminus 0$, and the vector $u^{n+1}$, as seen from formula (4.4.1) with $\gamma^n = \delta_1(\lambda^n - \alpha)$, depends continuously on $\alpha$. Finally, the vector $u^{n+1} \neq 0$ by virtue of the identity
\[
\begin{aligned}
\bigl\|B^{-1}\{M - \lambda(u)L\}u + \gamma u\bigr\|_B^2
&= \bigl\|B^{-1}\{M - \lambda(u)L\}u\bigr\|_B^2 + \gamma^2 \|u\|_B^2.
\end{aligned} \tag{4.4.17}
\]
Consequently, the function $\lambda^{n+1}$ depends continuously on $\alpha$. Theorem~4.4.1 is proved.

\textit{Proof of Theorem 4.4.2.} Let us denote
\[
u = u^n, \quad \lambda = \lambda^n, \quad \hat{u} = u^{n+1}, \quad \hat{\lambda} = \lambda^{n+1}, \quad \gamma = \gamma^n,
\]
and write one step of the method (4.4.1) in the form
\[
\hat{u} = \{B^{-1}M - \lambda B^{-1}L\}u + \gamma u. \tag{4.4.18}
\]
Let us define the subspace (three-dimensional in the nontrivial case $u \neq u_1$, with $\hat{u} \neq 0$ by virtue of (4.4.17))
\[
H_B^3 = \myspan\{u, \hat{u}, u_1\} \subset H_B.
\]
Let $P_B^3$ be the orthogonal projector in $H_B$ onto this subspace. Let us define the ``three-dimensional'' operators
\[
M^3 = P_B^3 B^{-1} M\big|_{H_B^3}, \quad L^3 = P_B^3 B^{-1} L\big|_{H_B^3},
\]
% === Orig. p. 139 [s71R] ===
and consider the three-dimensional eigenvalue problem
\[
M^3 u^3 = \lambda^3 L^3 u^3, \quad u^3 \in H_B^3, \tag{4.4.19}
\]
taking into account the fact that
\[
M^3 = (M^3)^* \quad \text{in } H_B^3 \quad \text{and} \quad L^3 = (L^3)^* > 0 \quad \text{in } H_B^3.
\]
The spectral problem (4.4.19) belongs to the class of symmetric generalized spectral problems of the form (0.1), and one can apply to it the modified gradient method (4.4.1) with $B^3 = E$. Taking the vector $u$ as the initial guess, let us write down the first step of this method:
\[
\hat{u}^3 = \{M^3 - \lambda^3 L\}u + \gamma u. \tag{4.4.20} % [math note: in the original, L instead of L^3]
\]
By the hypothesis of Theorem~4.4.2, the convergence rate of a single iteration of the method (4.4.20), if $\lambda^3 > \lambda_2^3$, is described by the inequality
\[
\frac{\lambda_1^3 - \hat{\lambda}^3}{\hat{\lambda}^3 - \lambda_2^3} \leqslant \bar{\sigma}^3 \frac{\lambda_1^3 - \lambda^3}{\lambda^3 - \lambda_2^3}, \quad \bar{\sigma}^3 < 1, \tag{4.4.21}
\]
where $\bar{\sigma}^3 = \bar{\sigma}(\xi^3)$, $\xi^3 = \delta^3(\lambda_1^3 - \lambda_2^3)/(\lambda_1^3 - \lambda_3^3)$, and $\lambda_1^3 > \lambda_2^3 \geqslant \lambda_3^3$ is the spectrum of problem (4.4.19). % [orig: comma after (4.4.19) in the original]

Now let us examine the relationship between problems (0.1) and (4.4.19).

The equalities
\[
\lambda(u) = \frac{(Mu, u)}{(Lu, u)} = \frac{(M^3 u, u)_B}{(L^3 u, u)_B}, \quad u \in H^3 % [math note: in the original, u \in H^3 instead of u \in H_B^3]
\]
show that the Rayleigh quotients % [orig: singular noun in the original]
for these problems coincide; hence, $\lambda^3 = \lambda$. Multiplying the left- and right-hand sides of (4.4.18) by $P_B^3$ and comparing with (4.4.20), we obtain $\hat{u}^3 = \hat{u}$, whence $\hat{\lambda}^3 = \hat{\lambda}$, and from (4.4.21) we have $\hat{\lambda} \geqslant \lambda$.

Since $u_1 \in H_B^3$, we have $\lambda_1^3 = \lambda_1$. Applying also the Courant--Fischer minimax principle to problem (0.1), we arrive at the relations
\[
\lambda_1^3 = \lambda_1 > \lambda_2 \geqslant \lambda_2^3 \geqslant \lambda_3^3 \geqslant \lambda_{\min}, \tag{4.4.22}
\]
% === Orig. p. 140 [s72L] ===
from which, taking into account the inequalities $\hat{\lambda} = \hat{\lambda}^3 \geqslant \lambda^3 \geqslant \lambda$, we obtain the bound
\[
\frac{\lambda_1 - \hat{\lambda}}{\hat{\lambda} - \lambda_2} \, \frac{\lambda - \lambda_2}{\lambda_1 - \lambda} \leqslant \frac{\lambda_1^3 - \hat{\lambda}^3}{\hat{\lambda}^3 - \lambda_2^3} \, \frac{\lambda^3 - \lambda_2^3}{\lambda_1^3 - \lambda^3}.
\]

To obtain the desired estimate
\[
\frac{\lambda_1 - \hat{\lambda}}{\hat{\lambda} - \lambda_2} \leqslant \bar{\sigma}(\xi) \frac{\lambda_1 - \lambda}{\lambda - \lambda_2} \tag{4.4.23}
\]
(cf.\ (4.4.12)), where the function $\bar{\sigma}(\cdot)$ is nonincreasing by the hypothesis of the theorem, it suffices to verify that the inequality
\[
\xi^3 = \delta^3(\lambda_1^3 - \lambda_2^3)/(\lambda_1^3 - \lambda_3^3) \geqslant \xi = \delta(\lambda_1 - \lambda_2)/(\lambda_1 - \lambda_{\min})
\]
holds.

Relations (4.4.22) show that for this purpose it suffices to check the validity of the inequality $\delta^3 \geqslant \delta$, which indeed holds because
\[
\delta^3 = \min_{u \in H_B^3} \delta(u) \Big/ \max_{u \in H_B^3} \delta(u) \geqslant \min_{u \in H} \delta(u) \Big/ \max_{u \in H} \delta(u) \geqslant \delta,
\]
where $\delta(u) \equiv (Lu, u)/(Bu, u)$. Theorem~4.4.2 is proved.

In conclusion, we note that Theorem~4.4.1 and, apparently, Theorem~4.4.2 generalize to the case $\lambda^0 \in (\lambda_{s+1}, \lambda_s]$; for instance, instead of (4.4.10) one then has the estimate
\begin{gather*}
(\lambda_s - \lambda^k)/(\lambda^k - \lambda_{s+1}) \leqslant (1 - \xi)^k (\lambda_s - \lambda^0)(\lambda^0 - \lambda_{s+1}). \tag{4.4.24} \\ % [math note: in the original, the division sign / before (\lambda^0 - \lambda_{s+1}) is omitted, cf. (4.2.21)]
\xi = \delta(\lambda_s - \lambda_{s+1})/(\lambda_s - \lambda_{\min})
\end{gather*}
(cf.\ (4.2.21)), showing that either $\lambda^k < \lambda_s$, % [math note: in the original, \lambda^k < \lambda_s (by the meaning of the estimate for \lambda^0 \in (\lambda_{s+1}, \lambda_s], the case \lambda^k > \lambda_s is possible, cf. p. 130)]
or $\lambda^k$ converges to $\lambda_s$ (see \citep{DyakonovOrekhov1980}).

% === Orig. p. 141 [s72R] ===
\monosection{4.5}{Block modified gradient methods}

To solve problem (0.1) with $p > 1$, let us consider the iterative method
\[
\begin{aligned}
U^{n+1} &= \bigl\{B^{-1}(M - \lambda^n L) + \gamma^n E\bigr\} U^n, \\
\gamma^n &\equiv \delta_1(\lambda^n - \lambda_{\min}), \quad n = 0, 1, \ldots,
\end{aligned} \tag{4.5.1}
\]
where $U^n$ are $p$-dimensional subspaces, $\lambda^n$ is the smallest Ritz value of problem (0.1) for the subspace $U^n$, and $U^0$ is the initial guess. The subspace iterations (4.5.1) are a block analogue of the modified gradient method $\{$(\textup{4.4.1}), (\textup{4.4.2})$\}$. Theorem~4.4.1 admits the following generalization.

\textbf{Theorem 4.5.1.} If $\lambda^0 > \lambda_{p+1}$, then $\dim U^k = p$ and
\[
\Delta_p \lambda_k \leqslant (1 - \xi_p)^k \Delta_p \lambda_0, \quad k \geqslant 1, \tag{4.5.2}
\]
where
\[
\begin{aligned}
\Delta_p \lambda_n &\equiv (\lambda_p - \lambda^n)/(\lambda^n - \lambda_{p+1}), \quad n = 0, k, \\
\xi_p &\equiv \delta(\lambda_p - \lambda_{p+1})/(\lambda_p - \lambda_{\min}).
\end{aligned} \tag{4.5.3}
\]

Let us denote the next Ritz value after $\lambda^n$ in nondecreasing order by $\lambda_{p-1}^n$. In Remark~4.1.1 it was pointed out that in order to determine the error of approximating the eigenvector $u_p$ by a corresponding Ritz vector $u_p^n$, one needs an estimate for $\lambda_{p-1}^n$ sharper than the inequality $\lambda_p \leqslant \lambda_{p-1}^n$. The following theorem provides a suitable result for large $n$.

\textbf{Theorem 4.5.2.} If $\lambda^n > \lambda_{p+1}$, then
\[
\lambda_{p-1}^{n+1} - \lambda_p \geqslant \delta(\lambda^n - \lambda_{p+1})(\lambda_{p-1} - \lambda^n)/(\lambda^n - \lambda_{\min}) - (\lambda_p - \lambda^n). \tag{4.5.4}
\]

The right-hand side of inequality (4.5.4) is positive for large $n$, since the subtrahend tends to zero by virtue of (4.5.2).

% === Orig. p. 142 [s73L] ===
Having an estimate of the form (4.5.4) alongside (4.5.2), on the one hand, ensures the proximity of $u_p^n$ to $u_p$ (see Theorem~2.2.2) and, on the other hand, allows one to take the subspace $U^n \ominus \myspan\{u_p^n\}$ as an initial guess for computing $\lambda_{p-1}$ and $u_{p-1}$ by a method of the form (4.5.1).

Let us propose a block analogue of the method of steepest ascent. In the subspace
\[
B^{-1}(M - \lambda^n L)U^n + U^n,
\]
whose dimension is obviously no less than $p = \dim U^n$, we carry out the Rayleigh--Ritz method for problem (0.1). We select the $p$ largest Ritz values and define $U^{n+1}$ as the linear span of the Ritz vectors corresponding to them. By definition, $\dim U^{n+1} = \dim U^n$, and the values $\lambda_p^{n+1} = \lambda^{n+1}$ and $\lambda_{p-1}^{n+1}$ in this method are no smaller than the corresponding ones in the method (4.5.1) for the same $U^n$. Consequently, the estimates of Theorems~4.5.1 and 4.5.2 also hold for the proposed block modified gradient method of steepest ascent.

\textit{Proof} of Theorems~4.5.1 and 4.5.2. The proof is analogous to that of Theorem~4.4.1. For problem (4.4.13), the iterative method proposed in \citep{DyakonovKnyazev1982} is equivalent to the method (4.5.1) with $\gamma^n = \delta_1(\lambda^n - \alpha)$, $\alpha < \lambda_{\min}$, and the estimates from that paper, written in the notation of problem (0.1), pass into the estimates (4.5.2)--(4.5.4) with $\xi_p = \delta(\lambda_p - \lambda_{p+1})/(\lambda_p - \alpha)$.

It remains to justify passing to the limit as $\alpha \to \lambda_{\min} - 0$. The nontrivial part of this proof consists in verifying the continuous dependence of the quantities $\lambda^{n+1}$ and $\lambda_{p-1}^{n+1}$ on $\alpha$. The verification relies on the following chain of arguments. The transition operator
\[
B^{-1}(M - \lambda^n L) + \gamma^n E \colon U^n \to U^{n+1}
\]
with $\gamma^n = \delta_1(\lambda^n - \alpha)$ depends continuously on $\alpha$. By virtue of representation (4.4.17), this operator is nonsingular; hence, the image $U^{n+1}$ is a $p$-dimensional subspace and also depends continuously on $\alpha$ (in the sense of the metric on the set of subspaces defined by the gap between subspaces). Consequently, the quantities $\lambda^{n+1}$ and $\lambda_{p-1}^{n+1}$ are continuous in $\alpha$, since they are Ritz values for the subspace $U^{n+1}$, and Ritz values depend continuously on the ``trial'' subspace in the Rayleigh--Ritz method. Theorems~4.5.1 and 4.5.2 are proved.

% === Orig. p. 143 [s73R] ===
\monosection{4.6}{Orthogonalization against previously computed eigenvectors}

Let $\widetilde{V}$ be the linear span of already computed approximations to eigenvectors of problem (0.1), i.e., the subspace $\widetilde{V}$ is close to some $L^{-1}M$-invariant subspace $V$. The eigenvectors not yet found lie in $V^{\perp_L}$---the orthogonal complement to $V$ in $H_L$. One can take advantage of this when computing them by iterative methods, supplementing the scheme of the methods with orthogonalization. Take, for example, the modified gradient method (4.4.1) with orthogonalization:
\[
u^{n+1} = \widetilde{P}^{\perp} B^{-1}\{M - \lambda^n L\} + \gamma^n u^n, \tag{4.6.1} % [math note: in the original, the factor u^n after \{M - \lambda^n L\} is omitted]
\]
where $\widetilde{P}^{\perp}$ is an orthogonal projector onto an orthogonal complement to $\widetilde{V}$. It is important in which inner products orthogonality is defined.

In \S\,2.7 we already discussed the choice of metric for finding the orthogonal complement to $\widetilde{V}$. For definiteness, below we consider choosing the metric of the space $H_L$.

To ensure the self-adjointness of the transition operator in (4.6.1), the projector onto the subspace $\widetilde{V}^{\perp_L}$ should be taken orthogonal in $H_B$ (rather than in $H_L$, as in \S\,2.7). Then the transition operator will be self-adjoint in $H_B$.

In this case, the backward error analysis approach is applicable, and all arguments of \S\,2.7 remain in force if problem (2.7.6) is replaced by the spectral problem
\[
\widetilde{P}_B^{\perp_L} B^{-1} M\big|_{\widetilde{V}^{\perp_L}} u = \tilde{\lambda}^{\perp} \widetilde{P}_B^{\perp_L} B^{-1} L\big|_{\widetilde{V}^{\perp_L}} u, \quad u \in \widetilde{V}^{\perp_L}.
\]

Let us discuss the practical feasibility of computing $\widetilde{P}_B^{\perp_L} u$, $u \in H$, from a given basis of the space $\widetilde{V}$. If $B = L$, the computation process, for instance by the Gram--Schmidt method, is well known. For $B \neq L$, we compute a basis of the subspace $\widetilde{W} \equiv B^{-1} L \widetilde{V}$ and recall the property of the orthogonal complement (2.6.2). We then conclude that the orthogonal projector $\widetilde{P}_B$ onto $\widetilde{V}^{\perp_L}$ coincides with the orthogonal projector $\widetilde{P}_B$ onto $\widetilde{W}^{\perp_B}$, whereas for computing the latter the standard algorithm is already applicable, since the projector $\widetilde{P}_B$ is orthogonal in $H_B$ and the orthogonal complement to the subspace $\widetilde{W}$ is also taken in $H_B$.

% =============================================================================
% UNIFIED BIBTEX BIBLIOGRAPHY
% =============================================================================
\nocite{*}
\bibliographystyle{k_monograph_en}
\bibliography{k_references}
\end{document}
